% Lesson 9 — Vocabulary: Words and expressions related to settling community disputes — Anglais 1ère A — Cours signé OnBuch+
% Programme officiel MINESEC. Compiler avec : tectonic ENA09-vocabulary-words-and-expressions-related-to-settling-co.tex
\newcommand{\DOCMATIERE}{Anglais}
\newcommand{\DOCNIVEAU}{1ère A}
\newcommand{\DOCMODULE}{Chapter 1 : Family and Social Life}
\newcommand{\DOCLECON}{ENA09}
\newcommand{\DOCTITRE}{Lesson 9 — Vocabulary: Words and expressions related to settling community disputes}
% OnBuch+ — préambule « cours signés » ENRICHI (compilé avec tectonic / XeLaTeX)
% Système visuel « Pop » d'OnBuch+ : fond crème (jamais blanc), bordures encre
% épaisses, ombres « dures » décalées, titres Archivo Black en capitales.
% Version MATHÉMATIQUES : graphes (pgfplots/tikz), boîtes pédagogiques
% (définition, propriété, méthode, attention, le savais-tu, exercice résolu,
% exercice, TP, bilan), page de garde et signature OnBuch+ ancrée en bas.
%
% Macros attendues dans main.tex AVANT \input{preamble} :
%   \DOCMATIERE  \DOCNIVEAU  \DOCMODULE  \DOCLECON (numéro)  \DOCTITRE
\documentclass[11pt]{article}

\usepackage{fontspec}
\usepackage[english,french]{babel} % français = langue principale ; l'anglais via \eng{..} / textebox
\usepackage[a4paper,margin=2cm,top=3.4cm,bottom=2.8cm,headheight=1.4cm,footskip=1.3cm]{geometry}
\usepackage{amsmath,amssymb,amsfonts}
\usepackage{mathtools} % \xmapsto, \coloneqq... souvent utilisés par les modèles
\usepackage{cancel} % simplifications barrées (bilans de forces)
% \abs{z} (module, valeur absolue) et \norm{u} (norme), utilisés par les modèles
\providecommand{\abs}{}\let\abs\relax\DeclarePairedDelimiter{\abs}{\lvert}{\rvert}
\providecommand{\norm}{}\let\norm\relax\DeclarePairedDelimiter{\norm}{\lVert}{\rVert}
\usepackage[table]{xcolor} % [table] : fournit aussi \rowcolors (alternance de lignes)
\usepackage{pagecolor}
\usepackage{tikz}
\usetikzlibrary{arrows.meta,positioning,calc,shapes.geometric,shapes.misc,shapes.arrows,decorations.pathmorphing,decorations.pathreplacing,decorations.markings,patterns,shadows,mindmap,trees,fit,backgrounds,matrix,intersections,angles,quotes}
\usepackage{pgfplots}
\usepackage[version=4]{mhchem} % équations nucléaires \ce{^{235}_{92}U}
\usepackage[european,siunitx]{circuitikz} % schémas électriques
\pgfplotsset{compat=1.18}
\usepgfplotslibrary{fillbetween,groupplots} % aires entre courbes (calcul intégral)
\usepackage{enumitem}
\usepackage{fancyhdr}
% [most] charge toutes les bibliothèques tcolorbox (listings, minted, raster,
% documentation...) et gonfle inutilement la mémoire TeX sur un document long
% (~300 boîtes) : on ne charge que ce qui est réellement utilisé.
\usepackage[skins,breakable]{tcolorbox}
\usepackage{graphicx}
\usepackage{colortbl}
\usepackage{booktabs}
\usepackage{tabularx}
\usepackage{diagbox} % case d'en-tête diagonale des tableaux à double entrée
% Colonnes extensibles centrée / gauche / droite (C, L, R), souvent utilisées par les modèles
\newcolumntype{C}{>{\centering\arraybackslash}X}
\newcolumntype{L}{>{\raggedright\arraybackslash}X}
\newcolumntype{R}{>{\raggedleft\arraybackslash}X}
\usepackage{array}
\usepackage{multicol}
\usepackage{multirow}
\usepackage{siunitx}
% parse-numbers=false : accepte aussi \SI{2,5 \times 10^{-3}}{...}
\sisetup{locale=FR,output-decimal-marker={,},group-minimum-digits=4,parse-numbers=false}
\DeclareSIUnit\ohm{\ensuremath{\symup{\Omega}}} % U+2126 absent de la police
\DeclareSIUnit\division{div} % réglages d'oscilloscope (ms/div, V/div)
\DeclareSIUnit\tr{tr} % tours (vitesse de rotation, ex. \SI{1500}{\tr\per\minute})
\DeclareSIUnit\molar{M} % concentration molaire (ex. \SI{3}{\molar}, \SI{1}{\micro\molar})
\DeclareSIUnit\an{an} % année (le modèle confond parfois avec \year, un compteur TeX) — ex. \SI{5}{\kilo\meter\per\an}
\DeclareSIUnit\FCFA{FCFA} % franc CFA (ex. \SI{500}{\FCFA\per\kilo\gram})
\DeclareSIUnit\degreeBrix{\textdegree Brix} % teneur en sucre d'un sirop/jus (réfractométrie)
\DeclareSIUnit\month{mois} % le modèle utilise parfois \month, un compteur TeX, comme unité de durée
\usepackage{qrcode}
% \degree : commande fréquemment utilisée par les modèles pour un angle ou une
% température en dehors de \SI{}{} (non fournie par les paquets chargés).
\providecommand{\degree}{^{\circ}}
% \qed (fin de démonstration), utilisé par les modèles sans amsthm
\providecommand{\qed}{\hfill$\square$}
% \barre : double barre des valeurs interdites dans un tableau de variations (tkz-tab)
\providecommand{\barre}{\ensuremath{\|}}
% environnement proof (amsthm non chargé) utilisé par les modèles
\ifdefined\proof\else\newenvironment{proof}[1][Démonstration]{\par\noindent\textit{#1.}\ }{\qed\par}\fi
\usepackage{lastpage}
\usepackage{hyperref}
\hypersetup{hidelinks}

% ============================================================
% Palette « Pop » OnBuch+ — mêmes hex que la classe P (plus_ui.dart)
% ============================================================
\definecolor{poporange}{HTML}{F59321}
\definecolor{popdark}{HTML}{E07A0C}
\definecolor{popcream}{HTML}{FFF6E8}
\definecolor{popink}{HTML}{1C1714}
\definecolor{poppurple}{HTML}{7A5AE0}
\definecolor{popgreen}{HTML}{2E9E6A}
\definecolor{poppink}{HTML}{E0457B}
\definecolor{popblue}{HTML}{2F7DE1}
\definecolor{popgold}{HTML}{C98A12}
\definecolor{popmuted}{HTML}{8A7F76}
\definecolor{popline}{HTML}{EADFD2}
\definecolor{popgray}{HTML}{8A7F76}
\colorlet{popgrey}{popgray}
% Alias tolérants (le modèle nomme parfois les couleurs différemment)
\colorlet{popwhite}{white}
\colorlet{popblack}{popink}
\colorlet{popred}{poppink}
\colorlet{popyellow}{popgold}
% Teintes claires (fonds de tableaux/encadrés) — le modèle nomme parfois
% les couleurs avec un suffixe L (ex. popblueL) sans les avoir définies.
\colorlet{poporangeL}{poporange!20}
\colorlet{popdarkL}{popdark!20}
\colorlet{poppurpleL}{poppurple!20}
\colorlet{popgreenL}{popgreen!20}
\colorlet{poppinkL}{poppink!20}
\colorlet{popblueL}{popblue!20}
\colorlet{popgoldL}{popgold!20}
\colorlet{popredL}{popred!20}
\colorlet{popgrayL}{popgray!20}
% Teintes claires pour fonds de tableaux / zones
\colorlet{popblueL}{popblue!10!white}
\colorlet{poporangeL}{poporange!14!white}
\colorlet{popgreenL}{popgreen!12!white}
\colorlet{poppurpleL}{poppurple!12!white}
\colorlet{poppinkL}{poppink!12!white}
\colorlet{popgoldL}{popgold!14!white}

\pagecolor{popcream}

% ============================================================
% Typographie
% ============================================================
\defaultfontfeatures{Ligatures=TeX}
\setmainfont{PlusJakartaSans-Medium.ttf}[Path=../../../fonts/,BoldFont=PlusJakartaSans-Bold.ttf]
% Symboles, API et emojis absents de la police principale : repli sur DejaVu Sans (caractères actifs)
\newfontface\symfont{DejaVuSans.ttf}[Path=../../../fonts/]
\makeatletter
\newcommand\symactive[1]{\catcode#1=13 \begingroup\lccode`\~=#1 \lowercase{\endgroup\def~}{{\symfont\char#1}}}
\newcommand\symblank[1]{\catcode#1=13 \begingroup\lccode`\~=#1 \lowercase{\endgroup\def~}{}}
\newcommand\symhyph[1]{\catcode#1=13 \begingroup\lccode`\~=#1 \lowercase{\endgroup\def~}{\mbox{-}}}
\newcount\symc
\newcommand\symrange[2]{\symc=#1 \@whilenum\symc<#2 \do{\expandafter\symactive\expandafter{\the\symc}\advance\symc by 1 }}
\newcommand\symblankrange[2]{\symc=#1 \@whilenum\symc<#2 \do{\expandafter\symblank\expandafter{\the\symc}\advance\symc by 1 }}
\makeatother
\symrange{"0250}{"02FF}\symactive{"2713}\symactive{"2714}\symactive{"2717}\symactive{"2718}\symactive{"2610}\symactive{"2611}\symactive{"2612}
\symhyph{"2011}\symhyph{"2E1A}\symblankrange{"1F300}{"1FAFF}\symblank{"FE0F}
% Maths en sans-serif (Fira Math) avec chiffres et lettres droites de Plus
% Jakarta, pour que les formules se fondent dans le texte.
\usepackage[math-style=ISO,bold-style=ISO]{unicode-math}
\setmathfont{FiraMath-Regular.otf}[Path=../../../fonts/]
\setmathfont{PlusJakartaSans-Medium.ttf}[Path=../../../fonts/,range={up/num,up/{Latin,latin},\mathup}]
\usepackage{icomma} % virgule décimale sans espace en mode math
% Caractères mathématiques Unicode absents des polices de texte (Plus Jakarta,
% Archivo Black) : rendus par leur équivalent mathématique, en texte comme en
% formule — sinon ils s'affichent en carré vide (ex. « Arithmétique dans ℤ »).
\usepackage{newunicodechar}
\newunicodechar{ℝ}{\ensuremath{\mathbb{R}}}
\newunicodechar{ℤ}{\ensuremath{\mathbb{Z}}}
\newunicodechar{ℂ}{\ensuremath{\mathbb{C}}}
\newunicodechar{ℕ}{\ensuremath{\mathbb{N}}}
\newunicodechar{ℚ}{\ensuremath{\mathbb{Q}}}
\newunicodechar{𝔻}{\ensuremath{\mathbb{D}}}
\newunicodechar{α}{\ensuremath{\alpha}}
\newunicodechar{β}{\ensuremath{\beta}}
\newunicodechar{γ}{\ensuremath{\gamma}}
\newunicodechar{δ}{\ensuremath{\delta}}
\newunicodechar{ε}{\ensuremath{\varepsilon}}
\newunicodechar{θ}{\ensuremath{\theta}}
\newunicodechar{λ}{\ensuremath{\lambda}}
\newunicodechar{μ}{\ensuremath{\mu}}
\newunicodechar{σ}{\ensuremath{\sigma}}
\newunicodechar{τ}{\ensuremath{\tau}}
\newunicodechar{φ}{\ensuremath{\varphi}}
\newunicodechar{ψ}{\ensuremath{\psi}}
\newunicodechar{ω}{\ensuremath{\omega}}
\newunicodechar{Σ}{\ensuremath{\Sigma}}
% majuscules produites par \MakeUppercase dans les titres (α-aminés -> Α)
\newunicodechar{Α}{\ensuremath{\alpha}}
\newunicodechar{Β}{\ensuremath{\beta}}
\newunicodechar{Γ}{\ensuremath{\Gamma}}
\newunicodechar{Δ}{\ensuremath{\Delta}}
\newunicodechar{Ε}{\ensuremath{\varepsilon}}
\newunicodechar{Θ}{\ensuremath{\Theta}}
\newunicodechar{Λ}{\ensuremath{\Lambda}}
\newunicodechar{Μ}{\ensuremath{\mu}}
\newunicodechar{Τ}{\ensuremath{\tau}}
\newunicodechar{Φ}{\ensuremath{\Phi}}
\newunicodechar{Ψ}{\ensuremath{\Psi}}
\newunicodechar{Ω}{\ensuremath{\Omega}}
\newunicodechar{↦}{\ensuremath{\mapsto}}
\newunicodechar{∈}{\ensuremath{\in}}
\newunicodechar{∉}{\ensuremath{\notin}}
\newunicodechar{∧}{\ensuremath{\wedge}}
\newunicodechar{⇔}{\ensuremath{\Leftrightarrow}}
\newunicodechar{⟺}{\ensuremath{\Longleftrightarrow}}
\newunicodechar{⇒}{\ensuremath{\Rightarrow}}
\newunicodechar{⟹}{\ensuremath{\Longrightarrow}}
\newunicodechar{∘}{\ensuremath{\circ}}
\newunicodechar{∪}{\ensuremath{\cup}}
\newunicodechar{∩}{\ensuremath{\cap}}
\newunicodechar{≡}{\ensuremath{\equiv}}
\newunicodechar{⊂}{\ensuremath{\subset}}
\newunicodechar{⊥}{\ensuremath{\perp}}
\newunicodechar{∃}{\ensuremath{\exists}}
\newunicodechar{∀}{\ensuremath{\forall}}
\newunicodechar{∛}{\ensuremath{\sqrt[3]{\,}}}
\newunicodechar{∜}{\ensuremath{\sqrt[4]{\,}}}
\newunicodechar{⃗}{\ensuremath{{}^{\to}}}
\newunicodechar{ⁿ}{\ifmmode^{n}\else\textsuperscript{\ensuremath{n}}\fi}
\newunicodechar{ˣ}{\ifmmode^{x}\else\textsuperscript{\ensuremath{x}}\fi}
\newunicodechar{ᵇ}{\ifmmode^{b}\else\textsuperscript{\ensuremath{b}}\fi}
\newunicodechar{ʳ}{\ifmmode^{r}\else\textsuperscript{\ensuremath{r}}\fi}
\newunicodechar{ᵘ}{\ifmmode^{u}\else\textsuperscript{\ensuremath{u}}\fi}
\newunicodechar{ʸ}{\ifmmode^{y}\else\textsuperscript{\ensuremath{y}}\fi}
\newunicodechar{ᵃ}{\ifmmode^{a}\else\textsuperscript{\ensuremath{a}}\fi}
\newunicodechar{ᵉ}{\ifmmode^{e}\else\textsuperscript{\ensuremath{e}}\fi}
\newunicodechar{ᵏ}{\ifmmode^{k}\else\textsuperscript{\ensuremath{k}}\fi}
\newunicodechar{ᵖ}{\ifmmode^{p}\else\textsuperscript{\ensuremath{p}}\fi}
\newunicodechar{ᴺ}{\ifmmode^{N}\else\textsuperscript{\ensuremath{N}}\fi}
\newunicodechar{ᶜ}{\ifmmode^{c}\else\textsuperscript{\ensuremath{c}}\fi}
\newunicodechar{ʰ}{\ifmmode^{h}\else\textsuperscript{\ensuremath{h}}\fi}
\newunicodechar{ᵠ}{\ifmmode^{q}\else\textsuperscript{\ensuremath{q}}\fi}
\newunicodechar{ₙ}{\ifmmode_{n}\else\textsubscript{\ensuremath{n}}\fi}
\newunicodechar{ₐ}{\ifmmode_{a}\else\textsubscript{\ensuremath{a}}\fi}
\newunicodechar{ᵢ}{\ifmmode_{i}\else\textsubscript{\ensuremath{i}}\fi}
\newunicodechar{ᵣ}{\ifmmode_{r}\else\textsubscript{\ensuremath{r}}\fi}
\newunicodechar{ₖ}{\ifmmode_{k}\else\textsubscript{\ensuremath{k}}\fi}
\newunicodechar{ᵦ}{\ifmmode_{\beta}\else\textsubscript{\ensuremath{\beta}}\fi}
\newunicodechar{ₓ}{\ifmmode_{x}\else\textsubscript{\ensuremath{x}}\fi}
\newfontfamily\popdisplay{ArchivoBlack-Regular.ttf}[Path=../../../fonts/]
\newfontfamily\poplogo{SpaceGrotesk-Bold.ttf}[Path=../../../fonts/]
\newfontfamily\popsemi{PlusJakartaSans-SemiBold.ttf}[Path=../../../fonts/]
\newfontfamily\popxbold{PlusJakartaSans-ExtraBold.ttf}[Path=../../../fonts/]
\newfontfamily\popxboldls{PlusJakartaSans-ExtraBold.ttf}[Path=../../../fonts/,LetterSpace=3.0]

\setlist[enumerate]{leftmargin=1.7em,itemsep=5pt,topsep=4pt}
\setlist[itemize]{leftmargin=1.7em,itemsep=4pt,topsep=4pt,label={\color{poporange}\textbullet}}
\setlength{\parindent}{0pt}
\setlength{\parskip}{5pt}
\renewcommand{\arraystretch}{1.25}

% pgfplots : style Pop par défaut
\pgfplotsset{
  every axis/.append style={axis line style={popink,line width=1pt},tick style={popink},
    grid style={popline,line width=0.5pt},label style={font=\small},tick label style={font=\footnotesize}},
  popcurve/.style={poporange,line width=1.8pt,no markers,smooth},
}
% Même style disponible dans un \draw TikZ simple (hors pgfplots)
\tikzset{popcurve/.style={poporange,line width=1.8pt}}
\tikzset{popgrid/.style={popline,very thin}}
\colorlet{popcurve}{poporange} % le nom du style est parfois utilisé comme couleur

% ============================================================
% Pill / squiggle
% ============================================================
\newcommand{\poppill}[3]{%
  \tikz[baseline=-0.55ex]\node[draw=popink,line width=1.2pt,rounded corners=2.6pt,%
    fill=#2,inner xsep=6.5pt,inner ysep=3pt]{%
    \popxboldls\fontsize{7}{7}\selectfont\color{#3}\MakeUppercase{#1}};%
}
\newcommand{\popsquiggle}[1]{%
  \tikz[baseline=-0.4ex]{
    \draw[#1,line width=2.6pt,line cap=round]
      plot [smooth] coordinates {(0,0.09) (0.34,0.01) (0.68,0.21) (1.02,0.01) (1.36,0.10)};
  }%
}
% Mot-clé surligné (marqueur orange sous le texte)
\newcommand{\cle}[1]{{\popxbold\color{popdark}#1}}

% ============================================================
% En-tête / pied
% ============================================================
\pagestyle{fancy}
\fancyhf{}
\renewcommand{\headrulewidth}{0pt}
\renewcommand{\footrulewidth}{0pt}
\lhead{\raisebox{-0.15ex}{\poplogo\fontsize{13}{13}\selectfont\color{popink}OnBuch\color{poporange}+}}
\rhead{\poppill{\DOCMATIERE\ \textbullet\ \DOCNIVEAU\ \textbullet\ Leçon \DOCLECON}{white}{popink}}
\lfoot{\popxbold\fontsize{7.6}{7.6}\selectfont\color{popmuted}\MakeUppercase{Cours signé OnBuch+}\ \textbullet\ {\popsemi\DOCTITRE}}
\rfoot{\tikz[baseline=-0.6ex]\node[rounded corners=8.5pt,fill=popink,minimum height=17pt,minimum width=17pt,inner xsep=5pt,inner sep=0]{\popxbold\fontsize{8}{8}\selectfont\color{popcream}\thepage\,/\,\pageref*{LastPage}};}

% ============================================================
% Titres
% ============================================================
\newcommand{\chaptitre}[2]{%
  \begin{tcolorbox}[enhanced,colback=poporange,colframe=popink,boxrule=2.6pt,arc=11pt,
      drop shadow=popink,left=15pt,right=15pt,top=12pt,bottom=11pt,before skip=3pt,after skip=18pt]
    \poppill{#2}{popink}{popcream}\par\vspace{9pt}
    {\raggedright\hyphenpenalty=10000\exhyphenpenalty=10000\popdisplay\fontsize{22}{24}\selectfont\color{popcream}\MakeUppercase{#1}\par}
  \end{tcolorbox}
}
% Section numérotée : pastille orange avec le numéro + titre Archivo
\newcounter{popsec}
\newcommand{\coursec}[1]{%
  \stepcounter{popsec}\setcounter{popsub}{0}%
  \par\vspace{16pt}\noindent
  \begin{minipage}{\linewidth}
  \tikz[baseline=-0.7ex]\node[draw=popink,line width=1.6pt,fill=poporange,rounded corners=4pt,minimum size=22pt,inner sep=0]{\popdisplay\fontsize{12}{12}\selectfont\color{popcream}\thepopsec};%
  \hspace{8pt}{\popdisplay\fontsize{15}{17}\selectfont\color{popink}\MakeUppercase{#1}}\par
  \vspace{4pt}\noindent{\color{poporange}\rule{36pt}{3.6pt}}
  \end{minipage}\par\nopagebreak\vspace{8pt}
}
\newcounter{popsub}
\newcommand{\courssub}[1]{%
  \stepcounter{popsub}%
  \par\vspace{9pt}\noindent{\popxbold\fontsize{11.5}{13}\selectfont\color{popink}{\color{poporange}\thepopsec.\thepopsub}\hspace{6pt}#1}\par\nopagebreak\vspace{4pt}
}

% ============================================================
% Boîtes pédagogiques — même recette : fond blanc, bordure épaisse colorée,
% ombre dure encre, badge plein. Titre optionnel [#1].
% ============================================================
\tcbset{popbase/.style={breakable,enhanced,colback=white,boxrule=2.3pt,arc=9pt,
  shadow={3pt}{-3pt}{0pt}{popink},left=14pt,right=14pt,top=11pt,bottom=11pt,before skip=11pt,after skip=15pt,
  fontupper=\popsemi\fontsize{10.6}{14.5}\selectfont\color{popink},
  fontlower=\popsemi\fontsize{10.6}{14.5}\selectfont\color{popink}}}
% Coche dessinée (le glyphe ✓ n'existe pas dans les polices du document)
\newcommand{\popcheck}[1]{\tikz[baseline=-0.5ex]\draw[#1,line width=1.6pt,line cap=round,line join=round](-0.13,0.0)--(-0.04,-0.09)--(0.13,0.1);}
\providecommand{\coche}{\popcheck{popgreen}}
\providecommand{\resultat}[1]{\par\smallskip\noindent\fcolorbox{popgreen}{popgreenL}{\parbox{\dimexpr\linewidth-2\fboxsep-2\fboxrule\relax}{\textbf{Résultat.} #1}}\par\smallskip}
\newcommand{\popboxhead}[3]{\poppill{#1}{#2}{white}\ifx\relax#3\relax\else\hspace{6pt}{\popxbold\fontsize{10.6}{12}\selectfont\color{popink}#3}\fi\par\vspace{7pt}}

\newtcolorbox{aretenir}[1][]{popbase,colframe=popgreen,before upper={\popboxhead{À retenir}{popgreen}{#1}}}
% Texte en anglais (document de lecture, dialogue, extrait) : cadre
% d'étude. Un titre optionnel (source, auteur) s'affiche dans l'en-tête.
\newtcolorbox{textebox}[1][]{popbase,colframe=popblue,before upper={\popboxhead{Text}{popblue}{#1}\begin{otherlanguage}{english}},after upper={\end{otherlanguage}}}
% Marques ✓ / ✗ (forme correcte / fautive) souvent utilisées par les modèles
\providecommand{\cross}{\textcolor{poppink}{\ensuremath{\times}}}
\providecommand{\xmark}{\cross}\providecommand{\crossmark}{\cross}\providecommand{\cmark}{\ensuremath{\checkmark}}
% Ligne de réponse / justification à compléter (inventée par les modèles)
\providecommand{\justification}[1]{\par\noindent\textit{Justification~:} \dotfill\par}
% \textipa (package tipa non chargé) : la police ne couvre pas l'API -> simple passage
\providecommand{\textipa}[1]{#1}
% Soulignement (ulem non chargé) : \uline, \ul
\providecommand{\uline}[1]{\underline{#1}}\providecommand{\ul}[1]{\underline{#1}}
% \glossaire{\item ...} inventée par les modèles : liste de glossaire
\providecommand{\glossaire}[1]{\begin{itemize}#1\end{itemize}}
% \alert (beamer) inventée par les modèles : texte mis en évidence
\providecommand{\alert}[1]{\textbf{\textcolor{poppink}{#1}}}
% variantes ASCII d'ordinaux écrites en macro
\providecommand{\iere}{\textsuperscript{re}}\providecommand{\ieme}{\textsuperscript{e}}\providecommand{\ere}{\textsuperscript{re}}\providecommand{\eme}{\textsuperscript{e}}\providecommand{\er}{\textsuperscript{er}}\providecommand{\ie}{\textsuperscript{e}}
% Ordinaux français écrits en macro par les modèles : 1\ière, 3\ième
\newcommand{\ière}{\textsuperscript{re}}\newcommand{\ième}{\textsuperscript{e}}
% \exemple (inventée par les modèles) : étiquette « Exemple : »
\providecommand{\exemple}{\textbf{Exemple~:}\ }
% \square utilisable aussi hors mode math (cases à cocher)
\AtBeginDocument{\let\squareMath\square\renewcommand{\square}{\ensuremath{\squareMath}}}
% Mot ou phrase en anglais à l'intérieur d'un paragraphe français
\newcommand{\eng}[1]{\ifmmode\text{\textit{#1}}\else\begingroup\selectlanguage{english}\itshape #1\endgroup\fi}
\let\esp\eng % alias toléré
% flèches d'intonation inventées par les modèles
\providecommand{\textsearrow}{}\renewcommand{\textsearrow}{\ensuremath{\searrow}}\providecommand{\textnearrow}{}\renewcommand{\textnearrow}{\ensuremath{\nearrow}}
% \fr{..} : mot français (inventé par les modèles)
\providecommand{\fr}[1]{\textit{#1}}
\newtcolorbox{exemplebox}[1][]{popbase,colframe=poporange,before upper={\popboxhead{Exemple}{poporange}{#1}}}
\newtcolorbox{definition}[1][]{popbase,colframe=popblue,before upper={\popboxhead{Définition}{popblue}{#1}}}
\newtcolorbox{propriete}[1][]{popbase,colframe=poppurple,before upper={\popboxhead{Propriété}{poppurple}{#1}}}
\newtcolorbox{methode}[1][]{popbase,colframe=popgold,before upper={\popboxhead{Méthode}{popgold}{#1}}}
\newtcolorbox{attention}[1][]{popbase,colframe=poppink,before upper={\popboxhead{Attention}{poppink}{#1}}}
\newtcolorbox{experience}[1][]{popbase,colframe=popblue,colback=popblueL,before upper={\popboxhead{Expérience}{popblue}{#1}}}
% « Le savais-tu ? » : carte violette pleine (contexte, culture, Cameroun)
\newtcolorbox{savaistu}[1][]{popbase,colback=poppurple,colframe=popink,
  fontupper=\popsemi\fontsize{10.4}{14.2}\selectfont\color{white},
  before upper={\poppill{Le savais-tu\,?}{white}{poppurple}\ifx\relax#1\relax\else\hspace{6pt}{\popxbold\color{white}#1}\fi\par\vspace{7pt}}}
% Exercice résolu : énoncé en haut, \tcblower puis la solution sur fond crème
\newcounter{popexo}\newcounter{popexr}
\newtcolorbox{exoresolu}[1][]{popbase,colframe=poporange,colbacklower=popcream,
  segmentation style={popink,line width=1.2pt,dashed},
  before upper={\stepcounter{popexr}\popboxhead{Exercice résolu \thepopexr}{poporange}{#1}},
  before lower={\poppill{Solution}{popink}{popcream}\par\vspace{6pt}}}
% Exercice (non résolu en place) — la correction est dans la partie Corrigés
\newtcolorbox{exercice}[1][]{popbase,colframe=popink,
  before upper={\stepcounter{popexo}\popboxhead{Exercice \thepopexo}{popink}{#1}}}
% Corrigé
\newtcolorbox{corrige}[1][]{popbase,colframe=popgreen,colback=popgreenL,
  before upper={\popboxhead{Corrigé}{popgreen}{#1}}}
% Objectifs / prérequis / bilan
\newtcolorbox{objectifs}{popbase,colframe=popink,colback=poporangeL,
  before upper={\popboxhead{Objectifs de la leçon}{popink}{}}}
\newtcolorbox{prerequis}{popbase,colframe=popink,colback=popblueL,
  before upper={\popboxhead{Prérequis}{popblue}{}}}
\newtcolorbox{bilan}{popbase,colframe=popink,colback=white,boxrule=2.8pt,
  before upper={\popboxhead{Fiche bilan}{poporange}{L'essentiel à connaître pour l'examen}}}
% Figure encadrée avec légende
\newtcolorbox{popfigure}[1][]{enhanced,breakable=false,colback=white,colframe=popline,boxrule=1.4pt,arc=8pt,
  left=10pt,right=10pt,top=10pt,bottom=8pt,before skip=10pt,after skip=12pt,halign=center,
  lower separated=false,halign lower=center,
  fontlower=\popsemi\fontsize{9}{11.5}\selectfont\color{popmuted},
  #1}
\newcommand{\legende}[1]{\par\vspace{6pt}{\popsemi\fontsize{9}{11.5}\selectfont\color{popmuted}#1}}

% ============================================================
% Signature OnBuch+
% ============================================================
\newcommand{\poplogoinline}[1]{{\poplogo\fontsize{#1}{#1}\selectfont\color{popink}OnBuch\color{poporange}+}}
\newcommand{\signatureonbuch}{%
  \begin{tcolorbox}[enhanced,colback=popink,colframe=popink,boxrule=2.2pt,arc=10pt,
      drop shadow=poporange,left=14pt,right=14pt,top=10pt,bottom=10pt,before skip=0pt,after skip=0pt]
    \tikz[baseline=-0.7ex]\node[circle,fill=popgreen,draw=popcream,line width=1.3pt,minimum size=22pt,inner sep=0]{\popcheck{white}};%
    \hspace{9pt}\begin{minipage}[c]{0.62\linewidth}
      {\popxbold\fontsize{12}{13}\selectfont\color{popcream}Signé\ \textbullet\ {\poplogo OnBuch\color{poporange}+}}\par\vspace{2pt}
      {\raggedright\popsemi\fontsize{8}{10.5}\selectfont\color{popline}\MakeUppercase{Équipe pédagogique OnBuch+}\\\MakeUppercase{Conforme au programme officiel MINESEC}\par}
    \end{minipage}\hfill
    {\poplogo\fontsize{22}{22}\selectfont\color{popcream}OnBuch\color{poporange}+}
  \end{tcolorbox}%
}

% ============================================================
% Page de garde
% #1 titre · #2 matière · #3 niveau · #4 module · #5 n° leçon · #6 durée ·
% #7 description / objectifs (texte ou itemize)
% ============================================================
\newcommand{\pagegarde}[7]{%
  \thispagestyle{empty}
  \newgeometry{a4paper,margin=1.7cm,top=1.6cm,bottom=1.4cm}%
  \begin{tikzpicture}[remember picture,overlay]
    \fill[poporange!18] ([xshift=-3.2cm,yshift=-3.6cm]current page.north east) circle (4.6cm);
    \fill[poppurple!14] ([xshift=2.4cm,yshift=7.6cm]current page.south west) circle (3.8cm);
  \end{tikzpicture}%
  {\poplogo\fontsize{30}{30}\selectfont\color{popink}OnBuch\color{poporange}+}\hfill
  \raisebox{0.6ex}{\poppill{Cours signé \textbullet\ Premium}{popink}{popcream}}\par
  \vspace{1pt}\popsquiggle{poppurple}\par
  \vspace{8pt}
  \begin{tcolorbox}[enhanced,colback=poporange,colframe=popink,boxrule=2.8pt,arc=13pt,
      drop shadow=popink,left=18pt,right=18pt,top=12pt,bottom=13pt,after skip=10pt]
    \poppill{#2 \textbullet\ #3}{popink}{popcream}\hspace{5pt}\poppill{#4}{white}{popink}\par\vspace{8pt}
    {\popdisplay\fontsize{14}{14}\selectfont\color{popink}LEÇON #5}\par\vspace{4pt}
    {\raggedright\hyphenpenalty=10000\exhyphenpenalty=10000\popdisplay\fontsize{25}{27}\selectfont\color{popcream}\MakeUppercase{#1}\par}
    \vspace{8pt}
    {\popxbold\fontsize{9.5}{9.5}\selectfont\color{popink}\MakeUppercase{Durée conseillée : #6}}
  \end{tcolorbox}
  \begin{tcolorbox}[enhanced,colback=white,colframe=popink,boxrule=2.2pt,arc=10pt,
      drop shadow=popink,left=16pt,right=16pt,top=10pt,bottom=10pt,after skip=8pt]
    \poppill{Dans cette leçon}{poppurple}{white}\par\vspace{5pt}
    \popsemi\fontsize{9.3}{12.6}\selectfont\color{popink}#7
  \end{tcolorbox}
  \noindent\begin{minipage}[t]{0.487\linewidth}
    \begin{tcolorbox}[enhanced,colback=popgreen,colframe=popink,boxrule=2.2pt,arc=10pt,
        drop shadow=popink,left=11pt,right=11pt,top=10pt,bottom=10pt,height=3.1cm,valign=center]
      \tcbox[colback=white,colframe=popink,boxrule=1.3pt,arc=5pt,boxsep=0pt,nobeforeafter,
        left=3pt,right=3pt,top=3pt,bottom=3pt,valign=center]{\qrcode[height=1.9cm]{https://play.google.com/store/apps/details?id=com.luvvix.onbuch&hl=fr}}%
      \hspace{8pt}\begin{minipage}[c]{0.5\linewidth}
        {\popdisplay\fontsize{11}{12.5}\selectfont\color{white}\MakeUppercase{Télécharge OnBuch}}\par\vspace{4pt}
        {\raggedright\popsemi\fontsize{8.4}{11}\selectfont\color{white}Gratuit \textbullet\ Google Play\\Tuteur IA, annales, concours, résultats\par}
      \end{minipage}
    \end{tcolorbox}
  \end{minipage}\hfill
  \begin{minipage}[t]{0.487\linewidth}
    \begin{tcolorbox}[enhanced,colback=poppurple,colframe=popink,boxrule=2.2pt,arc=10pt,
        drop shadow=popink,left=11pt,right=11pt,top=10pt,bottom=10pt,height=3.1cm,valign=center]
      \tikz[baseline=-0.7ex]\node[circle,fill=white,draw=popink,line width=1.3pt,minimum size=1.6cm,inner sep=0]{\popdisplay\fontsize{24}{24}\selectfont\color{poppurple}+};%
      \hspace{8pt}\begin{minipage}[c]{0.6\linewidth}
        {\popdisplay\fontsize{11}{12.5}\selectfont\color{white}\MakeUppercase{Passe à OnBuch+}}\par\vspace{4pt}
        {\raggedright\popsemi\fontsize{8.4}{11}\selectfont\color{white}Tous les cours signés, Tuteur IA illimité, fiches PDF — onglet OnBuch+ de l'app\par}
      \end{minipage}
    \end{tcolorbox}
  \end{minipage}
  \vspace{8pt}
  \popcontenu
  \vfill
  \signatureonbuch
  \restoregeometry
  \newpage
}

% Rangée « Ce cours contient » (page de garde)
\newcommand{\poptile}[3]{% #1 couleur #2 gros texte #3 légende
  \begin{tcolorbox}[enhanced,colback=#1,colframe=popink,boxrule=2pt,arc=9pt,shadow={2.5pt}{-2.5pt}{0pt}{popink},
      left=6pt,right=6pt,top=9pt,bottom=9pt,halign=center,valign=center,height=1.75cm,width=\linewidth,nobeforeafter]
    {\popdisplay\fontsize{12.5}{13}\selectfont\color{white}\MakeUppercase{#2}}\par\vspace{3pt}
    {\centering\popsemi\fontsize{7.8}{9.5}\selectfont\color{white}#3\par}
  \end{tcolorbox}}
\newcommand{\popcontenu}{%
  \par\noindent{\popxboldls\fontsize{7.5}{7.5}\selectfont\color{popmuted}CE COURS SIGNÉ CONTIENT}\par\vspace{7pt}
  \noindent\begin{minipage}[t]{0.235\linewidth}\poptile{popblue}{Cours}{complet, illustré, pas à pas}\end{minipage}\hfill
  \begin{minipage}[t]{0.235\linewidth}\poptile{poporange}{Méthodes}{et exercices résolus}\end{minipage}\hfill
  \begin{minipage}[t]{0.235\linewidth}\poptile{poppink}{Exercices}{type examen + corrigés}\end{minipage}\hfill
  \begin{minipage}[t]{0.235\linewidth}\poptile{popgreen}{Fiche bilan}{l'essentiel à retenir}\end{minipage}\par}

% Sommaire
\newtcolorbox{sommaire}{enhanced,colback=white,colframe=popink,boxrule=2.2pt,arc=10pt,
  drop shadow=popink,left=18pt,right=18pt,top=13pt,bottom=13pt,before skip=6pt,after skip=20pt,
  before upper={\poppill{Sommaire}{poppurple}{white}\par\vspace{9pt}\popsemi\fontsize{10.6}{14.5}\selectfont\color{popink}}}

% Signature de fin de cours — poussée tout en bas de la dernière page
\newcommand{\findecours}{%
  \par\vfill\nopagebreak
  \begin{center}\popsquiggle{poporange}\end{center}\vspace{4pt}
  \signatureonbuch
}
\AtBeginDocument{\def\llbracket{\mathopen{[\mkern-3.5mu[}}\def\rrbracket{\mathclose{]\mkern-3.5mu]}}} % intervalles d'entiers (glyphe ⟦ absent de la police)
% Glyphes absents de Fira Math : redéfinis à partir de glyphes disponibles
\AtBeginDocument{%
  \def\setminus{\mathbin{\backslash}}\let\smallsetminus\setminus
  \def\vdots{\mathord{\vbox{\baselineskip3.2pt\lineskiplimit0pt\kern4pt\hbox{.}\hbox{.}\hbox{.}}}}%
  \def\ddots{\mathinner{\mkern1mu\raise6.5pt\hbox{.}\mkern2mu\raise3.6pt\hbox{.}\mkern2mu\raise0.7pt\hbox{.}\mkern1mu}}%
  \def\triangle{\mathord{\scalebox{0.9}{$\symup{\Delta}$}}}%
  \def\nabla{\mathord{\raisebox{\depth}{\scalebox{1}[-1]{$\symup{\Delta}$}}}}%
  \def\checkmark{\popcheck{popdark}}%
  \def\star{\mathbin{\ast}}\def\bigstar{\mathord{\ast}}%
  \def\diamond{\mathbin{\scalebox{0.6}{$\Diamond$}}}%
  \def\bigcup{\mathop{\mathchoice{\raisebox{-0.3em}{\scalebox{1.7}{$\cup$}}}{\scalebox{1.25}{$\cup$}}{\cup}{\cup}}\displaylimits}%
  \def\bigcap{\mathop{\mathchoice{\raisebox{-0.3em}{\scalebox{1.7}{$\cap$}}}{\scalebox{1.25}{$\cap$}}{\cap}{\cap}}\displaylimits}%
  \def\lesseqqgtr{\mathrel{\lessgtr}}\def\bigoplus{\mathop{\oplus}\displaylimits}%
}
\providecommand{\ssi}{si et seulement si} % abréviation fréquente
\AtBeginDocument{\providecommand{\wideparen}{\overparen}} % arc de cercle

%% FIGFIX-BEGIN v4 — correctifs globaux des figures (docs/onbuch-generation/tools/figfix)
\usepackage{adjustbox}\usepackage{etoolbox}
% toute figure TikZ/pgfplots plus large que la ligne est réduite à la largeur disponible
\BeforeBeginEnvironment{tikzpicture}{\begin{adjustbox}{max width=\linewidth}}
\AfterEndEnvironment{tikzpicture}{\end{adjustbox}}
% légendes pgfplots lisibles par-dessus les courbes
\pgfplotsset{every axis/.append style={legend style={fill=white,fill opacity=0.92,text opacity=1,draw=black!15,inner sep=3pt}}}
% étiquettes d'arêtes (syntaxe edge["..."]) avec fond blanc
\tikzset{every edge quotes/.append style={fill=white,inner sep=1.2pt}}
%% FIGFIX-END
\begin{document}

\pagegarde{Lesson 9 — Vocabulary: Words and expressions related to settling community disputes}{Anglais}{1ère A}{Chapter 1 : Family and Social Life}{ENA09}{2 h}{Cette leçon de vocabulaire présente les mots et expressions liés au règlement des conflits communautaires : conflits, médiation, arbitrage, réconciliation. Elle fournit le lexique nécessaire pour décrire, discuter et résoudre un différend à l'oral comme à l'écrit.
\vspace{4pt}
\begin{multicols}{2}\small
\begin{itemize}
\item À la fin de cette leçon, je sais nommer les types de conflits et de différends en anglais
\item Je sais identifier les acteurs du règlement des conflits (mediator, arbitrator, elder, chief...)
\item Je sais employer les verbes et collocations clés : settle a dispute, resolve a conflict, reach an agreement...
\item Je sais utiliser les expressions idiomatiques liées à la paix et à la réconciliation
\item Je sais éviter les faux amis (to demand, an issue, a sentence...)
\item Je sais former des mots de la même famille (agree / agreement / disagree / disagreement...)
\end{itemize}\end{multicols}}

\begin{sommaire}
\begin{multicols}{2}
\begin{enumerate}
\item Types of conflicts and disputes
\item Actors and places of dispute settlement
\item The settlement process: key verbs and collocations
\item Idiomatic expressions and false friends
\item Word families, spelling and pronunciation
\item Using the vocabulary in context
\item Activité d'intégration
\item Méthodes et exercices résolus
\item Exercices
\item Corrigés des exercices
\item Fiche bilan
\end{enumerate}
\end{multicols}
\end{sommaire}

\begin{objectifs}
\begin{itemize}
\item À la fin de cette leçon, je sais nommer les types de conflits et de différends en anglais
\item Je sais identifier les acteurs du règlement des conflits (mediator, arbitrator, elder, chief...)
\item Je sais employer les verbes et collocations clés : settle a dispute, resolve a conflict, reach an agreement...
\item Je sais utiliser les expressions idiomatiques liées à la paix et à la réconciliation
\item Je sais éviter les faux amis (to demand, an issue, a sentence...)
\item Je sais former des mots de la même famille (agree / agreement / disagree / disagreement...)
\item Je sais prononcer correctement le vocabulaire du thème (dispute, reconcile, verdict...)
\item Je sais réemployer ce lexique dans des phrases, des dialogues et un court texte
\end{itemize}
\end{objectifs}

\begin{prerequis}
\begin{itemize}
\item Vocabulaire général de la vie sociale et familiale (family, community, neighbour)
\item Présent simple et prétérit simple pour raconter un conflit
\item Verbes modaux de base (should, must, can) pour donner un conseil
\item Notions de Seconde sur les champs lexicaux et les mots de la même famille
\end{itemize}
\end{prerequis}

\newpage

\coursec{Types of conflicts and disputes}

\courssub{Situation d'accroche : a village near Bamenda}

Imaginez la scène suivante, très courante dans les régions anglophones du Cameroun :

\begin{quote}
\emph{In a village near Bamenda, two farmers quarrel over a piece of land. Instead of going straight to court, the village chief calls both men, listens to them, and helps them reach an agreement. In Cameroon, as in Ghana or Nigeria, traditional rulers, elders and family heads often settle disputes before judges do.}
\end{quote}

« Dans un village près de Bamenda, deux fermiers se disputent un morceau de terre. Au lieu d'aller directement au tribunal, le chef du village convoque les deux hommes, les écoute et les aide à parvenir à un accord. Au Cameroun, comme au Ghana ou au Nigeria, les chefs traditionnels, les anciens et les chefs de famille règlent souvent les différends avant même que les juges n'interviennent. »

Comment parler des conflits, de la médiation et de la réconciliation en anglais ? Quelle différence y a-t-il entre \eng{a dispute}, \eng{a quarrel} et \eng{a conflict} ? C'est toute la question de cette leçon. Commençons par le point de départ de tout règlement : les \cle{types de conflits} et leurs noms.

\courssub{Observation : reading a short text}

Lisons d'abord ce court texte original. Soulignez mentalement tous les mots qui désignent un conflit ou une dispute.

\begin{textebox}[A bitter quarrel over land]
The two brothers had a bitter quarrel over their late father's land. The dispute lasted three years and divided the whole family. It all started with a simple misunderstanding: the elder brother believed the will gave him the biggest field, while the younger one claimed their father had promised it to him. Soon, the argument became violent. Neighbours stopped visiting them, and the two men would not even greet each other at the market. Some relatives took sides, and a long-standing feud began between the two branches of the family. Finally, the village chief called a meeting under the big mango tree. He listened patiently to both brothers, asked the elders for advice, and helped them reach a peaceful agreement. “A family that fights over land loses both the land and the family,” he said. The brothers shook hands, and the conflict was over at last.
\end{textebox}

\textbf{Glossaire}

\begin{center}
\begin{tabularx}{\linewidth}{|l|l|X|}
\hline
\rowcolor{popblueL} Mot anglais & Nature & Traduction française \\
\hline
bitter & adjectif & amer, acharné \\
\hline
late & adjectif & défunt, feu (\eng{her late father} : feu son père) \\
\hline
will & nom & testament \\
\hline
to claim & verbe & prétendre, affirmer \\
\hline
to take sides & expression & prendre parti \\
\hline
feud & nom & querelle durable (entre familles, clans) \\
\hline
to reach an agreement & expression & parvenir à un accord \\
\hline
to shake hands & expression & se serrer la main \\
\hline
\end{tabularx}
\end{center}

\textbf{Comprehension questions}

\begin{enumerate}
\item What are the two brothers quarrelling over?
\item How long did the dispute last?
\item Who finally helped them settle the conflict?
\item Quote two different words used in the text to name the conflict.
\end{enumerate}

\begin{corrige}[Comprehension questions]
1. \eng{They are quarrelling over their late father's land.} « Ils se disputent la terre de leur père défunt. »

2. \eng{The dispute lasted three years.} « Le différend a duré trois ans. »

3. \eng{The village chief helped them settle the conflict.} « Le chef du village les a aidés à régler le conflit. »

4. Par exemple : \eng{quarrel}, \eng{dispute}, \eng{argument}, \eng{feud}, \eng{conflict} (deux de ces mots suffisent).
\end{corrige}

\courssub{Naming conflicts : dispute, quarrel, conflict, argument}

Le texte contient cinq noms différents pour parler du même phénomène : \eng{quarrel}, \eng{dispute}, \eng{argument}, \eng{feud}, \eng{conflict}. En français, nous utilisons souvent le seul mot « dispute » ; l'anglais, lui, choisit le mot selon le \cle{registre} (familier ou formel) et la \cle{gravité} du conflit.

\begin{definition}[dispute]
\eng{A dispute} est \emph{a serious argument or disagreement, especially one that lasts a long time} : un différend, un litige, souvent sérieux et durable, fréquemment juridique ou officiel. Prononciation du nom : « DIS-pioute » (accent sur la première syllabe).
\end{definition}

\begin{definition}[quarrel]
\eng{A quarrel} est une dispute \emph{personnelle}, souvent au sein de la famille ou entre amis : une querelle, une brouille. Prononciation : « KWO-rel ».
\end{definition}

\begin{definition}[conflict]
\eng{A conflict} est le terme le plus \emph{général} : il désigne toute opposition, d'un désaccord familial à une guerre entre pays : un conflit.
\end{definition}

Voici les autres noms utiles à ce niveau :

\begin{itemize}
\item \eng{an argument} : une dispute verbale, une prise de bec (prononciation : « AR-gue-ment ») ;
\item \eng{a disagreement} : un désaccord (plus doux, plus formel qu'\eng{argument}) ;
\item \eng{a feud} : une querelle de longue durée entre deux familles, deux clans, deux villages (« fioud ») ;
\item \eng{a row} : une dispute bruyante (anglais britannique familier, prononcé « raou ») ;
\item \eng{a clash} : un affrontement, souvent bref et violent ;
\item \eng{a fight} : une bagarre, une lutte (souvent physique).
\end{itemize}

\begin{propriete}[Choisir le bon mot : dispute / quarrel / conflict]
\begin{itemize}
\item \eng{dispute} : registre \textbf{formel}, souvent juridique ou administratif : \eng{a land dispute} (un litige foncier), \eng{a trade dispute} (un conflit commercial).
\item \eng{quarrel} : registre \textbf{courant}, conflit \textbf{personnel ou familial} : \eng{a lovers' quarrel} (une dispute d'amoureux).
\item \eng{conflict} : terme \textbf{général et large}, de toutes tailles : \eng{an armed conflict} (un conflit armé), \eng{a conflict of interest} (un conflit d'intérêts).
\end{itemize}
\end{propriete}

\courssub{Naming the causes of conflicts}

D'où viennent les conflits ? Voici les collocations essentielles pour nommer les \cle{causes} :

\begin{itemize}
\item \eng{a land dispute} : un conflit foncier (à propos de la terre) ;
\item \eng{a boundary conflict} : un conflit de limites, de frontières (\eng{boundary} = limite, prononcé « BAOUN-dri ») ;
\item \eng{an inheritance quarrel} : une querelle d'héritage (\eng{inheritance} = héritage) ;
\item \eng{a misunderstanding} : un malentendu (\eng{to misunderstand} = mal comprendre) ;
\item \eng{jealousy} : la jalousie ;
\item \eng{a breach of trust} : un abus de confiance, une rupture de confiance (\eng{breach} = violation).
\end{itemize}

\begin{exemplebox}[Exemples traduits]
\eng{They fell out over money.} « Ils se sont brouillés à cause de l'argent. »

\eng{It was just a minor misunderstanding.} « Ce n'était qu'un petit malentendu. »

\eng{The inheritance quarrel divided the whole family.} « La querelle d'héritage a divisé toute la famille. »

\eng{His lies caused a breach of trust between the two friends.} « Ses mensonges ont provoqué une rupture de confiance entre les deux amis. »
\end{exemplebox}

\courssub{Adjectives to describe a conflict}

Pour qualifier un conflit, l'anglais utilise des adjectifs placés \textbf{avant} le nom (comme toujours en anglais) :

\begin{center}
\begin{tabularx}{\linewidth}{|l|X|X|}
\hline
\rowcolor{popblueL} Adjectif & Traduction & Exemple \\
\hline
serious & sérieux, grave & \eng{a serious conflict} : un conflit grave \\
\hline
minor & mineur, sans gravité & \eng{a minor disagreement} : un petit désaccord \\
\hline
long-standing & de longue date & \eng{a long-standing feud} : une vieille querelle \\
\hline
bitter & amer, acharné & \eng{a bitter dispute} : un litige acharné \\
\hline
violent & violent & \eng{a violent clash} : un affrontement violent \\
\hline
peaceful & pacifique & \eng{a peaceful settlement} : un règlement pacifique \\
\hline
\end{tabularx}
\end{center}

\begin{attention}[Ordre des mots]
En français on dit « un conflit \emph{violent} », l'adjectif après le nom. En anglais, l'adjectif se place \textbf{toujours avant} le nom : \eng{a violent conflict}, jamais \eng{a conflict violent}. De même : \eng{a serious dispute} (un litige sérieux).
\end{attention}

\courssub{Verbs of conflict}

Passons aux verbes. Observez d'abord ces phrases :

\begin{itemize}
\item \eng{The two brothers quarrel every day.} « Les deux frères se disputent tous les jours. »
\item \eng{They argued about the boundary for hours.} « Ils se sont disputés à propos de la limite pendant des heures. »
\item \eng{She fell out with her sister over the inheritance.} « Elle s'est brouillée avec sa sœur à propos de l'héritage. »
\item \eng{The two villages clashed over water.} « Les deux villages se sont affrontés à propos de l'eau. »
\item \eng{I disagree with you.} « Je ne suis pas d'accord avec vous. »
\end{itemize}

\begin{propriete}[Constructions à retenir]
\begin{itemize}
\item \eng{to quarrel with somebody about/over something} : se disputer avec quelqu'un à propos de quelque chose ;
\item \eng{to argue with somebody about something} : se disputer, débattre avec quelqu'un sur quelque chose ;
\item \eng{to fall out with somebody (over something)} : se brouiller avec quelqu'un (à cause de quelque chose) — verbe à particule ; prétérit : \eng{fell out}, participe passé : \eng{fallen out} ;
\item \eng{to fight (with somebody) over something} : se battre pour quelque chose ; prétérit : \eng{fought} ;
\item \eng{to clash (with somebody) over something} : s'affronter violemment ;
\item \eng{to disagree with somebody} : ne pas être d'accord avec quelqu'un.
\end{itemize}
Attention aux prépositions : elles font partie du verbe et doivent être apprises avec lui.
\end{propriete}

\begin{attention}[Pièges fréquents des francophones]
\begin{itemize}
\item Ne dites pas \eng{to dispute with somebody} dans la conversation courante : le verbe \eng{to dispute} signifie « contester » (\eng{to dispute a decision} : contester une décision). Pour « se disputer », utilisez \eng{to quarrel} ou \eng{to argue}.
\item \eng{argument} ne signifie pas d'abord « argument » au sens de raisonnement logique : c'est avant tout une \textbf{dispute verbale}. \eng{They had an argument} = « Ils se sont disputés » (et non « ils ont eu un argument »). Pour « avancer une raison », on dit souvent \eng{a point} : \eng{That's a good point.} « C'est un bon argument. »
\item \eng{to disagree} se construit avec \eng{with}, jamais avec \eng{to} : \eng{I disagree with him}, pas \eng{I disagree to him}.
\item Orthographe britannique : \eng{quarrel} fait \eng{quarrelled}, \eng{quarrelling} au prétérit et en \emph{-ing} (doublement du \emph{l}, comme \eng{travelled}).
\end{itemize}
\end{attention}

\courssub{Pronunciation : where to put the stress}

\begin{propriete}[L'accent tonique change le sens]
Le mot \eng{dispute} se prononce différemment selon sa nature :
\begin{itemize}
\item le \textbf{nom} : \eng{a DISpute} — « DIS-pioute » (accent sur la première syllabe) ;
\item le \textbf{verbe} : \eng{to disPUTE} — « dis-PIOUTE » (accent sur la deuxième syllabe).
\end{itemize}
C'est un phénomène fréquent en anglais (comparez \eng{a REcord} / \eng{to reCORD}).
\end{propriete}

Autres prononciations à maîtriser : \eng{argument} « AR-gue-ment » (le \emph{g} se prononce), \eng{quarrel} « KWO-rel », \eng{jealousy} « DJE-leu-si », \eng{boundary} « BAOUN-dri », \eng{feud} « fioud ».

\begin{attention}[Deux mots, deux prononciations : row]
Le nom \eng{a row} (une dispute bruyante) se prononce « raou ». Il ne faut pas le confondre avec \eng{a row} prononcé « ro », qui signifie « une rangée » (\eng{a row of houses} : une rangée de maisons). Même orthographe, deux prononciations, deux sens différents.
\end{attention}

\courssub{Mind map : the lexical field of conflicts}

\begin{popfigure}
\begin{center}
\begin{tikzpicture}[
  mindmap,
  grow cyclic,
  every node/.style={concept, align=center, text=white},
  root concept/.append style={concept color=popblue, font=\bfseries},
  level 1/.append style={level distance=3.2cm, sibling angle=120},
  level 2/.append style={level distance=2.2cm, sibling angle=45},
]
\node[root concept]{CONFLICTS}
  child[concept color=poporange]{ node {names}
      child { node {dispute} }
      child { node {quarrel} }
      child { node {feud} }
  }
  child[concept color=popgreen]{ node {causes}
      child { node {land} }
      child { node {money} }
      child { node {inheritance} }
  }
  child[concept color=poppurple]{ node {verbs}
      child { node {argue} }
      child { node {fall out} }
      child { node {clash} }
  };
\end{tikzpicture}
\end{center}
\legende{Carte mentale du champ lexical des conflits : les noms, les causes et les verbes.}
\end{popfigure}

\courssub{Tableau récapitulatif anglais / français}

\begin{center}
\begin{tabularx}{\linewidth}{|l|X|X|}
\hline
\rowcolor{popblueL} English & Français & Remarque \\
\hline
a dispute & un différend, un litige & formel, souvent juridique \\
\hline
a quarrel & une querelle, une brouille & personnel, familial \\
\hline
a conflict & un conflit & terme général \\
\hline
an argument & une dispute (verbale) & pas un « argument » logique \\
\hline
a feud & une querelle durable & entre familles ou clans \\
\hline
a row & une dispute bruyante & familier, britannique, « raou » \\
\hline
a clash & un affrontement & souvent bref et violent \\
\hline
a misunderstanding & un malentendu & cause fréquente \\
\hline
a breach of trust & un abus de confiance & expression figée \\
\hline
to fall out with & se brouiller avec & verbe à particule, prétérit \eng{fell out} \\
\hline
to disagree with & ne pas être d'accord avec & préposition \eng{with} \\
\hline
\end{tabularx}
\end{center}

\begin{savaistu}[Land disputes in anglophone Cameroon]
Dans les régions anglophones du Cameroun (le Nord-Ouest et le Sud-Ouest), les \eng{land disputes} — les conflits fonciers — sont l'une des causes les plus fréquentes de conflits villageois. La terre appartient souvent collectivement à la communauté ou à la famille, et les limites entre les parcelles ne sont pas toujours écrites. Avant d'aller devant un tribunal (\eng{a court of law}), les plaignants s'adressent aux \eng{fons} (au Nord-Ouest) et aux \eng{chiefs} (au Sud-Ouest), les \eng{traditional rulers} (chefs traditionnels), qui rendent la justice coutumière sous l'arbre à palabres. Ce système de médiation existe aussi au Ghana et au Nigeria.
\end{savaistu}

\begin{exoresolu}[Choosing the right word]
Énoncé : complétez chaque phrase avec le mot qui convient parmi : \eng{dispute}, \eng{quarrel}, \eng{conflict}, \eng{misunderstanding}, \eng{feud}.

1. The two families have had a \ldots\ for twenty years over a piece of land.
2. It was only a \ldots : she thought the meeting was at ten, not at nine.
3. The lawyer is dealing with a land \ldots\ between two villages.
4. The children had a \ldots\ over a toy, but they are friends again now.
5. The armed \ldots\ in the region has lasted too long.

\tcblower
Solution détaillée :

1. \eng{feud} : une querelle qui dure vingt ans entre deux familles est une \eng{feud} (querelle durable entre clans).
2. \eng{misunderstanding} : une confusion d'horaires est un malentendu.
3. \eng{dispute} : contexte juridique (\eng{lawyer}) ; la collocation est \eng{a land dispute}.
4. \eng{quarrel} : dispute personnelle et légère entre enfants.
5. \eng{conflict} : la collocation figée est \eng{an armed conflict} (un conflit armé).
\end{exoresolu}

\begin{exercice}[Your turn]
Traduisez en anglais, en choisissant le mot juste :

1. Les deux voisins se sont brouillés à cause d'une limite de terrain.
2. Ce n'était qu'un malentendu mineur.
3. Le chef a réglé un vieux conflit entre les deux familles.
4. Je ne suis pas d'accord avec toi.
\end{exercice}

\begin{corrige}[Exercice « Your turn »]
1. \eng{The two neighbours fell out over a land boundary.} (ou \eng{quarrelled over a boundary})
2. \eng{It was just a minor misunderstanding.}
3. \eng{The chief settled a long-standing conflict between the two families.}
4. \eng{I disagree with you.}
\end{corrige}

\begin{aretenir}[L'essentiel de la section]
\begin{itemize}
\item \eng{dispute} = différend sérieux, souvent juridique ; \eng{quarrel} = querelle personnelle ; \eng{conflict} = terme général.
\item Causes : \eng{a land dispute}, \eng{a boundary conflict}, \eng{an inheritance quarrel}, \eng{a misunderstanding}, \eng{jealousy}, \eng{a breach of trust}.
\item Verbes : \eng{to quarrel with}, \eng{to argue about}, \eng{to fall out with} (prétérit \eng{fell out}), \eng{to clash over}, \eng{to disagree with}.
\item Prononciation : \eng{a DISpute} (nom) mais \eng{to disPUTE} (verbe) ; \eng{ARgument}, \eng{KWOrel} ; \eng{a row} (dispute) se prononce « raou ».
\item Piège majeur : \eng{an argument} est une dispute, pas un raisonnement logique.
\end{itemize}
\end{aretenir}

\coursec{Actors and places of dispute settlement}

Dans la section précédente, nous avons appris à nommer les différents \cle{types of conflicts and disputes} : \eng{a quarrel}, \eng{a disagreement}, \eng{a land dispute}, \eng{a family conflict}. Mais une fois le conflit déclaré, une question se pose immédiatement : \emph{who settles it, and where?} Qui règle le différend, et où ? Au Cameroun comme dans les pays anglophones, il existe deux grandes voies : la voie traditionnelle (les anciens, le chef, le conseil de famille) et la voie judiciaire (l'avocat, le tribunal, le juge). Cette section vous donne tout le vocabulaire des \cle{actors} (acteurs) et des \cle{places} (lieux) du règlement des conflits.

\courssub{Discovering the vocabulary in context}

Lisons d'abord ce court texte. Soulignez mentalement tous les mots qui désignent des personnes ou des lieux.

\begin{textebox}[A dispute in Mbankomo]
When neighbours cannot agree, the village chief acts as a mediator. If mediation fails, the case may go to the customary court.

Last month, in Mbankomo, a farmer called Mr Essomba accused his neighbour, Mrs Abena, of planting cassava on his land. The two families quarrelled for weeks. Finally, the family head called a family meeting, but nobody could calm the two opponents. The parties concerned then went to the palace, where the traditional ruler received them with two elders. The elders heard both parties patiently. One witness testified that the boundary had been moved during the rainy season. The chief, acting as a mediator, proposed a fair solution: each farmer would keep the land he had cultivated for ten years. Both sides accepted.

But not every dispute ends so peacefully. In a nearby village, a trader decided to sue his partner who had refused to pay a debt. The plaintiff hired a lawyer, and the defendant hired one too. The case was heard at the tribunal in Yaoundé. Several witnesses came to testify. The judge listened to everybody and gave his judgment after two weeks. The defendant had to pay the money, plus the court costs. “Going to court is slow and expensive,” the trader said later. “Next time, I will first ask a counsellor or an arbitrator to help us. An arbitrator's decision is binding, and it is much faster than a trial.”
\end{textebox}

\textbf{Glossaire}\par

\begin{center}
\begin{tabularx}{\linewidth}{|l|l|X|}
\hline
\rowcolor{popblueL} \textbf{Word} & \textbf{Nature} & \textbf{French} \\
\hline
a mediator & noun & un médiateur \\
\hline
a customary court & noun & un tribunal coutumier \\
\hline
a family head & noun & un chef de famille \\
\hline
a witness & noun & un témoin \\
\hline
to testify & verb & témoigner, déposer \\
\hline
a boundary & noun & une limite, une frontière \\
\hline
to sue & verb & poursuivre en justice \\
\hline
the plaintiff & noun & le plaignant \\
\hline
the defendant & noun & le défendeur, l'accusé \\
\hline
binding & adjective & contraignant, qui oblige \\
\hline
a trial & noun & un procès \\
\hline
\end{tabularx}
\end{center}

\textbf{Comprehension questions}\par

\begin{enumerate}
\item Who acted as a mediator in the land dispute between Mr Essomba and Mrs Abena?
\item What did the witness testify?
\item Why did the trader prefer an arbitrator to a trial?
\item Where was the second case heard?
\end{enumerate}

\courssub{The people: actors of dispute settlement}

\begin{definition}[Mediator]
A \cle{mediator} is a person who helps two sides reach an agreement without taking sides. Un médiateur : une personne neutre et impartiale qui aide deux parties à trouver un accord. Prononciation : « MI-di-é-teur ». Le verbe correspondant est \eng{to mediate} (« MI-di-éte »), et l'action est \eng{mediation} (« mi-di-É-chonne »).
\end{definition}

Voici les principaux acteurs, classés par fonction :

\begin{center}
\begin{tabularx}{\linewidth}{|l|X|X|}
\hline
\rowcolor{popblueL} \textbf{English} & \textbf{Français} & \textbf{Role / Rôle} \\
\hline
a mediator & un médiateur & helps the parties talk and agree (propose) \\
\hline
an arbitrator & un arbitre & listens and takes a binding decision (décide) \\
\hline
a negotiator & un négociateur & discusses to reach an agreement \\
\hline
a judge & un juge & gives judgment in a court of law \\
\hline
a lawyer & un avocat & defends a client in court \\
\hline
a witness & un témoin & says what he or she saw or heard \\
\hline
an elder & un ancien, un notable & gives advice in the community \\
\hline
a traditional ruler / a chief & un chef traditionnel & settles disputes in the village \\
\hline
a family head & un chef de famille & settles family conflicts \\
\hline
a counsellor & un conseiller & guides and advises people \\
\hline
\end{tabularx}
\end{center}

\begin{propriete}[Mediator or arbitrator?]
Attention à la différence essentielle :
\begin{itemize}
\item The \cle{mediator} \textbf{proposes} : il aide les parties à dialoguer, mais il ne décide rien. Les parties restent libres d'accepter ou de refuser.
\item The \cle{arbitrator} \textbf{decides} : après avoir entendu les deux parties, il rend une décision qui s'impose à tous. \eng{His decision is binding.} = Sa décision est contraignante.
\end{itemize}
\eng{The mediator suggested a compromise.} = Le médiateur a suggéré un compromis.\\
\eng{The arbitrator ordered the company to pay the workers.} = L'arbitre a ordonné à l'entreprise de payer les ouvriers.
\end{propriete}

\begin{definition}[The parties in a dispute]
Les \cle{parties} sont les personnes engagées dans le conflit :
\begin{itemize}
\item \cle{the plaintiff} ou \cle{the complainant} : le plaignant, celui qui porte plainte. Prononciation : « PLEIN-tif ».
\item \cle{the defendant} : le défendeur, l'accusé, celui qui se défend.
\item \cle{the parties concerned} : les parties concernées, toutes les personnes impliquées.
\item \cle{the opponent} : l'adversaire, la partie adverse.
\end{itemize}
\eng{The plaintiff accused the defendant of destroying his fence.} = Le plaignant a accusé le défendeur d'avoir détruit sa clôture.
\end{definition}

\courssub{The places and institutions}

\begin{center}
\begin{tabularx}{\linewidth}{|l|X|X|}
\hline
\rowcolor{popblueL} \textbf{English} & \textbf{Français} & \textbf{Remarque} \\
\hline
a court (of law) & un tribunal (de droit) & general word ; \eng{to go to court} = aller au tribunal \\
\hline
a tribunal & un tribunal & often for special matters \\
\hline
a customary court & un tribunal coutumier & applies local customs \\
\hline
a village council & un conseil de village & elders and notables \\
\hline
a family meeting & une réunion de famille & first step for family disputes \\
\hline
a palace & une chefferie & the chief's residence and court \\
\hline
\end{tabularx}
\end{center}

\begin{attention}[Palace : un faux sens possible]
\eng{A palace} signifie d'abord « un palais » (Buckingham Palace), mais dans le contexte camerounais et africain anglophone, \eng{the palace} désigne couramment la \textbf{chefferie}, c'est-à-dire la résidence du chef traditionnel où se tiennent les audiences. \eng{The case was settled at the palace.} = L'affaire a été réglée à la chefferie. Ne traduisez donc pas toujours \eng{palace} par « palais » !
\end{attention}

\courssub{The verbs: what the actors do}

Chaque acteur a son verbe. Apprenez ces collocations verbe + acteur :

\begin{center}
\begin{tabularx}{\linewidth}{|l|X|X|}
\hline
\rowcolor{popblueL} \textbf{Verb} & \textbf{Français} & \textbf{Example} \\
\hline
to mediate & servir de médiateur & The chief mediated between the two families. \\
\hline
to arbitrate & arbitrer & An arbitrator arbitrated the trade dispute. \\
\hline
to judge / to try & juger / instruire & The judge tried the case fairly. \\
\hline
to hear a case & entendre une affaire & The court heard the case on Monday. \\
\hline
to testify & témoigner & Three witnesses testified in court. \\
\hline
to sue & poursuivre en justice & She decided to sue her neighbour. \\
\hline
to accuse (of) & accuser (de) & He accused his brother of lying. \\
\hline
to defend & défendre & The lawyer defended the accused man. \\
\hline
\end{tabularx}
\end{center}

\begin{exemplebox}[Exemples traduits]
\eng{The elders heard both parties.} = Les anciens ont entendu les deux parties.\\
\eng{She decided to sue her neighbour.} = Elle a décidé de poursuivre son voisin en justice.\\
\eng{The witness testified that he had seen everything.} = Le témoin a déclaré sous serment qu'il avait tout vu.\\
\eng{The lawyer defended her client brilliantly.} = L'avocate a défendu son client avec brio.\\
\eng{The customary court heard the land case last Friday.} = Le tribunal coutumier a entendu l'affaire de terrain vendredi dernier.
\end{exemplebox}

\begin{attention}[To accuse OF, to sue FOR]
Les prépositions changent tout :
\begin{itemize}
\item \eng{to accuse somebody \textbf{of} something} : \eng{He accused her of stealing his money.} (Il l'a accusée d'avoir volé son argent.) Jamais \emph{to accuse somebody for}.
\item \eng{to sue somebody \textbf{for} something} : \eng{She sued the company for unfair dismissal.} (Elle a poursuivi l'entreprise pour licenciement abusif.)
\end{itemize}
\end{attention}

\courssub{Word families (Dérivation)}

\begin{propriete}[Noms, verbes et adjectifs apparentés]
Maîtriser les familles de mots permet d'enrichir son expression écrite et orale.
\begin{center}
\begin{tabular}{|l|l|l|l|}
\hline
\rowcolor{popblueL} \textbf{Verb} & \textbf{Noun (person)} & \textbf{Noun (action/state)} & \textbf{Adjective} \\
\hline
to mediate & a mediator & mediation & mediatory \\
\hline
to arbitrate & an arbitrator & arbitration & arbitral \\
\hline
to negotiate & a negotiator & negotiation & negotiable \\
\hline
to judge & a judge & judgment / judgement & judicial \\
\hline
to sue & a plaintiff / a suer & a lawsuit / a suit & — \\
\hline
to defend & a defendant / a defender & defence & defensive \\
\hline
to witness & a witness & witness (act of) & — \\
\hline
to counsel & a counsellor & counsel / counselling & — \\
\hline
\end{tabular}
\end{center}
\emph{Note} : \eng{judgment} (orthographe juridique standard) et \eng{judgement} (usage courant) sont acceptés. \eng{Defence} (nom) s'écrit avec un \eng{c} en anglais britannique (\eng{defense} en anglais américain).
\end{propriete}

\courssub{Idiomatic expressions and false friends}

\begin{exemplebox}[Expressions idiomatiques utiles]
\eng{to settle out of court} = régler à l'amiable (hors tribunal).\\
\eng{to reach a compromise / an agreement} = parvenir à un compromis / un accord.\\
\eng{to drop the charges} = laisser tomber les charges / abandonner les poursuites.\\
\eng{to take someone to court} = traîner quelqu'un en justice.\\
\eng{to give evidence} = déposer, témoigner (synonyme de \eng{to testify}).\\
\eng{to hold a hearing} = tenir une audience.
\end{exemplebox}

\begin{attention}[To assist : faux ami célèbre]
\eng{To assist} ne veut PAS dire « assister à » (être présent) !
\begin{itemize}
\item \eng{to attend a meeting / a hearing} = assister à une réunion / une audience : \eng{The elders attended the family meeting.}
\item \eng{to assist somebody} = aider quelqu'un : \eng{The counsellor assisted the young couple.}
\end{itemize}
De même, \eng{to assist at} n'existe pas en anglais standard ; dites toujours \eng{to attend}.
\end{attention}

\begin{attention}[Damage vs Damages]
\begin{itemize}
\item \eng{Damage} (singulier, indénombrable) = dommage, dégâts : \eng{The storm caused extensive damage.}
\item \eng{Damages} (pluriel, terme juridique) = dommages-intérêts : \eng{The plaintiff was awarded 500,000 francs CFA in damages.}
\end{itemize}
\end{attention}

\courssub{Traditional way or legal way?}

\begin{popfigure}
\begin{tikzpicture}[
box/.style={draw=popblue, fill=popblueL, rounded corners=3pt, align=center, text width=3.4cm, minimum height=1.1cm, font=\small},
bigbox/.style={draw=poporange, fill=poporangeL, rounded corners=3pt, align=center, text width=3.2cm, minimum height=1cm, font=\small\bfseries},
fleche/.style={-{Stealth[length=3mm]}, thick, popdark}
]
\node[bigbox] (d) at (0,0) {DISPUTE};
\node[box, align=center] (t) at (-3.6,-2.4){Traditional way\\ elders, chief,\\ family council,\\ palace};
\node[box, align=center] (l) at (3.6,-2.4){Legal way\\ lawyer, court,\\ judge, tribunal};
\node[bigbox] (s) at (0,-5) {SETTLEMENT};
\draw[fleche] (d) -- (t);
\draw[fleche] (d) -- (l);
\draw[fleche] (t) -- (s);
\draw[fleche] (l) -- (s);
\end{tikzpicture}
\legende{Two paths to settlement : la voie traditionnelle et la voie judiciaire convergent vers le règlement du conflit.}
\end{popfigure}

\begin{savaistu}[Two systems side by side]
Au Cameroun anglophone comme dans de nombreux pays africains anglophones, le droit coutumier (\eng{customary law}) coexiste avec le droit moderne hérité de la colonisation britannique (\eng{common law}). Beaucoup de conflits de voisinage, de mariage ou de terres sont d'abord portés devant le \eng{traditional ruler} ou le \eng{village council} ; ce n'est qu'en cas d'échec que l'affaire va devant un \eng{court of law}. En anglais, on dit : \eng{The matter was settled out of court.} = L'affaire a été réglée à l'amiable (hors du tribunal).
\end{savaistu}

\courssub{Pronunciation and spelling focus}

\begin{center}
\begin{tabular}{|l|l|l|}
\hline
\rowcolor{popblueL} \textbf{Word} & \textbf{Prononciation} & \textbf{Piège} \\
\hline
mediator & « MI-di-é-teur » & accent sur la 1\textsuperscript{re} syllabe \\
\hline
arbitrator & « AR-bi-tré-teur » & accent sur \emph{ar-} \\
\hline
witness & « WIT-ness » & \emph{w} se prononce « ou » \\
\hline
plaintiff & « PLEIN-tif » & orthographe difficile : \emph{-tiff} \\
\hline
defendant & « di-FEN-dant » & accent sur \emph{-fen-} \\
\hline
counsellor & « KAUN-se-leur » & double \emph{l} en anglais britannique \\
\hline
lawsuit & « LO-suite » & composé \emph{law} + \emph{suit} \\
\hline
\end{tabular}
\end{center}

\begin{exoresolu}[Complete with the right actor or place]
Énoncé — Complétez chaque phrase avec le mot correct choisi dans la liste : \eng{mediator, witness, plaintiff, customary court, lawyer, defendant}.
\begin{enumerate}
\item The \ldots\ldots\ldots told the judge exactly what he had seen that night.
\item The \ldots\ldots\ldots complained that his neighbour had taken his land, so he went to court.
\item The \ldots\ldots\ldots helped the two brothers talk calmly, but he did not take any decision.
\item Land and marriage disputes are often heard at the \ldots\ldots\ldots in the village.
\item The \ldots\ldots\ldots said he was innocent and hired a \ldots\ldots\ldots to defend him.
\end{enumerate}
\tcblower
Solution détaillée :
\begin{enumerate}
\item \eng{witness} — c'est le témoin qui dit ce qu'il a vu ; collocation : \eng{a witness testifies / tells the court}.
\item \eng{plaintiff} — celui qui se plaint et qui va au tribunal est le plaignant.
\item \eng{mediator} — il aide à dialoguer sans décider ; s'il décidait, ce serait un \eng{arbitrator}.
\item \eng{customary court} — les affaires de terres et de mariage relèvent du tribunal coutumier.
\item \eng{defendant} puis \eng{lawyer} — l'accusé se défend ; l'avocat le défend (\eng{to defend a client}).
\end{enumerate}
\end{exoresolu}

\begin{exercice}[Your turn]
Traduisez en anglais :
\begin{enumerate}
\item Les anciens ont entendu les deux parties à la chefferie.
\item Elle a décidé de poursuivre son voisin en justice.
\item Le médiateur a proposé une solution, mais la décision de l'arbitre est contraignante.
\item Trois témoins ont témoigné devant le tribunal.
\end{enumerate}
\end{exercice}

\begin{corrige}[Exercice « Your turn »]
\begin{enumerate}
\item \eng{The elders heard both parties at the palace.}
\item \eng{She decided to sue her neighbour.}
\item \eng{The mediator proposed a solution, but the arbitrator's decision is binding.}
\item \eng{Three witnesses testified before the court.}
\end{enumerate}
\end{corrige}

\begin{aretenir}[L'essentiel de la section]
\begin{itemize}
\item Les acteurs : \eng{mediator} (propose), \eng{arbitrator} (décide, \eng{binding decision}), \eng{judge, lawyer, witness, elder, chief, family head, counsellor}.
\item Les parties : \eng{plaintiff / complainant} contre \eng{defendant} ; \eng{the parties concerned}.
\item Les lieux : \eng{court of law, tribunal, customary court, village council, family meeting, palace} (chefferie).
\item Les verbes : \eng{to mediate, to arbitrate, to try / to judge, to hear a case, to testify, to sue (for), to accuse (of), to defend}.
\item Familles de mots : \eng{mediate/mediator/mediation}, \eng{arbitrate/arbitrator/arbitration}, \eng{negotiate/negotiator/negotiation}, \eng{judge/judgment/judicial}, \eng{sue/lawsuit}, \eng{defend/defendant/defence}, \eng{counsel/counsellor/counselling}.
\item Expressions : \eng{settle out of court, reach a compromise, drop the charges, take to court, give evidence}.
\item Pièges : \eng{to assist} = aider (≠ assister à = \eng{to attend}) ; \eng{palace} = chefferie en contexte africain ; \eng{damage} (dégâts) vs \eng{damages} (dommages-intérêts) ; \eng{accuse of} / \eng{sue for}.
\end{itemize}
\end{aretenir}

\coursec{The settlement process: key verbs and collocations}

Dans la section précédente, nous avons rencontré les \cle{acteurs} du règlement des conflits — the \eng{chief}, the \eng{elders}, the \eng{judge}, the \eng{mediator} — et les \cle{lieux} où ils interviennent — the \eng{palace}, the \eng{village square}, the \eng{court}. Mais que font exactement ces acteurs ? Quelles actions enchaînent-ils pour transformer un conflit en paix ? C'est précisément l'objet de cette section : les \cle{verbes} et les \cle{collocations} qui décrivent le \eng{settlement process}, c'est-à-dire le processus de règlement, étape par étape.

\courssub{Observing the process: a short text}

\begin{textebox}[Peace returns to Mankon village]
After hours of negotiation, both families agreed to compromise. They finally reached an agreement and shook hands in front of the elders.

The dispute had started three months earlier, when Mr. Che's goats destroyed part of Mrs. Ngong's maize farm. Mrs. Ngong demanded compensation, but Mr. Che refused to pay. Angry words were exchanged, and the two families stopped speaking to each other.

The village chief first invited both parties to discuss the matter calmly. When the discussion failed, the elders helped them to negotiate. Each side had to make concessions: Mr. Che agreed to pay half of the damage, and Mrs. Ngong agreed to forgive him. The chief then ruled that the case was closed. “Let us bury the hatchet and restore peace in our quarter,” he said. The two neighbours apologised to each other, made up, and reconciliation was complete. The whole village celebrated that evening.
\end{textebox}

\textbf{Glossaire}

\begin{center}
\begin{tabularx}{\linewidth}{|l|l|X|l|}
\hline
\rowcolor{popblueL} \textbf{Mot anglais} & \textbf{Nature} & \textbf{Traduction française} & \textbf{Prononciation (approx.)} \\
\hline
to compromise & verb & transiger, faire un compromis & « kom-pro-maïz » \\
\hline
to reach an agreement & expression & aboutir à un accord & « ri-tch an a-gri-ment » \\
\hline
compensation & noun & compensation, dédommagement & « kom-pen-seï-chn » \\
\hline
to rule & verb & statuer, décider (autorité) & « roul » \\
\hline
to bury the hatchet & idiom & enterrer la hache de guerre & « be-ri ze hâ-tchit » \\
\hline
to make up & phrasal verb & se réconcilier & « meïk ape » \\
\hline
reconciliation & noun & réconciliation & « ré-kon-si-li-é-chn » \\
\hline
\end{tabularx}
\end{center}

\textbf{Comprehension questions}

\begin{enumerate}
\item What started the dispute between the two families?
\item Who helped the parties to negotiate?
\item What compromise did each side make?
\item What did the chief rule at the end?
\end{enumerate}

\courssub{The five stages of the settlement process}

Observez le texte : le règlement ne se fait pas d'un coup. Il suit des \cle{étapes} ordonnées, chacune portée par un verbe précis.

\begin{popfigure}
\begin{tikzpicture}[node distance=0.35cm]
\tikzset{stage/.style={rectangle, rounded corners=3pt, draw=popblue, fill=popblueL, text=popdark, align=center, minimum height=0.9cm, minimum width=2.15cm, font=\small\bfseries}}
\node[stage, fill=poporangeL, draw=poporange] (s1) {CONFLICT};
\node[stage, right=of s1] (s2) {DISCUSSION};
\node[stage, right=of s2] (s3) {NEGOTIATION};
\node[stage, right=of s3] (s4) {COMPROMISE};
\node[stage, right=of s4] (s5) {AGREEMENT};
\node[stage, below=0.6cm of s3, fill=popgreenL, draw=popgreen] (s6) {RECONCILIATION};
\draw[-{Stealth}, thick, popblue] (s1) -- (s2);
\draw[-{Stealth}, thick, popblue] (s2) -- (s3);
\draw[-{Stealth}, thick, popblue] (s3) -- (s4);
\draw[-{Stealth}, thick, popblue] (s4) -- (s5);
\draw[-{Stealth}, thick, popgreen] (s5) -- (s6);
\end{tikzpicture}
\legende{Les étapes du règlement d'un conflit : du conflit à la réconciliation.}
\end{popfigure}

\begin{definition}[The key verbs of each stage]
\begin{itemize}
\item \eng{to discuss} (the matter) \emph{(« di-skeuss »)} : discuter du problème, première étape informelle. \eng{The elders asked both parties to discuss the matter calmly.} « Les anciens ont demandé aux deux parties de discuter calmement de l'affaire. »
\item \eng{to negotiate} \emph{(« né-go-chi-eït ») } : négocier, chercher un accord par des concessions mutuelles. \eng{The two families negotiated for three hours.} « Les deux familles ont négocié pendant trois heures. »
\item \eng{to compromise} / \eng{to make a compromise} \emph{(« kom-pro-maïz »)} : transiger, faire un compromis. \eng{Neither side won everything; both had to compromise.} « Aucune des parties n'a tout obtenu ; toutes deux ont dû transiger. »
\item \eng{to reach an agreement} \emph{(« ri-tch an a-gri-ment »)} : aboutir à un accord. \eng{They finally reached an agreement at midnight.} « Ils sont finalement parvenus à un accord à minuit. »
\item \eng{to reconcile} / \eng{to be reconciled} \emph{(« re-kon-saïl »)} : (se) réconcilier. \eng{The two brothers were reconciled after the chief's decision.} « Les deux frères se sont réconciliés après la décision du chef. »
\end{itemize}
\end{definition}

\courssub{Essential collocations: the verb and its partner noun}

En anglais, un verbe ne se combine pas avec n'importe quel nom : il a des \cle{partenaires privilégiés}. C'est ce qu'on appelle une \cle{collocation}. Voici les collocations essentielles du thème.

\begin{center}
\begin{tabularx}{\linewidth}{|X|X|X|}
\hline
\rowcolor{popblueL} \textbf{Collocation anglaise} & \textbf{Traduction française} & \textbf{Exemple en contexte} \\
\hline
to settle a dispute & régler un différend & \eng{They settled the dispute peacefully.} « Ils ont réglé le différend pacifiquement. » \\
\hline
to resolve a conflict & résoudre un conflit & \eng{The mediator resolved the conflict.} « Le médiateur a résolu le conflit. » \\
\hline
to solve a problem & résoudre un problème & \eng{Dialogue solves many problems.} « Le dialogue résout bien des problèmes. » \\
\hline
to reach an agreement & aboutir à un accord & \eng{Both parties reached an agreement.} « Les deux parties sont parvenues à un accord. » \\
\hline
to make a compromise & faire un compromis & \eng{Each side made a compromise.} « Chaque partie a fait un compromis. » \\
\hline
to find a solution & trouver une solution & \eng{The elders found a solution.} « Les anciens ont trouvé une solution. » \\
\hline
to restore peace & rétablir la paix & \eng{The chief restored peace in the village.} « Le chef a rétabli la paix dans le village. » \\
\hline
to bury the hatchet & enterrer la hache de guerre & \eng{They buried the hatchet after ten years.} « Ils ont enterré la hache de guerre après dix ans. » \\
\hline
\end{tabularx}
\end{center}

\begin{propriete}[Collocations figées : le bon verbe avec le bon nom]
Une collocation est une combinaison \cle{figée} : on ne peut pas intervertir les verbes librement.
\begin{itemize}
\item On dit \eng{to SETTLE a dispute} et \eng{to RESOLVE a conflict}, rarement l'inverse.
\item On dit \eng{to REACH an agreement}, jamais \eng{to arrive an agreement} (calque du français « arriver à un accord » ; on peut dire \eng{to arrive AT an agreement}, plus formel, mais \eng{reach} est la forme standard).
\item On dit \eng{to MAKE a compromise} et \eng{to MAKE concessions}, jamais \eng{to do a compromise}.
\item On dit \eng{to FIND a solution}, pas \eng{to find a resolution} dans le langage courant.
\end{itemize}
\eng{They settled the dispute peacefully.} « Ils ont réglé le différend pacifiquement. »
\end{propriete}

\begin{methode}[Comment mémoriser les collocations]
\begin{enumerate}
\item N'apprends jamais un mot isolé : apprends le \cle{verbe avec son nom partenaire} (\eng{settle + dispute}, \eng{reach + agreement}, \eng{make + compromise}).
\item Crée une \cle{phrase personnelle} pour chaque paire, liée à ta vie : \eng{My father settled a land dispute in our quarter last year.}
\item Regroupe les paires par étape du processus (discussion, négociation, accord, réconciliation) et récite la chaîne complète.
\item Teste-toi en cachant le verbe : « to \ldots\ a dispute ? » → \eng{settle}.
\end{enumerate}
\end{methode}

\courssub{Verbs of decision: when an authority rules}

Quand le règlement est confié à une autorité (chef, juge), d'autres verbes apparaissent : les \cle{verbes de décision}.

\begin{definition}[Verbs of decision]
\begin{itemize}
\item \eng{to rule (that...)} \emph{(« roul »)} : statuer, décider officiellement. \eng{The chief ruled that the land belonged to the family.} « Le chef a statué que la terre appartenait à la famille. »
\item \eng{to give a verdict} \emph{(« guiv a vé-rdikt »)} : rendre un verdict (décision sur la culpabilité). \eng{The jury gave its verdict: not guilty.} « Le jury a rendu son verdict : non coupable. »
\item \eng{to pass a sentence} \emph{(« passe a sé-ntenç ») } : prononcer une peine. \eng{The judge passed a sentence of two years in prison.} « Le juge a prononcé une peine de deux ans de prison. »
\item \eng{to fine somebody} \emph{(« faïn sam-ba-di ») } : infliger une amende à quelqu'un. \eng{The judge fined him 50,000 francs.} « Le juge lui a infligé une amende de 50 000 francs. »
\item \eng{to order compensation} \emph{(« or-der kom-pen-seï-chn ») } : ordonner un dédommagement. \eng{The court ordered compensation for the victim.} « Le tribunal a ordonné un dédommagement pour la victime. »
\end{itemize}
\end{definition}

\begin{attention}[Faux amis : sentence, verdict, demand]
\begin{itemize}
\item « Une sentence » au sens de peine de prison = \eng{a sentence} ; mais le \cle{verdict} est la décision de culpabilité (\eng{guilty} / \eng{not guilty}), pas la peine. Le juge \eng{gives a verdict} d'abord, puis \eng{passes a sentence}.
\item \eng{to demand} = \cle{exiger} (ton fort), et non « demander » au sens poli. « Demander » = \eng{to ask (for)}. \eng{She asked for compensation politely.} « Elle a demandé poliment un dédommagement. » / \eng{He demanded an apology immediately.} « Il a exigé des excuses immédiatement. »
\item \eng{compensation} = dédommagement (argent) ; ne pas confondre avec « compensation » au sens psychologique.
\end{itemize}
\end{attention}

\courssub{Verbs of reconciliation: making peace again}

La dernière étape est la plus humaine : se pardonner et renouer les liens.

\begin{definition}[Verbs of reconciliation]
\begin{itemize}
\item \eng{to forgive (somebody for something)} \emph{(« for-giv » — irrégulier : forgive, forgave, forgiven)} : pardonner (à quelqu'un pour quelque chose). \eng{She forgave him for his mistake.} « Elle lui a pardonné son erreur. »
\item \eng{to apologise} (GB) / \eng{to apologize} (US) \emph{(« a-po-lo-jaïz ») } : s'excuser. \eng{He apologised for his behaviour.} « Il s'est excusé pour son comportement. »
\item \eng{to make up (with somebody)} \emph{(« meïk ape »)} : se réconcilier (avec quelqu'un). \eng{The two friends made up after their quarrel.} « Les deux amis se sont réconciliés après leur querelle. »
\item \eng{to shake hands} \emph{(« cheïk hændz ») } : se serrer la main. \eng{They shook hands in front of the elders.} « Ils se sont serré la main devant les anciens. »
\item \eng{to mend fences} (idiom) \emph{(« mend fen-siz ») } : renouer des relations, réparer les liens. \eng{After the election, the two politicians tried to mend fences.} « Après l'élection, les deux hommes politiques ont tenté de renouer des relations. »
\end{itemize}
\end{definition}

\begin{attention}[Orthographe et grammaire des verbes de réconciliation]
\begin{itemize}
\item \eng{apologise} est l'orthographe britannique, \eng{apologize} l'américaine ; les deux sont correctes, mais sois cohérent dans une même copie. Le nom est toujours \eng{an apology} (une seule graphie).
\item \eng{to apologise} est suivi de \eng{for} + nom ou \emph{-ing} : \eng{He apologised for being late.} « Il s'est excusé d'être en retard. » Jamais \eng{to apologise his behaviour}.
\item \eng{to forgive} est irrégulier : \eng{forgive — forgave — forgiven}.
\item \eng{to make up} est inséparable quand il signifie « se réconcilier » : \eng{They made up} (pas *\eng{They made it up}).
\end{itemize}
\end{attention}

\courssub{From verbs to nouns: building the word family}

Chaque verbe du processus a un \cle{nom correspondant}. Les connaître permet de varier les formulations à l'écrit (style plus formel).

\begin{center}
\begin{tabularx}{\linewidth}{|X|X|X|l|}
\hline
\rowcolor{popgreenL} \textbf{Verbe} & \textbf{Nom} & \textbf{Traduction du nom} & \textbf{Prononciation du nom} \\
\hline
to settle & a settlement & un règlement (d'un différend) & « setteul-ment » \\
\hline
to resolve & a resolution & une résolution & « ré-zo-lou-chn » \\
\hline
to agree & an agreement & un accord & « a-gri-ment » \\
\hline
to compromise & a compromise & un compromis & « kom-pro-maïz » \\
\hline
to give a verdict & a verdict & un verdict & « vé-rdikt » \\
\hline
to fine & a fine & une amende & « faïn » \\
\hline
to compensate & compensation & un dédommagement & « kom-pen-seï-chn » \\
\hline
to apologise & an apology & des excuses & « a-po-lo-dji » \\
\hline
to reconcile & reconciliation & la réconciliation & « ré-kon-si-li-é-chn » \\
\hline
\end{tabularx}
\end{center}

Compare : \eng{They reached an agreement.} « Ils sont parvenus à un accord. » et \eng{The agreement was signed by both families.} « L'accord a été signé par les deux familles. » Le passage du verbe au nom rend la phrase plus formelle — très utile dans une composition.

\courssub{Practice}

\begin{exoresolu}[Complete with the correct collocation]
Complète chaque phrase avec le verbe qui convient, à la forme correcte : \emph{settle, reach, make, apologise, fine, restore}.

\begin{enumerate}
\item The elders helped the two families to \ldots\ the dispute peacefully.
\item After long negotiations, they finally \ldots\ an agreement.
\item Each party had to \ldots\ a compromise.
\item The judge \ldots\ the trader 25,000 francs for damaging his neighbour's fence.
\item The young man \ldots\ to his mother for his rude words.
\item The chief's decision \ldots\ peace in the quarter.
\end{enumerate}
\tcblower
\textbf{Solution détaillée}
\begin{enumerate}
\item \eng{settle} — collocation figée : \eng{to settle a dispute}. « Les anciens ont aidé les deux familles à régler le différend pacifiquement. »
\item \eng{reached} — \eng{to reach an agreement} (passé : \eng{reached}). « Après de longues négociations, ils sont finalement parvenus à un accord. »
\item \eng{make} — \eng{to make a compromise}. « Chaque partie a dû faire un compromis. »
\item \eng{fined} — \eng{to fine somebody + montant}. « Le juge a infligé au commerçant une amende de 25 000 francs. »
\item \eng{apologised} — \eng{to apologise to somebody for something}. « Le jeune homme s'est excusé auprès de sa mère pour ses paroles grossières. »
\item \eng{restored} — \eng{to restore peace}. « La décision du chef a rétabli la paix dans le quartier. »
\end{enumerate}
\end{exoresolu}

\begin{exercice}[Put the process in order and speak]
\begin{enumerate}
\item Remets les étapes suivantes dans l'ordre logique du processus de règlement : \eng{to reconcile — to discuss — to reach an agreement — to negotiate — to compromise}.
\item Pour chaque étape, écris une phrase complète au passé en utilisant une collocation de la leçon, dans le contexte d'un conflit de terrain dans ton village.
\item À l'oral, raconte le règlement complet en cinq phrases, en utilisant \eng{first, then, after that, finally}.
\end{enumerate}
\end{exercice}

\begin{corrige}[Exercice : Put the process in order]
\begin{enumerate}
\item Ordre correct : \eng{to discuss → to negotiate → to compromise → to reach an agreement → to reconcile}.
\item Exemple de production : \eng{First, the two farmers discussed the land problem with the chief. Then, they negotiated for several hours. After that, each farmer made a compromise on the boundary. Finally, they reached an agreement and were reconciled in front of the whole village.} « D'abord, les deux agriculteurs ont discuté du problème foncier avec le chef. Ensuite, ils ont négocié pendant plusieurs heures. Après cela, chaque agriculteur a fait un compromis sur la limite. Finalement, ils sont parvenus à un accord et se sont réconciliés devant tout le village. »
\item Production orale libre : vérifier la présence des connecteurs chronologiques et des collocations correctes (\eng{settle a dispute}, \eng{reach an agreement}, \eng{make a compromise}).
\end{enumerate}
\end{corrige}

\begin{savaistu}[Justice traditionnelle et justice moderne au Cameroun]
Au Cameroun, le règlement des conflits repose souvent sur une \cle{double juridiction} : la justice coutumière (chef de village, notables, \eng{customary court}) et la justice moderne (tribunaux de première instance, \eng{high court}). Dans les zones rurales, plus de 80\% des litiges fonciers et familiaux sont d'abord portés devant le chef ou le \eng{council of elders}. La procédure y est orale, basée sur la \eng{mediation} et la \eng{reconciliation} plutôt que sur la peine. L'expression \eng{“bury the hatchet”} trouve tout son sens ici : l'objectif n'est pas de punir mais de \eng{restore social harmony}. Cependant, tout justiciable peut faire appel de la décision coutumière devant les tribunaux étatiques. Cette coexistence est consacrée par la Constitution et la loi n°2006/015 sur l'organisation judiciaire.
\end{savaistu}

\begin{aretenir}[L'essentiel de la section]
\begin{itemize}
\item Le processus suit cinq étapes : \eng{discuss → negotiate → compromise → reach an agreement → reconcile}.
\item Les collocations sont figées : \eng{settle a dispute}, \eng{resolve a conflict}, \eng{reach an agreement}, \eng{make a compromise}, \eng{restore peace}.
\item Décision d'une autorité : \eng{rule, give a verdict, pass a sentence, fine somebody, order compensation}.
\item Réconciliation : \eng{forgive — forgave — forgiven}, \eng{apologise for}, \eng{make up with}, \eng{shake hands}, \eng{bury the hatchet}, \eng{mend fences}.
\item Chaque verbe a son nom : \eng{a settlement, an agreement, a verdict, a fine, an apology, reconciliation}.
\end{itemize}
\end{aretenir}

\coursec{Idiomatic expressions and false friends}

In the previous section, we studied the key verbs and collocations of the settlement process (\eng{to lodge a complaint}, \eng{to mediate a dispute}, \eng{to reach an agreement}...). But when Cameroonians — and English speakers in general — talk about conflicts and peace in everyday conversation, they rarely use only plain, literal vocabulary. They use \cle{idioms}: fixed expressions whose meaning cannot be guessed from the individual words. They also fall into the trap of \cle{false friends}, words that look like French words but mean something different. This section will train you to master both.

\courssub{Observing idioms in context}

Read the following short newspaper-style report. Pay attention to the expressions in bold: can you guess their meaning from the context?

\begin{textebox}[Peace returns to Mankon and Chomba — a community report]
After years of being \textbf{at loggerheads} over a piece of farmland, the villages of Mankon and Chomba have finally \textbf{buried the hatchet}. For almost a decade, every meeting between the two traditional councils ended in shouting, and careless speeches on both sides only \textbf{added fuel to the fire}. Some elders even \textbf{had a bone to pick with} the local administrator, whom they accused of favouring one side.

Last month, however, a respected mediator from Bamenda invited both delegations to \textbf{clear the air}. Speaking frankly for the first time, the two chiefs admitted that the quarrel had cost both communities too much: blocked roads, abandoned farms and broken marriages between families. “It is time to \textbf{turn over a new leaf},” declared Chief Fomuki of Chomba. “Our children deserve peace.”

After three days of negotiation, the two sides agreed to \textbf{meet each other halfway}: the disputed land will be shared, with a common market built on the boundary line. “Peace is always a question of \textbf{give and take},” the mediator explained. “Nobody wins everything. What matters is that both communities have found \textbf{common ground}.” On Saturday, the two villages celebrated the agreement with a common meal — a powerful symbol that they had truly \textbf{made peace}.
\end{textebox}

\textbf{Glossaire :}

\begin{center}
\begin{tabularx}{\linewidth}{|l|l|X|}
\hline
\rowcolor{popblueL} \textbf{English} & \textbf{Nature} & \textbf{Français} \\
\hline
a quarrel & noun & une querelle, une dispute \\
\hline
to favour & verb & favoriser \\
\hline
frankly & adverb & franchement \\
\hline
a boundary line & noun & une limite, une frontière \\
\hline
to deserve & verb & mériter \\
\hline
to share & verb & partager \\
\hline
\end{tabularx}
\end{center}

\textbf{Comprehension questions:}
\begin{enumerate}
\item What were the two villages quarrelling about?
\item What did the mediator ask the delegations to do first?
\item How was the dispute finally settled? Quote one idiom from the text in your answer.
\end{enumerate}

\begin{corrige}[Comprehension questions]
\begin{enumerate}
\item They were quarrelling about a piece of farmland.
\item He invited both delegations to clear the air — that is, to speak frankly and openly about the problem.
\item The dispute was settled by sharing the land and building a common market on the boundary line: the two sides agreed \eng{to meet each other halfway} and found \eng{common ground}.
\end{enumerate}
\end{corrige}

\courssub{Idioms of conflict}

\begin{definition}[Key idioms of conflict]
\begin{itemize}
\item \eng{to be at loggerheads (with somebody)} = to disagree strongly and for a long time (être en désaccord profond, être à couteaux tirés). \emph{Prononciation : « bi at lo-gueur-rède ».}
\item \eng{to add fuel to the fire} = to make a bad situation worse (jeter de l'huile sur le feu). \emph{Prononciation : « ad fiou-èl tou de faïeur ».}
\item \eng{to have a bone to pick with somebody} = to have a complaint or a grievance to discuss with somebody (avoir un compte à régler avec quelqu'un). \emph{Prononciation : « rav è bone tou pic ouiz ».}
\end{itemize}
\end{definition}

\begin{exemplebox}[Idioms of conflict in sentences]
\eng{The two families have been at loggerheads since the land dispute began.} « Les deux familles sont à couteaux tirés depuis le début du conflit foncier. »

\eng{Spreading rumours only adds fuel to the fire.} « Répandre des rumeurs ne fait que jeter de l'huile sur le feu. »

\eng{I have a bone to pick with you: why did you tell the chief I was responsible?} « J'ai un compte à régler avec toi : pourquoi as-tu dit au chef que j'étais responsable ? »
\end{exemplebox}

\begin{savaistu}[Where does “a bone to pick” come from?]
The expression \eng{to have a bone to pick with you} comes from the image of two dogs fighting over a single bone: each one wants to “pick” the meat off it, and neither will give up. It is extremely common in everyday British and American conversation — much more common than its formal equivalent \eng{to have a grievance}. Use it in informal speech, not in an official letter!
\end{savaistu}

\courssub{Idioms of peace and reconciliation}

\begin{definition}[To bury the hatchet]
\eng{To bury the hatchet} means \emph{to end a quarrel and become friendly again} (enterrer la hache de guerre). The expression comes from a Native American custom: when two tribes made peace, their chiefs buried a hatchet (une hache) in the ground as a symbol that the war was over. \emph{Prononciation : « bé-ri de ra-tchit ».}
\end{definition}

\begin{propriete}[Other idioms of peace]
\begin{itemize}
\item \eng{to make peace (with somebody)} = to end a conflict and restore friendly relations (faire la paix). Note the verb: we \emph{make} peace, never \emph{do} peace.
\item \eng{to turn over a new leaf} = to change your behaviour and start a better life (tourner la page, prendre un nouveau départ). A “leaf” here is a page of a book, not a leaf of a tree!
\item \eng{to clear the air} = to discuss a problem openly in order to remove tension and bad feelings (apaiser les tensions, mettre les choses au clair).
\end{itemize}
\end{propriete}

\begin{exemplebox}[Idioms of peace in sentences]
\eng{After ten years of hostility, the two villages finally buried the hatchet.} « Après dix ans d'hostilité, les deux villages ont finalement enterré la hache de guerre. »

\eng{Let's clear the air: tell me exactly what happened at the meeting.} « Apaisons les tensions : dis-moi exactement ce qui s'est passé à la réunion. »

\eng{He apologised and promised to turn over a new leaf.} « Il s'est excusé et a promis de tourner la page. »

\eng{The two brothers made peace at their father's funeral.} « Les deux frères ont fait la paix aux funérailles de leur père. »
\end{exemplebox}

\courssub{Idioms of compromise}

\begin{definition}[Key idioms of compromise]
\begin{itemize}
\item \eng{to meet somebody halfway} = to accept part of what the other person wants in order to reach an agreement (couper la poire en deux, faire des concessions mutuelles).
\item \eng{give and take} = a situation in which both sides make concessions (des concessions réciproques ; « donnant-donnant »). Note: this idiom is used as a \emph{noun phrase}.
\item \eng{to find common ground} = to discover opinions or interests that both sides share (trouver un terrain d'entente).
\end{itemize}
\end{definition}

\begin{exemplebox}[Idioms of compromise in sentences]
\eng{They agreed to meet each other halfway.} « Ils ont accepté de couper la poire en deux. »

\eng{Don't add fuel to the fire — try to find common ground instead!} « Ne jette pas d'huile sur le feu — essaie plutôt de trouver un terrain d'entente ! »

\eng{A successful negotiation is always a matter of give and take.} « Une négociation réussie est toujours une affaire de concessions réciproques. »
\end{exemplebox}

\begin{attention}[Idioms are fixed expressions]
An idiom cannot be changed freely. Say \eng{to bury the hatchet}, never \emph{to bury the axe} or \emph{to hide the hatchet}. Say \eng{to add fuel to the fire}, never \emph{to add petrol to the fire}. If you change the words, native speakers will not understand you — or will smile! Learn each idiom as a complete block, exactly as it is.
\end{attention}

\courssub{False friends: words that betray you}

A \cle{false friend} (\eng{a false cognate}) is an English word that looks like a French word but has a different meaning. In the field of disputes and meetings, four false friends are particularly dangerous for francophone learners.

\begin{definition}[Four dangerous false friends]
\begin{itemize}
\item \eng{to demand} = to ask for something forcefully, as a right (\textbf{exiger}). It is much stronger than French \emph{demander}. To translate \emph{demander}, use \eng{to ask (for)} or, more formally, \eng{to request}.
\item \eng{an issue} = a problem, a topic under discussion (\textbf{un problème, une question}). It never means \emph{une issue} in the sense of an exit: an exit is \eng{an exit} or \eng{a way out}.
\item \eng{to achieve} = to succeed in doing something after effort (\textbf{accomplir, réaliser, atteindre}). It does not mean \emph{achever} in the sense of finishing: to finish is \eng{to finish} or \eng{to complete}.
\item French \emph{une déception} is not English \emph{a deception}. The correct word is \eng{a disappointment} (une déception). \eng{Deception} in English means \emph{tromperie, duperie}!
\end{itemize}
\end{definition}

\begin{attention}[“To demand” is much stronger than “demander”]
\eng{I demand an apology.} = « J'exige des excuses. » — This is aggressive and confrontational. If you simply want to ask politely, say \eng{I would like to ask for an explanation} or \eng{Could I request a meeting with the chief?} Using \eng{to demand} in a polite request can sound rude and may damage a negotiation!
\end{attention}

\begin{exemplebox}[False friends used correctly]
\eng{The demonstrators demanded the release of their leader.} « Les manifestants ont exigé la libération de leur chef. »

\eng{The boundary issue was discussed for three hours.} « La question de la limite a été discutée pendant trois heures. »

\eng{The mediators achieved remarkable results in Buea.} « Les médiateurs ont accompli des résultats remarquables à Buea. »

\eng{To everyone's disappointment, the talks failed.} « À la déception de tous, les pourparlers ont échoué. »
\end{exemplebox}

\begin{popfigure}
\begin{center}
\begin{tikzpicture}[
  box/.style={rectangle, rounded corners=2pt, draw=popline, fill=poppinkL, text width=4.2cm, align=center, font=\small, minimum height=1.1cm},
  corr/.style={rectangle, rounded corners=2pt, draw=popline, fill=popgreenL, text width=5.2cm, align=center, font=\small, minimum height=1.1cm},
  arr/.style={-{Stealth[length=2.5mm]}, thick, popdark}
]
\node[box, align=center] (f1) at (0,0){\textbf{to demand}\\ \emph{looks like} « demander »};
\node[corr, align=center] (c1) at (7.2,0){\textbf{= to insist forcefully (exiger)}\\ polite: \eng{to ask for / to request}};
\draw[arr] (f1) -- (c1);

\node[box, align=center] (f2) at (0,-1.7){\textbf{an issue}\\ \emph{looks like} « une issue »};
\node[corr, align=center] (c2) at (7.2,-1.7){\textbf{= a problem, a topic (un problème)}\\ « une issue/sortie » = \eng{an exit}};
\draw[arr] (f2) -- (c2);

\node[box, align=center] (f3) at (0,-3.4){\textbf{to achieve}\\ \emph{looks like} « achever »};
\node[corr, align=center] (c3) at (7.2,-3.4){\textbf{= to accomplish (accomplir)}\\ « achever/finir » = \eng{to finish}};
\draw[arr] (f3) -- (c3);

\node[box, align=center] (f4) at (0,-5.1){\textbf{a deception}\\ \emph{looks like} « une déception »};
\node[corr, align=center] (c4) at (7.2,-5.1){\textbf{= a trick (une tromperie)}\\ « une déception » = \eng{a disappointment}};
\draw[arr] (f4) -- (c4);
\end{tikzpicture}
\end{center}
\legende{False friends: what the English word really means, and how to correct the mistake.}
\end{popfigure}

\courssub{Calques du français à éviter}

Beyond single words, francophone learners often translate whole French structures word for word. These \cle{calques} produce incorrect or unnatural English.

\begin{center}
\begin{tabularx}{\linewidth}{|X|X|X|}
\hline
\rowcolor{popblueL} \textbf{Calque du français (incorrect)} & \textbf{Correct English} & \textbf{Français} \\
\hline
to assist a meeting & to attend a meeting & assister à une réunion \\
\hline
to make a dispute & to have a dispute & avoir un différend \\
\hline
to do peace & to make peace & faire la paix \\
\hline
to demand an information & to ask for information & demander une information \\
\hline
to arrange a conflict & to settle / to resolve a conflict & régler un conflit \\
\hline
to take a decision (acceptable but less natural) & to make a decision & prendre une décision \\
\hline
\end{tabularx}
\end{center}

\begin{attention}[“Assister à” is not “to assist”]
French \emph{assister à une réunion} means \eng{to attend a meeting} (être présent). English \eng{to assist} means \textbf{aider} : \eng{The secretary assisted the mediator.} « La secrétaire a aidé le médiateur. » Similarly, \emph{régler un conflit} is \eng{to settle a dispute}, never \emph{to arrange a conflict} (\eng{to arrange} = organiser, arranger). Note also that \eng{information} is uncountable: never \emph{an information}, say \eng{some information} or \eng{a piece of information}.
\end{attention}

\courssub{Summary table: idioms at a glance}

\begin{center}
\begin{tabularx}{\linewidth}{|X|X|}
\hline
\rowcolor{popblueL} \textbf{English idiom} & \textbf{French equivalent} \\
\hline
to be at loggerheads & être à couteaux tirés, en désaccord profond \\
\hline
to add fuel to the fire & jeter de l'huile sur le feu \\
\hline
to have a bone to pick with somebody & avoir un compte à régler avec quelqu'un \\
\hline
to bury the hatchet & enterrer la hache de guerre \\
\hline
to make peace & faire la paix \\
\hline
to turn over a new leaf & tourner la page, prendre un nouveau départ \\
\hline
to clear the air & apaiser les tensions, mettre les choses au clair \\
\hline
to meet somebody halfway & couper la poire en deux \\
\hline
give and take & concessions réciproques, donnant-donnant \\
\hline
to find common ground & trouver un terrain d'entente \\
\hline
\end{tabularx}
\end{center}

\begin{exoresolu}[Guided practice: idioms and false friends]
\textbf{Exercise.} Complete each sentence with the correct idiom or word from the box, in the correct form: \eng{to bury the hatchet} — \eng{to add fuel to the fire} — \eng{to meet halfway} — \eng{an issue} — \eng{to attend} — \eng{a disappointment}.

\begin{enumerate}
\item After twenty years of quarrelling, the two families finally decided to \dots
\item Please stop shouting; you are only \dots
\item The water \dots{} was discussed at the council meeting.
\item Both sides had to \dots{} to reach an agreement.
\item Did you \dots{} the reconciliation meeting in Douala last week?
\item The failure of the talks was a great \dots{} for the whole community.
\end{enumerate}

\tcblower
\textbf{Solution détaillée.}
\begin{enumerate}
\item \eng{to bury the hatchet} — contexte : fin d'une longue querelle (enterrer la hache de guerre). La structure \eng{decided to} est suivie de l'infinitif : \eng{decided to bury the hatchet}.
\item \eng{adding fuel to the fire} — le présent continu est nécessaire après \eng{you are} : crier aggrave la situation (jeter de l'huile sur le feu).
\item \eng{issue} — faux ami : « la question de l'eau » = \eng{the water issue} ; jamais \emph{the water exit}!
\item \eng{meet (each other) halfway} — après \eng{had to}, infinitif : les deux parties ont dû couper la poire en deux.
\item \eng{attend} — calque à éviter : « assister à une réunion » = \eng{to attend a meeting}, jamais \emph{to assist a meeting}.
\item \eng{disappointment} — faux ami : « une déception » = \eng{a disappointment} ; \eng{a deception} signifierait « une tromperie ».
\end{enumerate}
\end{exoresolu}

\begin{exercice}[Your turn]
Translate into English, using an idiom or the correct word:
\begin{enumerate}
\item Ne jette pas d'huile sur le feu !
\item Les deux chefs ont enterré la hache de guerre après la médiation.
\item J'ai un compte à régler avec mon voisin.
\item Ils ont finalement trouvé un terrain d'entente sur la question des terres.
\item J'exige des excuses ! (utiliser \eng{to demand})
\end{enumerate}
\end{exercice}

\begin{corrige}[Exercice « Your turn »]
\begin{enumerate}
\item \eng{Don't add fuel to the fire!}
\item \eng{The two chiefs buried the hatchet after the mediation.} (passé : \eng{buried})
\item \eng{I have a bone to pick with my neighbour.}
\item \eng{They finally found common ground on the land issue.} (faux ami : \eng{issue} = question, problème)
\item \eng{I demand an apology!}
\end{enumerate}
\end{corrige}

\begin{aretenir}[What you must remember]
\begin{itemize}
\item Idioms are \textbf{fixed blocks}: learn \eng{to bury the hatchet}, \eng{to add fuel to the fire}, \eng{to meet somebody halfway} exactly as they are, with their exact words.
\item \eng{To demand} = exiger (strong!); to ask politely = \eng{to ask for}, \eng{to request}.
\item \eng{An issue} = un problème ; \eng{to achieve} = accomplir ; « une déception » = \eng{a disappointment}.
\item « Assister à » = \eng{to attend} ; « régler un conflit » = \eng{to settle a dispute} ; « faire la paix » = \eng{to make peace}.
\end{itemize}
\end{aretenir}

\coursec{Word families, spelling and pronunciation}

\courssub{Transition : des expressions figées à la construction du mot}

Dans la section précédente, nous avons exploré les \emph{idioms} et les \emph{faux amis} qui colorent le discours sur les conflits communautaires. Maîtriser ces blocs figés est indispensable pour la compréhension et l'expression naturelle. Cependant, pour atteindre l'autonomie lexicale requise en classe de Première — notamment pour l'épreuve de \emph{Word Formation} au Baccalauréat — il faut comprendre la \textbf{morphologie dérivationnelle} : comment une racine unique génère verbes, noms, adjectifs et adverbes grâce aux préfixes et suffixes. C'est l'objet de cette section : observer les familles de mots, fixer l'orthographe et verrouiller la prononciation.

\courssub{The \cle{AGREE} family : from consensus to conflict}

Commençons par l'observation d'un court extrait de procès-verbal de médiation à Buea.

\begin{textebox}{Minutes of a mediation session (South-West Region)}
The \textbf{agreement} was reached after three hours of tense discussion. Both parties \textbf{agreed} to respect the traditional boundaries. However, a minor \textbf{disagreement} arose regarding the cocoa farm access road. The mediator found the proposal \textbf{agreeable}, but one elder deemed it \textbf{disagreeable} because it ignored ancestral customs. Finally, they \textbf{agreed to disagree} on that specific point and signed the document.
\end{textebox}

\begin{definition}[Glossaire]
\begin{itemize}
\item \textbf{agreement} (n) : accord ; \textbf{to agree} (v) : être d'accord ; \textbf{disagreement} (n) : désaccord ; \textbf{agreeable} (adj) : agréable, acceptable ; \textbf{disagreeable} (adj) : désagréable, inacceptable ; \textbf{to agree to disagree} (expr.) : convenir de ne pas être d'accord.
\end{itemize}
\end{definition}

\begin{propriete}[Famille morphologique de AGREE]
\emph{Racine} : \cle{agree} (verbe de base).
\begin{center}
\begin{tabularx}{\linewidth}{|X|X|X|X|}
\hline
\rowcolor{popblueL}
\textbf{Forme} & \textbf{Catégorie} & \textbf{Sens} & \textbf{Exemple traduit} \\
\hline
agree & verbe & consentir, être d'avis conforme & \eng{They agree on the terms.} « Ils sont d'accord sur les termes. » \\
\hline
agreement & nom (concret/abstrait) & accord, entente & \eng{The agreement is binding.} « L'accord est contraignant. » \\
\hline
disagree & verbe (préfixe \cle{dis-}) & être en désaccord & \eng{I disagree with the verdict.} « Je suis en désaccord avec le verdict. » \\
\hline
disagreement & nom & désaccord, dispute & \eng{A minor disagreement delayed the signature.} « Un léger désaccord a retardé la signature. » \\
\hline
agreeable & adjectif (suffixe \cle{-able}) & acceptable, plaisant & \eng{The settlement was agreeable to all.} « Le règlement était acceptable pour tous. » \\
\hline
disagreeable & adjectif (préfixe \cle{dis-} + \cle{-able}) & inacceptable, désagréable & \eng{His tone was disagreeable.} « Son ton était désagréable. » \\
\hline
\end{tabularx}
\end{center}
\end{propriete}

\begin{popfigure}
\begin{tikzpicture}[
  mindmap,
  concept color=popblue!60,
  level 1/.style={level distance=3.5cm, sibling angle=90},
  level 2/.style={level distance=3cm, sibling angle=45},
  every node/.style={font=\small, text=popdark},
  edge from parent/.style={draw=popblue!50, thick}
]
\node[concept, text width=2.5cm, align=center] {AGREE\\ (root)}
  child[concept color=popgreen!60] {node[concept, text width=2.2cm, align=center] {VERB\\ agree\\ disagree}}
  child[concept color=poporange!60] {node[concept, text width=2.2cm, align=center] {NOUN\\ agreement\\ disagreement}}
  child[concept color=poppurple!60] {node[concept, text width=2.2cm, align=center] {ADJECTIVE\\ agreeable\\ disagreeable}}
  child[concept color=poppink!60] {node[concept, text width=2.2cm, align=center] {ADVERB\\ agreeably\\ disagreeably}};
\end{tikzpicture}
\legende{Arbre dérivationnel de la famille AGREE. Notez la productivité du préfixe \texttt{dis-} et du suffixe \texttt{-able}.}
\end{popfigure}

\begin{exemplebox}{Collocations clés (AGREE)}
\begin{itemize}
\item \eng{to reach an agreement} (parvenir à un accord) — \textbf{pas} *make an agreement.
\item \eng{to agree on / upon sth} (s'accorder sur qch).
\item \eng{to agree with sb} (être d'accord avec qn).
\item \eng{agreeable to both parties} (acceptable pour les deux parties).
\end{itemize}
\end{exemplebox}

\courssub{The \cle{RECONCILE} family : restoring harmony}

\begin{textebox}{Community radio announcement (Bamenda)}
After the land dispute that divided the quarter for months, the \textbf{reconciliation} ceremony took place yesterday at the fon's palace. The rival families were \textbf{reconciled} by the traditional council. The elders stressed that no offence is \textbf{irreconcilable} if dialogue is sincere. A \textbf{reconciliatory} gesture — the sharing of kola nuts — sealed the peace.
\end{textebox}

\begin{definition}[Glossaire]
\begin{itemize}
\item \textbf{reconciliation} (n) : réconciliation ; \textbf{to reconcile} (v) : réconcilier ; \textbf{reconcilable} (adj) : réconciliable ; \textbf{irreconcilable} (adj, préfixe \cle{ir-}) : irréconciliable ; \textbf{reconciliatory} (adj) : de réconciliation, conciliant.
\end{itemize}
\end{definition}

\begin{propriete}[Famille morphologique de RECONCILE]
\emph{Racine} : \cle{reconcile} (verbe). Attention au \textbf{e} final muet qui tombe devant \emph{-iation} et \emph{-iatory}.
\begin{center}
\begin{tabularx}{\linewidth}{|X|X|X|X|}
\hline
\rowcolor{popblueL}
\textbf{Forme} & \textbf{Catégorie} & \textbf{Formation} & \textbf{Exemple traduit} \\
\hline
reconcile & verbe & base & \eng{They reconciled their differences.} « Ils ont aplani leurs différends. » \\
\hline
reconciliation & nom & -iation (chgt y $\to$ i) & \eng{Reconciliation takes time.} « La réconciliation prend du temps. » \\
\hline
reconcilable & adjectif & -able & \eng{The issues are reconcilable.} « Les questions sont réconciliables. » \\
\hline
irreconcilable & adjectif & ir- + -able & \eng{Their positions remain irreconcilable.} « Leurs positions restent irréconciliables. » \\
\hline
reconciliatory & adjectif & -iatory & \eng{A reconciliatory speech.} « Un discours de réconciliation. » \\
\hline
\end{tabularx}
\end{center}
\end{propriete}

\begin{attention}[Piège orthographique : \textbf{reconciliation} vs *reconcilliation]
En français, on écrit « réconciliation » avec deux \textbf{L}. En anglais, la racine \emph{reconcile} n'a qu'un \textbf{L} $\to$ \textbf{reconciliation} (un seul L). De même : \emph{reconciliatory}. \textbf{N'écrivez jamais} *reconcilliation.
\end{attention}

\begin{methode}[Word Formation : transformer RECONCILE]
\emph{Consigne type Bac :} \eng{Complete the sentence with the correct form of the word in brackets.}
\begin{enumerate}
\item Identifier la catégorie manquante (nom, adjectif, verbe, adverbe).
\item \eng{The \_\_\_\_\_\_ (RECONCILE) process was slow.} $\to$ Il faut un adjectif $\to$ \textbf{reconciliatory} (ou \emph{reconciling} comme participe présent adjectival).
\item \eng{We hope for a lasting \_\_\_\_\_\_ (RECONCILE).} $\to$ Il faut un nom après \emph{a lasting} $\to$ \textbf{reconciliation}.
\item Vérifier l'orthographe : pas de double L, passage \emph{y} $\to$ \emph{i} devant \emph{-ation/-atory}.
\end{enumerate}
\end{methode}

\courssub{The \cle{JUDGE} family : authority and decision}

\begin{textebox}{Court report (Yaoundé High Court)}
The \textbf{judge} listened carefully to both counsels. His \textbf{judgement} was delivered at noon. He urged the plaintiff not to be \textbf{judgmental} but to accept the compensation. The defendant's lawyer argued that the previous ruling was a \textbf{misjudgement}. The \textbf{adjudication} of land disputes often requires customary law expertise.
\end{textebox}

\begin{definition}[Glossaire]
\begin{itemize}
\item \textbf{judge} (n personne / v) : juge / juger ; \textbf{judgement} (n chose) : jugement, décision ; \textbf{judgmental} (adj) : qui juge hâtivement, critique ; \textbf{misjudgement} (n) : erreur de jugement ; \textbf{adjudication} (n) : adjudication, décision judiciaire.
\end{itemize}
\end{definition}

\begin{propriete}[Famille morphologique de JUDGE]
L'orthographe britannique \textbf{judgement} (avec \emph{e}) est la norme au Cameroun et dans le Commonwealth. L'américain \emph{judgment} (sans \emph{e}) est aussi accepté en anglais international juridique. Le verbe \emph{to judge} garde son \emph{e}.
\begin{center}
\begin{tabularx}{\linewidth}{|X|X|X|}
\hline
\rowcolor{popblueL}
\textbf{Anglais} & \textbf{Français} & \textbf{Note} \\
\hline
a judge & un juge (personne) & nom d'agent \\
to judge & juger & verbe \\
judgement & un jugement (décision) & nom d'action ; orthographe GB standard au Cameroun \\
judgmental & critique, moralisateur & adj péjoratif : \eng{Don't be so judgmental!} \\
misjudgement & erreur de jugement & préfixe \cle{mis-} (mauvais) \\
adjudication & décision d'un arbitre/tribunal & préfixe \cle{ad-} + racine \emph{judic-} \\
\hline
\end{tabularx}
\end{center}
\end{propriete}

\begin{attention}[Le piège \textbf{ARGUMENT} — Erreur n°1 des francophones]
Le verbe est \textbf{to argue} (pas de \emph{e} final). Le nom est \textbf{argument} (\textbf{sans e} après le \emph{g}).
\begin{itemize}
\item \eng{They had a heated argument.} (Correct)
\item \eng{*They had a heated arguement.} (\textbf{INCORRECT} — calque sur le français « argument » ou confusion avec \emph{judgement}).
\end{itemize}
Autres dérivés : \emph{arguable} (adj), \emph{arguably} (adv).
\end{attention}

\begin{exoresolu}{Word Formation (JUDGE family)}
\emph{Complétez avec la forme correcte du mot entre parenthèses.}
\begin{enumerate}
\item The village chief acted as an impartial \_\_\_\_\_\_ (JUDGE) during the dispute.
\item His \_\_\_\_\_\_ (JUDGE) was questioned by the losing party.
\item It was a serious \_\_\_\_\_\_ (JUDGE) to underestimate the tension.
\item Avoid being too \_\_\_\_\_\_ (JUDGE) when listening to both sides.
\end{enumerate}
\tcblower
\textbf{Solution détaillée :}
\begin{enumerate}
\item \textbf{judge} (nom personne, sujet de \emph{acted}). \eng{The village chief acted as an impartial judge.}
\item \textbf{judgement} (nom chose, possessif \emph{His}). \eng{His judgement was questioned...}
\item \textbf{misjudgement} (nom, après \emph{a serious}, préfixe \emph{mis-} pour exprimer l'erreur). \eng{It was a serious misjudgement...}
\item \textbf{judgmental} (adjectif après \emph{being too}, décrit une attitude). \eng{Avoid being too judgmental...}
\end{enumerate}
\end{exoresolu}

\courssub{The \cle{PEACE} family : the ultimate goal}

\begin{textebox}{Excerpt from a peace treaty preamble (Far North Region)}
\eng{We, the representatives of the farming and herding communities, commit to lasting \textbf{peace}. We shall live \textbf{peacefully} side by side. The traditional authorities will \textbf{pacify} any rising tension immediately. A \textbf{peaceful} coexistence benefits our children and our cattle.}
\end{textebox}

\begin{definition}[Glossaire]
\begin{itemize}
\item \textbf{peace} (n) : paix ; \textbf{peaceful} (adj) : paisible, pacifique ; \textbf{peacefully} (adv) : paisiblement ; \textbf{to pacify} (v) : apaiser, pacifier.
\end{itemize}
\end{definition}

\begin{propriete}[Famille morphologique de PEACE]
\begin{center}
\begin{tabularx}{\linewidth}{|X|X|X|X|}
\hline
\rowcolor{popblueL}
\textbf{Forme} & \textbf{Catégorie} & \textbf{Suffixation} & \textbf{Exemple traduit} \\
\hline
peace & nom (base) & — & \eng{Peace returned to the village.} « La paix est revenue au village. » \\
\hline
peaceful & adjectif & \cle{-ful} (plein de) & \eng{A peaceful resolution.} « Une résolution pacifique. » \\
\hline
peacefully & adverbe & \cle{-ful} + \cle{-ly} & \eng{They settled the matter peacefully.} « Ils ont réglé l'affaire paisiblement. » \\
\hline
pacify & verbe & racine latine \emph{pac-} + \cle{-ify} & \eng{The army was sent to pacify the region.} « L'armée fut envoyée pour pacifier la région. » \\
\hline
pacification & nom & \cle{-ication} & \eng{The pacification took years.} « La pacification a pris des années. » \\
\hline
\end{tabularx}
\end{center}
\end{propriete}

\begin{aretenir}[Distinction \textbf{peaceful} vs \textbf{pacific}]
\begin{itemize}
\item \textbf{Peaceful} : qui aime la paix, sans violence (\eng{a peaceful protest}, \eng{a peaceful night}). C'est l'adjectif courant.
\item \textbf{Pacific} : 1. Qui apporte la paix (\eng{pacific intentions}) — formel. 2. Relatif à l'océan Pacifique (\eng{the Pacific Ocean}). \textbf{Pas} *a pacific protest.
\end{itemize}
\end{aretenir}

\courssub{The \cle{NEGOTIATE} family : the art of compromise}

\begin{textebox}{News article (The Cameroon Tribune - Business)}
\eng{The \textbf{negotiation} between the cocoa cooperative and the exporters lasted three days. The chief \textbf{negotiator} insisted on a minimum price per ton. Both sides showed flexibility; the final price was \textbf{negotiable}. A deadlock was avoided because neither party adopted a \textbf{non-negotiable} stance from the start.}
\end{textebox}

\begin{definition}[Glossaire]
\begin{itemize}
\item \textbf{negotiation} (n) : négociation ; \textbf{to negotiate} (v) : négocier ; \textbf{negotiator} (n personne) : négociateur ; \textbf{negotiable} (adj) : négociable ; \textbf{non-negotiable} (adj) : non négociable.
\end{itemize}
\end{definition}

\begin{propriete}[Famille morphologique de NEGOTIATE]
\emph{Racine} : \cle{negotiate}. Le \textbf{e} final tombe devant \emph{-ion} et \emph{-or}.
\begin{center}
\begin{tabularx}{\linewidth}{|X|X|X|X|}
\hline
\rowcolor{popblueL}
\textbf{Forme} & \textbf{Catégorie} & \textbf{Formation} & \textbf{Exemple traduit} \\
\hline
negotiate & verbe & base & \eng{They negotiate in good faith.} « Ils négocient de bonne foi. » \\
\hline
negotiation & nom & -ation (e $\to$ Ø) & \eng{The negotiation failed.} « La négociation a échoué. » \\
\hline
negotiator & nom agent & -or (e $\to$ Ø) & \eng{A skilled negotiator.} « Un négociateur habile. » \\
\hline
negotiable & adjectif & -able (e $\to$ Ø) & \eng{The terms are negotiable.} « Les termes sont négociables. » \\
\hline
non-negotiable & adjectif & non- + -able & \eng{Human rights are non-negotiable.} « Les droits de l'homme ne se négocient pas. » \\
\hline
\end{tabularx}
\end{center}
\end{propriete}

\begin{propriete}[Règles d'orthographe récapitulatives (Suffixation)]
\begin{center}
\begin{tabularx}{\linewidth}{|X|X|X|}
\hline
\rowcolor{popblueL}
\textbf{Règle} & \textbf{Mots du thème} & \textbf{Attention} \\
\hline
\cle{-ment} : pas de changement (sauf \emph{argue} $\to$ \emph{argument}) & \textbf{agreement, settlement, judgement, adjudgement} & *arguement $\times$ \\
\hline
\cle{-tion / -sion} : \emph{e} final tombe ; \emph{te} $\to$ \emph{tion} ; \emph{de/se} $\to$ \emph{sion} & \textbf{negotiation, reconciliation, mediation, arbitration, pacification} & *negociation $\times$ (pas de \emph{e}) ; *reconcilliation $\times$ (1 seul \emph{l}) \\
\hline
\cle{-able / -ible} : \emph{e} final tombe généralement & \textbf{negotiable, agreeable, reconcilable, enforceable} & *negotiable $\checkmark$ ; *agreeable $\checkmark$ \\
\hline
\cle{Double consonne} (GB) : verbe 1 syllabe, voyelle courte + consonne finale $\to$ double au \emph{-ed/-ing} & \textbf{quarrel $\to$ quarrelled (GB), quarreled (US)} ; \textbf{settle $\to$ settled} (pas de double, \emph{e} muet) & \textbf{refer $\to$ referred} (accent sur dernière syllabe) \\
\hline
\end{tabularx}
\end{center}
\end{propriete}

\courssub{Prononciation : guide en lettres françaises}

La prononciation anglaise ne se devine pas à l'écrit. Voici les transcriptions \textbf{approchées en lettres françaises} pour les mots-clés de cette leçon. La syllabe en \textbf{MAJUSCULES} porte l'accent tonique principal.

\begin{center}
\begin{tabularx}{\linewidth}{|X|X|X|}
\hline
\rowcolor{popblueL}
\textbf{Mot (IPA simplifiée)} & \textbf{Prononciation (lettres FR)} & \textbf{Conseil} \\
\hline
reconcile /ˈrɛkənsaɪl/ & \textbf{RÉ}-kon-saïl & \emph{cile} se lit « saïl » (comme \emph{profile}). Pas « ré-con-si-lé ». \\
\hline
negotiation /nɪˌɡəʊʃiˈeɪʃn/ & né-go-chi-\textbf{É}-chone & \emph{ti} = /ʃ/ (ch) ; \emph{a} long /eɪ/ (éï). 5 syllabes. \\
\hline
apology /əˈpɒlədʒi/ & a-\textbf{PO}-lo-dji & Accent sur \textbf{PO}. \emph{gy} = /dʒi/ (dji). Pas « a-po-lo-gui ». \\
\hline
compromise /ˈkɒmprəmaɪz/ & \textbf{KOM}-pro-maïze & \emph{pro} = /prə/ (peu) ; \emph{mise} = /maɪz/ (maïze). \\
\hline
mediation /ˌmiːdiˈeɪʃn/ & mé-di-\textbf{É}-chone & Même finale que \emph{negotiation}. \\
\hline
arbitration /ˌɑːbɪˈtreɪʃn/ & ar-bi-\textbf{TRÉ}-chone & \emph{tra} = /treɪ/ (traï). \\
\hline
judgement /ˈdʒʌdʒmənt/ & \textbf{DJADJ}-meunte & \emph{dge} = /dʒ/ (dj) ; \emph{ment} = /mənt/ (meunte). \\
\hline
\end{tabularx}
\end{center}

\begin{methode}[Entraînement à la prononciation (Méthode \cle{3R})]
\begin{enumerate}
\item \textbf{Repérage} : Souligne la syllabe accentuée dans ton cahier (ex. \underline{recon}\textbf{cile}).
\item \textbf{Répétition chorale} : En classe, répète 3 fois le mot isolé, puis dans une phrase courte (\eng{We need reconciliation.}).
\item \textbf{Regroupement par famille} : Lis à voix haute la colonne « Verbe » puis « Nom » puis « Adjectif » du tableau \emph{NEGOTIATE} : \eng{negotiate — negotiation — negotiator — negotiable}. L'accent bouge souvent : \eng{NEGotiate} $\to$ \eng{negoTIation} $\to$ \eng{neGOtiator} $\to$ \eng{neGOtiable}.
\end{enumerate}
\end{methode}

\courssub{Synthèse des préfixes et suffixes productifs}

\begin{propriete}[Boîte à outils morphologiques pour \cle{settling disputes}]
\begin{center}
\begin{tabularx}{\linewidth}{|X|X|X|X|}
\hline
\rowcolor{popblueL}
\textbf{Afixe} & \textbf{Sens / Fonction} & \textbf{Exemples du thème} & \textbf{Autres exemples} \\
\hline
\cle{dis-} & négation, opposition & \textbf{disagree, dispute, disarm} & dislike, dishonest \\
\hline
\cle{ir-} (devant r) & négation & \textbf{irreconcilable, irregular} & illegal (devant l), impossible (devant p/m) \\
\hline
\cle{un-} & négation (adj/participe) & \textbf{unfair, unresolved, unbiased} & unhappy, unknown \\
\hline
\cle{mis-} & mauvais, mal & \textbf{misjudgement, misunderstand} & mistake, mislead \\
\hline
\cle{non-} & non- (adj composé) & \textbf{non-negotiable, non-violent} & non-existent, non-profit \\
\hline
\cle{-ment} & nom d'action/résultat (V $\to$ N) & \textbf{agreement, settlement, judgement} & development, government \\
\hline
\cle{-tion / -sion} & nom d'action (V $\to$ N) & \textbf{negotiation, mediation, reconciliation, arbitration} & information, decision \\
\hline
\cle{-or / -er} & nom d'agent (personne) & \textbf{negotiator, mediator, arbitrator, lawyer, judge} & teacher, actor \\
\hline
\cle{-able / -ible} & adjectif « qui peut être... » (V $\to$ Adj) & \textbf{negotiable, agreeable, reconcilable, enforceable} & readable, visible \\
\hline
\cle{-ful} & « plein de » (N $\to$ Adj) & \textbf{peaceful, respectful, harmful} & useful, careful \\
\hline
\cle{-less} & « sans » (N $\to$ Adj) & \textbf{peaceless (rare), lawless, hopeless} & useless, careless \\
\hline
\end{tabularx}
\end{center}
\end{propriete}

\courssub{Exercice d'application intégré : Texte à trous (Word Formation)}

\begin{exercice}{Community Dispute Resolution — Word Formation}
\emph{Complétez le texte ci-dessous en utilisant la forme correcte du mot entre parenthèses. Écrivez vos réponses sur votre cahier.}

In the village of Mbankomo, a land (1) \_\_\_\_\_\_ (DISPUTE) between two families threatened the social (2) \_\_\_\_\_\_ (PEACE). The village chief, acting as a (3) \_\_\_\_\_\_ (MEDIATE), called for a (4) \_\_\_\_\_\_ (RECONCILE) meeting. He urged the parties to (5) \_\_\_\_\_\_ (NEGOTIATION) in good faith. After hours of discussion, a (6) \_\_\_\_\_\_ (SETTLE) was reached. The terms of the (7) \_\_\_\_\_\_ (AGREE) were (8) \_\_\_\_\_\_ (BIND) on both sides. One elder, however, remained (9) \_\_\_\_\_\_ (RECONCILE) to the outcome, claiming the process had been (10) \_\_\_\_\_\_ (FAIR). The chief warned that such an (11) \_\_\_\_\_\_ (JUDGE) attitude could (12) \_\_\_\_\_\_ (DANGER) the fragile peace.
\end{exercice}

\begin{corrige}{Exercice : Community Dispute Resolution}
\begin{enumerate}
\item \textbf{dispute} (nom, après \emph{a land}). Le verbe est \emph{dispute}, le nom s'écrit pareil.
\item \textbf{peace} (nom, après \emph{social}). L'adjectif \emph{peaceful} ne convient pas (pas de nom derrière).
\item \textbf{mediator} (nom agent, après \emph{a}). Suffixe \emph{-or}.
\item \textbf{reconciliation} (nom, après \emph{a}). Suffixe \emph{-tion} (attention : 1 seul L).
\item \textbf{negotiate} (verbe infinitif après \emph{to}). Pas \emph{negotiation} (nom).
\item \textbf{settlement} (nom, après \emph{a}). Suffixe \emph{-ment}.
\item \textbf{agreement} (nom, après \emph{the}). Suffixe \emph{-ment}.
\item \textbf{binding} (adjectif / participe présent, après \emph{were}). \emph{Bound} est le participe passé, mais \emph{binding} = contraignant (adj). \emph{Bind} est le verbe.
\item \textbf{irreconcilable} (adjectif, après \emph{remained}). Préfixe \emph{ir-} + \emph{reconcilable}. Sens : « qui ne peut être réconcilié / qui refuse la réconciliation ».
\item \textbf{unfair} (adjectif, après \emph{been}). Préfixe \emph{un-}.
\item \textbf{judgmental} (adjectif, après \emph{an}). Pas \emph{judgement} (nom). \emph{Judgmental} = critique, qui juge.
\item \textbf{endanger} (verbe, après modal \emph{could}). Verbe base. \emph{Danger} est un nom.
\end{enumerate}
\end{corrige}

\begin{savaistu}{Le Bac et la \cle{Word Formation}}
À l'épreuve écrite du Baccalauréat (série A), l'exercice de \textbf{Word Formation} (souvent 5 à 10 items) puise massivement dans les familles du programme : \emph{agree, judge, negotiate, reconcile, peace, settle, mediate, arbitrate}.
\begin{itemize}
\item \textbf{Stratégie gagnante :} Apprends les \textbf{triplets} (Verbe / Nom en \emph{-tion/-ment} / Adjectif en \emph{-able/-ful}) par cœur.
\item \textbf{Exemple classique :} \eng{The \_\_\_\_\_\_ (MEDIATE) proposed a fair solution.} $\to$ \textbf{mediator} (nom personne).
\item \textbf{Exemple classique :} \eng{The conflict was \_\_\_\_\_\_ (SOLVE) peacefully.} $\to$ \textbf{solved} (participe passé) ou \textbf{solvable} (adjectif) selon le contexte. Ici, passif \emph{was solved}.
\end{itemize}
Maîtriser l'orthographe de \emph{judgement}, \emph{argument}, \emph{reconciliation} et \emph{negotiation} te rapporte des points faciles.
\end{savaistu}

\begin{experience}[Atelier oral : \cle{The Peace Summit Simulation}]
\textbf{Objectif :} Réemployer le vocabulaire des familles \emph{AGREE, RECONCILE, NEGOTIATE, PEACE, JUDGE} en interaction orale.
\textbf{Durée :} 20 min.
\textbf{Déroulement :}
\begin{enumerate}
\item \textbf{Groupes de 4 :} 1 Médiateur, 2 Parties adverses (ex. éleveur / cultivateur), 1 Observateur/Note-taker.
\item \textbf{Scénario :} Un conflit d'accès à l'eau à la saison sèche. Le médiateur ouvre : \eng{``We are here to reach an \textbf{agreement}. Let us \textbf{negotiate} calmly.''}
\item \textbf{Contraintes linguistiques :} Chaque élève doit utiliser \textbf{au moins 3 mots dérivés} différents (ex. \emph{negotiator, negotiation, negotiable} ; \emph{reconcile, reconciliation, irreconcilable} ; \emph{judge, judgement, judgmental}).
\item \textbf{Rôle de l'observateur :} Coche les mots employés sur une grille. Note les erreurs de prononciation (ex. \emph{reconciliaTIon} vs \emph{reconCILIation}).
\item \textbf{Retour classe :} L'observateur fait un rapport : \eng{``The mediator used \emph{negotiation} and \emph{reconciliation} correctly. The farmer mispronounced \emph{judgement}.''}
\end{enumerate}
\end{experience}

\begin{aretenir}[Check-list de révision — Section 5]
\begin{itemize}
\item \textbf{AGREE family} : agree, agreement, disagree, disagreement, agreeable, disagreeable. \textbf{Colloc :} \emph{reach an agreement}.
\item \textbf{RECONCILE family} : reconcile, reconciliation (1 L), reconcilable, irreconcilable (ir-), reconciliatory. \textbf{Piège :} pas de double L.
\item \textbf{JUDGE family} : judge (n/v), judgement (GB standard), judgmental, misjudgement. \textbf{Piège :} \emph{argument} (pas de e).
\item \textbf{PEACE family} : peace, peaceful, peacefully, pacify, pacification. \textbf{Distinction :} peaceful $\neq$ pacific.
\item \textbf{NEGOTIATE family} : negotiate, negotiation, negotiator, negotiable, non-negotiable. \textbf{Règle :} e final tombe.
\item \textbf{Orthographe} : -ment (pas de changement), -tion (e tombe), -able (e tombe), double consonne GB (quarrelled).
\item \textbf{Prononciation} : reconcile (RÉ-kon-saïl), negotiation (né-go-chi-É-chone), apology (a-PO-lo-dji), compromise (KOM-pro-maïze).
\item \textbf{Affixes clés} : dis-, ir-, un-, mis-, non- / -ment, -tion, -or, -able, -ful.
\end{itemize}
\end{aretenir}

\coursec{Using the vocabulary in context}

Dans les sections précédentes, nous avons étudié les types de conflits, les acteurs du règlement, les verbes et collocations clés, les expressions idiomatiques, les faux amis ainsi que les familles de mots (\eng{dispute / disputant / disputable}), l'orthographe et la prononciation du lexique. Mais il ne suffit pas de \emph{connaître} les mots : il faut savoir les \cle{employer} dans des phrases, des dialogues et des textes. C'est l'objectif de cette dernière section : passer du vocabulaire \emph{reconnu} (à la lecture, à l'écoute) au vocabulaire \emph{réutilisé} (à l'oral et à l'écrit), comme l'exigent les épreuves de compréhension et de composition de la classe de Première.

\courssub{Reading: The Land Dispute in Mankon}

Lisons d'abord un court récit qui met en scène, dans un contexte camerounais, presque tout le lexique étudié dans cette leçon. Lisez le texte une première fois pour le sens général (\emph{skimming}), puis une seconde fois en soulignant les mots de la leçon (\emph{scanning}).

\begin{textebox}[The Land Dispute in Mankon]
For ten years, the Ndi and the Fong families had been at loggerheads over a piece of farmland behind the old market. The quarrel started when Mr Ndi planted plantains on a strip of land that the Fongs claimed as their own. Insults were exchanged, and one night, somebody destroyed the young plants. The two families stopped greeting each other, and the whole neighbourhood felt the tension.

Finally, the case was brought before the village council. The elders listened carefully to both parties, visited the land, and questioned two old witnesses who remembered the original boundaries. After a long discussion, they proposed a compromise: each family would farm half of the field, and a line of trees would mark the new boundary.

Both sides accepted the decision, apologised to each other, and buried the hatchet with a shared meal of eru and fufu. The chief praised the families for choosing dialogue instead of violence, and reminded everyone that a community that settles its disputes peacefully is a community that prospers.
\end{textebox}

\textbf{Glossaire}\par
\begin{itemize}
\item \eng{to be at loggerheads} (expression idiomatique) : être en conflit ouvert, « être à couteaux tirés »
\item \eng{a strip of land} (nom) : une bande de terrain
\item \eng{to claim} (verbe) : revendiquer, prétendre posséder
\item \eng{a boundary} (nom) : une limite, une frontière (entre deux terrains) ; prononciation « baoun-da-ri »
\item \eng{a witness} (nom) : un témoin ; prononciation « ouit-nèss »
\item \eng{to bury the hatchet} (idiome) : enterrer la hache de guerre, se réconcilier
\item \eng{to settle a dispute} (collocation) : régler un différend
\item \eng{to prosper} (verbe) : prospérer
\end{itemize}

\textbf{Comprehension questions}\par
\begin{enumerate}
\item What were the two families quarrelling about?
\item Who settled the dispute, and how?
\item What compromise did the elders propose?
\item Which idiom in the text means « to make peace »?
\end{enumerate}

\begin{corrige}[Comprehension questions]
\begin{enumerate}
\item They were quarrelling over a piece of farmland (a strip of land behind the old market).
\item The village council (the elders) settled it: they listened to both parties, visited the land and questioned two witnesses.
\item Each family would farm half of the field, and a line of trees would mark the new boundary.
\item \eng{To bury the hatchet}.
\end{enumerate}
\end{corrige}

\courssub{Listening and speaking: a model dialogue}

Voici maintenant le même lexique à l'oral, dans une médiation entre deux voisins devant le chef. Lisez le dialogue à voix haute en respectant la prononciation étudiée : \eng{dispute} se prononce « dis-pioute », \eng{compromise} « kom-pro-maïze », \eng{reconcile} « ré-kon-saïl », \eng{boundary} « baoun-da-ri ».

\begin{textebox}[Before the Chief]
Chief: I have called this meeting because your quarrel is dividing the whole quarter. Mr Ekane, speak first.

Mr Ekane: Chief, my neighbour built his fence two metres inside my plot. I want him to remove it.

Mr Bate: That is not true! The fence follows the old boundary. My father planted those trees himself.

Chief: Enough. I will not listen to accusations. I propose that we all go to the land tomorrow morning with two elders as witnesses. We shall measure the plots and check the documents.

Mr Ekane: I accept, Chief. I only want justice.

Mr Bate: I also accept. I am ready to compromise if I am wrong.

Chief: Good. Whatever we find, you will both shake hands and end this dispute like brothers. A fence should separate two plots, not two families.
\end{textebox}

\begin{exemplebox}[Phrases utiles du dialogue, traduites]
\eng{The chief called a meeting.} = Le chef a convoqué une réunion.\par
\eng{Both parties accepted the compromise.} = Les deux parties ont accepté le compromis.\par
\eng{I am ready to compromise.} = Je suis prêt à transiger, à faire un compromis.\par
\eng{You will both shake hands and end this dispute.} = Vous vous serrerez tous les deux la main et mettrez fin à ce différend.\par
\eng{A fence should separate two plots, not two families.} = Une clôture devrait séparer deux terrains, pas deux familles.
\end{exemplebox}

\begin{attention}[Orthographe : \eng{quarrel} au passé]
En orthographe britannique, le \emph{l} final de \eng{quarrel} double devant une terminaison : \eng{quarrelled}, \eng{quarrelling}. De même : \eng{travel} → \eng{travelled}, \eng{level} → \eng{levelled}. En anglais américain, on écrit \eng{quarreled} avec un seul \emph{l} : les deux formes existent, mais restez cohérent dans une même copie — au Cameroun, on privilégie la norme britannique.
\end{attention}

\courssub{From conflict to reconciliation: the three stages}

Tout récit de règlement de conflit suit la même architecture en trois temps. Mémorisez ce schéma : il vous servira de plan pour toutes vos productions écrites et orales.

\begin{popfigure}
\begin{center}
\begin{tikzpicture}[
box/.style={rectangle, rounded corners=4pt, draw=popblue, thick, fill=popblueL, text width=3.4cm, align=center, minimum height=2.8cm, font=\small},
arr/.style={-{Stealth[length=3mm]}, very thick, poporange}]
\node[box, align=center] (p) at (0,0){\textbf{PROBLEM}\\[2pt] quarrel over land\\ \emph{dispute, quarrel, tension, to be at loggerheads}};
\node[box, fill=poporangeL, draw=poporange, align=center] (a) at (5.2,0){\textbf{ACTION}\\[2pt] elders mediate, both sides compromise\\ \emph{to mediate, to listen to both parties, to propose a compromise}};
\node[box, fill=popgreenL, draw=popgreen, align=center] (r) at (10.4,0){\textbf{RESULT}\\[2pt] agreement, reconciliation\\ \emph{to settle the dispute, to apologise, to bury the hatchet}};
\draw[arr] (p) -- (a);
\draw[arr] (a) -- (r);
\end{tikzpicture}
\end{center}
\legende{Les trois étapes d'un récit de règlement de conflit, avec le lexique associé}
\end{popfigure}

\begin{methode}[How to write a conflict-settlement narrative]
Pour raconter un conflit et son règlement en 6 à 8 phrases, suivez ces cinq étapes :
\begin{enumerate}
\item \textbf{Introduce the parties and the cause.} Présentez les parties et la cause du conflit : \eng{The X and Y families had been at loggerheads over...}
\item \textbf{Describe the escalation.} Décrivez l'escalade : \eng{Insults were exchanged. The tension grew.}
\item \textbf{Name the mediator.} Nommez le médiateur : \eng{The case was brought before the village council / the chief.}
\item \textbf{Explain the solution.} Expliquez la solution : \eng{The elders proposed a compromise: ...}
\item \textbf{Conclude with reconciliation.} Concluez sur la réconciliation : \eng{Both sides apologised and buried the hatchet.}
\end{enumerate}
\end{methode}

\courssub{Guided practice: gap-fill and sentence transformation}

\begin{exoresolu}[Gap-fill with the lesson word bank]
\textbf{Énoncé.} Complétez avec : \eng{dispute — mediator — compromise — witnesses — settled — apologised}.
\begin{enumerate}
\item The two brothers had a \ldots\ldots over their father's land.
\item The chief acted as a \ldots\ldots between the two parties.
\item Three \ldots\ldots confirmed that the boundary was near the river.
\item After long talks, they reached a \ldots\ldots : each brother kept half of the field.
\item The case was finally \ldots\ldots last week.
\item The younger brother \ldots\ldots for his angry words.
\end{enumerate}
\tcblower
\textbf{Solution détaillée.}
\begin{enumerate}
\item \eng{dispute} — nom comptable précédé de l'article indéfini \eng{a} ; collocation \eng{to have a dispute over something}.
\item \eng{mediator} — le chef joue le rôle de médiateur ; structure \eng{to act as a + nom de fonction}.
\item \eng{witnesses} — trois témoins ; pluriel en \emph{-es} après \emph{-ss} (\eng{witness} → \eng{witnesses}).
\item \eng{compromise} — collocation : \eng{to reach a compromise}.
\item \eng{settled} — collocation : \eng{to settle a case / a dispute} ; ici au passif, prétérit : \eng{was settled}.
\item \eng{apologised} — \eng{to apologise for something} ; attention à la préposition \eng{for}, jamais *\eng{apologise his words}.
\end{enumerate}
\end{exoresolu}

\begin{exoresolu}[Sentence transformation with learned collocations]
\textbf{Énoncé.} Reformulez chaque phrase avec la collocation entre parenthèses, sans changer le sens.
\begin{enumerate}
\item They ended the conflict. (\eng{to resolve})
\item The two families made peace. (\eng{to bury the hatchet})
\item The chief helped them to find an agreement. (\eng{to mediate})
\item Both sides said yes to the solution. (\eng{to accept a compromise})
\end{enumerate}
\tcblower
\textbf{Solution détaillée.}
\begin{enumerate}
\item \eng{They resolved the conflict.}
\item \eng{The two families buried the hatchet.}
\item \eng{The chief mediated between them and helped them reach an agreement.}
\item \eng{Both sides accepted the compromise.}
\end{enumerate}
Remarquez que la reformulation avec une collocation apprise rend la phrase plus \emph{naturelle} et plus \emph{précise} : c'est exactement ce qu'attend l'examinateur à l'épreuve de composition.
\end{exoresolu}

\begin{attention}[Faux ami : « régler » n'est pas « to regulate »]
Ne traduisez jamais mot à mot « ils ont réglé le problème » par *\eng{they regulated the problem}. \eng{To regulate} signifie « réglementer » (par des lois, des règles : \eng{The government regulates prices.}). Pour « régler un problème / un conflit », dites : \eng{they solved the problem}, \eng{they settled the dispute} ou \eng{they resolved the conflict}.\par
Autre piège : « un compromis » se dit bien \eng{a compromise}, mais « compromettre » au sens de « mettre en danger » se dit \eng{to jeopardise} ou \eng{to compromise} uniquement suivi d'un nom abstrait (\eng{to compromise someone's reputation} = compromettre la réputation de quelqu'un), jamais pour dire « faire un compromis » sans complément : on dira \eng{to reach a compromise} ou \eng{to be ready to compromise}.
\end{attention}

\courssub{Guided writing: your own narrative}

\begin{center}
\begin{tabularx}{\linewidth}{|X|X|X|}
\hline
\rowcolor{popblueL} Étape & Modèle anglais & Français \\
\hline
Parties + cause & The Ndi and Fong families had been at loggerheads over a piece of farmland. & Les familles Ndi et Fong étaient à couteaux tirés à propos d'un terrain. \\
\hline
Escalade & Insults were exchanged and the tension grew. & Des insultes furent échangées et la tension s'aggrava. \\
\hline
Médiateur & The case was brought before the village council. & L'affaire fut portée devant le conseil du village. \\
\hline
Solution & The elders proposed a compromise: each family would farm half of the field. & Les anciens proposèrent un compromis : chaque famille cultiverait la moitié du champ. \\
\hline
Réconciliation & Both sides apologised and buried the hatchet. & Les deux parties s'excusèrent et enterrèrent la hache de guerre. \\
\hline
\end{tabularx}
\end{center}

\begin{exercice}[Your turn: write your narrative]
Racontez en 6 à 8 phrases un conflit (réel ou imaginaire) et son règlement, en suivant le plan en cinq étapes de la boîte méthode.\par
\textbf{Critères de réussite :}
\begin{itemize}
\item au moins \textbf{6 mots} du lexique de la leçon (\eng{dispute, quarrel, mediator, witness, boundary, compromise, reconciliation...}) ;
\item au moins \textbf{2 collocations} (\eng{to settle a dispute, to reach a compromise, to bring a case before...}) ;
\item au moins \textbf{1 idiome} (\eng{to be at loggerheads, to bury the hatchet}) ;
\item \textbf{aucun faux ami} (pas de *\eng{regulate the problem} !).
\end{itemize}
\end{exercice}

\begin{corrige}[Exercice de rédaction — production possible]
\eng{For five years, the Abang and the Tabi families had been at loggerheads over a cocoa farm near the river. The quarrel became so serious that the two families stopped speaking to each other. One day, the case was brought before the village chief, who acted as a mediator. He listened to both parties and called two old witnesses who knew the original boundaries. After a long discussion, the elders proposed a compromise: the farm would be divided into two equal parts. Both sides accepted the decision and apologised to each other. They finally settled the dispute and buried the hatchet with a shared meal. Since that day, peace has returned to the village.}\par
Vérification des critères : 6 mots du lexique (\eng{loggerheads, quarrel, mediator, witnesses, compromise, dispute}), 2 collocations (\eng{to propose a compromise, to settle the dispute}), 1 idiome (\eng{to bury the hatchet}), aucun faux ami.
\end{corrige}

\begin{aretenir}[L'essentiel de la section]
Un récit de règlement de conflit suit toujours trois temps : \textbf{problem} (la dispute) → \textbf{action} (la médiation, le compromis) → \textbf{result} (l'accord, la réconciliation). Pour bien réutiliser le lexique : employez les collocations apprises (\eng{settle a dispute, reach a compromise, bury the hatchet}), méfiez-vous des faux amis (\eng{regulate} ≠ « régler »), respectez l'orthographe britannique (\eng{quarrelled, apologised}) et structurez votre texte en cinq étapes : parties et cause, escalade, médiateur, solution, réconciliation.
\end{aretenir}

\newpage

\coursec{Activité d'intégration}

\courssub{Objectifs de l'activité}

À l'issue de cette activité d'intégration, tu seras capable de :

\begin{itemize}
\item réemployer \textbf{à l'oral} au moins huit mots du lexique du règlement des conflits communautaires (\eng{dispute, settle, mediate, compromise, witness, agreement, reconcile}...) dans une situation de communication authentique ;
\item utiliser correctement \textbf{trois collocations} clés (\eng{settle a dispute, reach an agreement, listen to both sides}...) et \textbf{une expression idiomatique} (\eng{bury the hatchet, meet halfway}...) ;
\item jouer un rôle social précis (\eng{chief, mediator, elder, witness, neighbour}) en respectant le registre de langue adapté ;
\item produire \textbf{à l'écrit} un court \eng{settlement report} de six phrases, structuré et rédigé au passé ;
\item évaluer la prestation des autres groupes et voter pour la médiation la plus convaincante.
\end{itemize}

\courssub{Situation}

Tu es élève en Première A au lycée de Mankon, près de Bamenda, dans la région du Nord-Ouest du Cameroun. Dans le village voisin, deux voisins, Mr Ndi et Mr Fomba, se disputent depuis des mois une bande de terre située entre leurs deux concessions. La tension monte : les deux familles ne se parlent plus et des voisins craignent que le conflit ne dégénère. Le chef du village décide alors de réunir le \eng{village council} (le conseil du village) pour régler le différend à l'amiable, selon la tradition, avant que l'affaire ne soit portée devant le tribunal.

Ta classe va reconstituer ce conseil. Lis d'abord la convocation officielle rédigée par le chef :

\begin{textebox}[Notice from the Village Chief — Mankon Traditional Council]
\textbf{PUBLIC NOTICE}

To all members of the Mankon community,

The Traditional Council will sit this Saturday at 10 a.m. under the big tree near the market square. The council will hear the land dispute between Mr Ndi and Mr Fomba, two neighbours of Mile 4 quarter.

Both parties will present their case. Witnesses and elders will be heard. The council expects everyone to speak calmly, to tell the truth, and to respect the mediator.

Our aim is not to punish anyone, but to settle the dispute peacefully, to help both sides reach an agreement, and to restore harmony in our community. After the hearing, the two neighbours will be invited to bury the hatchet and share the traditional kola nut.

Any person with useful information about the boundary between the two plots is welcome to testify.

Signed,
Chief Angwafor III
\end{textebox}

\begin{definition}[Glossaire utile]
\begin{itemize}
\item \eng{to sit} (v., ici) : siéger, tenir une séance — « The council will sit at 10 a.m. »
\item \eng{a hearing} (n.) : une audience, une audition des parties
\item \eng{a boundary} (n.) : une limite, une frontière (entre deux terrains)
\item \eng{a plot} (n.) : une parcelle de terrain
\item \eng{to testify} (v.) : témoigner, déposer
\item \eng{the kola nut} (n.) : la noix de cola (symbole de paix et de réconciliation)
\end{itemize}
\end{definition}

\courssub{Consignes}

Travaille en groupe de quatre ou cinq élèves. Suis les étapes dans l'ordre.

\begin{enumerate}
\item \textbf{Distribue les rôles} (2 minutes) : un élève joue le \eng{chief / mediator}, deux élèves jouent les voisins en conflit (\eng{the two parties}), un ou deux élèves jouent les \eng{elders} ou les \eng{witnesses}. Choisissez ensemble le type de conflit : \eng{land dispute} (dispute de terrain), \eng{noise} (bruit), ou \eng{inheritance} (héritage).

\item \textbf{Prépare le conflit} (5 minutes) : les deux voisins rédigent chacun trois phrases pour présenter leur version des faits (\eng{present your case}). Utilise le présent parfait et le prétérit : \eng{He has destroyed my fence. He planted cassava on my land last year.}

\item \textbf{Joue la médiation} (8 à 10 minutes) : le \eng{mediator} ouvre la séance, invite chaque partie à parler (\eng{listen to both sides}), pose des questions aux témoins, puis propose un compromis (\eng{propose a compromise}). Chaque groupe doit employer \textbf{au moins huit mots du lexique}, \textbf{trois collocations} et \textbf{un idiome} (\eng{bury the hatchet} ou \eng{meet halfway}).

\item \textbf{Conclus par un accord} : la séance se termine obligatoirement par un accord oral (\eng{reach an agreement}) et un geste de réconciliation (\eng{reconcile, shake hands, share the kola nut}).

\item \textbf{Rédige le rapport} (10 minutes, à l'écrit) : un secrétaire par groupe rédige le \eng{settlement report} en \textbf{six phrases exactement}, au prétérit, en suivant le plan : nature du conflit → version de chaque partie → témoignages → compromis proposé → accord final → réconciliation.

\item \textbf{Vote de la classe} : chaque groupe joue sa scène devant la classe. À la fin, votez pour la médiation la plus convaincante selon trois critères : richesse du lexique, correction de la langue, réalisme du compromis.
\end{enumerate}

\begin{attention}[Pendant le jeu de rôle]
Évite les calques du français : on ne dit pas \eng{*assist the meeting} pour « assister à la réunion » mais \eng{attend the meeting} ; on ne dit pas \eng{*make a compromise} mais \eng{propose / reach a compromise}. Le mot \eng{eventually} signifie « finalement », pas « éventuellement ».
\end{attention}

\courssub{Ressources à mobiliser}

\begin{aretenir}[Lexique et collocations à réemployer]
\begin{itemize}
\item \textbf{Noms} : \eng{a dispute, a conflict, a disagreement, a complaint, a boundary, a witness, an elder, a mediator, a settlement, an agreement, a compromise}
\item \textbf{Verbes} : \eng{to settle, to resolve, to mediate, to complain, to accuse, to apologize, to forgive, to reconcile, to compromise}
\item \textbf{Collocations} : \eng{settle a dispute, reach an agreement, listen to both sides, propose a compromise, restore peace, bear a grudge} (garder rancune)
\item \textbf{Idiomes} : \eng{bury the hatchet} (enterrer la hache de guerre), \eng{meet halfway} (faire des concessions mutuelles), \eng{clear the air} (apaiser les tensions)
\end{itemize}
\end{aretenir}

\begin{propriete}[Structures grammaticales utiles]
\begin{itemize}
\item \textbf{Prétérit} pour raconter les faits : \eng{Mr Ndi planted cassava on the disputed plot.}
\item \textbf{Present perfect} pour un fait avec conséquence présente : \eng{He has destroyed my fence.}
\item \textbf{Should / ought to} pour les recommandations du médiateur : \eng{Both parties should meet halfway.}
\item \textbf{Discours rapporté} dans le rapport : \eng{The witness said that the boundary was near the old mango tree.}
\end{itemize}
\end{propriete}

\courssub{Proposition de production et corrigé commenté}

\begin{exoresolu}[Modèle de settlement report]
Énoncé : à l'issue du jeu de rôle, rédige en six phrases le rapport de la médiation entre Mr Ndi et Mr Fomba (conflit de terrain).
\tcblower
\textbf{Modèle de production (en anglais) :}

\eng{Last Saturday, the village council met to settle a land dispute between Mr Ndi and Mr Fomba, two neighbours of Mile 4 quarter. Mr Ndi accused his neighbour of planting cassava on his plot, while Mr Fomba claimed that the land belonged to his late father. An elder, who had witnessed the original sharing of the land, testified that the boundary was near the old mango tree. The mediator listened to both sides carefully and proposed a compromise: each man would keep the part he had already cultivated. After a long discussion, the two parties reached an agreement and accepted the boundary proposed by the council. Finally, Mr Ndi and Mr Fomba shook hands, shared the kola nut and decided to bury the hatchet, so peace was restored in the community.}

\textbf{Commentaire des choix de langue :}
\begin{itemize}
\item \textbf{Lexique} : dix mots du champ lexical sont réemployés (\eng{settle, dispute, accused, witnessed, testified, boundary, mediator, compromise, agreement, peace}) — la consigne en exigeait huit.
\item \textbf{Collocations} : \eng{settle a land dispute, listen to both sides, propose a compromise, reach an agreement, restore peace} — cinq au lieu des trois exigées.
\item \textbf{Idiome} : \eng{bury the hatchet} clôt le récit, placé naturellement à la fin.
\item \textbf{Grammaire} : prétérit dominant (\eng{met, accused, testified, reached}) ; discours rapporté correct (\eng{testified that the boundary was...}, sans recul de temps nécessaire car le fait reste vrai) ; \eng{while} oppose les deux versions ; \eng{after + nom} et \eng{finally} structurent la chronologie.
\item \textbf{Format} : exactement six phrases, comme demandé.
\end{itemize}
\end{exoresolu}

\courssub{Bilan de l'activité}

Cette activité t'a permis de passer du lexique appris en classe à une \textbf{utilisation réelle en contexte} : tu as joué un rôle social authentique, proche des pratiques de médiation traditionnelle du Cameroun anglophone, et tu as transformé une performance orale en production écrite structurée. Pour vérifier ta réussite, demande-toi :

\begin{itemize}
\item Ai-je employé au moins huit mots du lexique et trois collocations sans calquer le français ?
\item Mon idiome était-il placé naturellement dans le discours ?
\item Mon rapport respecte-t-il les six phrases, le prétérit et le plan imposé ?
\item Mon compromis était-il réaliste et formulé poliment (\eng{should, I suggest that...}) ?
\end{itemize}

Ces compétences — écouter les deux parties, proposer un compromis, formuler un accord — te seront utiles bien au-delà de la classe d'anglais : elles sont au cœur de la vie en communauté.

\newpage

\coursec{Méthodes et exercices résolus}

\courssub{Méthodes et exercices résolus}

\begin{methode}[Identifier et employer le vocabulaire du thème]
Pour bien employer le vocabulaire du règlement des conflits communautaires, suis ces étapes :
\begin{enumerate}
\item \textbf{Classe le mot dans sa catégorie} : s'agit-il d'un type de conflit (\eng{a land dispute}, \eng{a family quarrel}), d'un acteur (\eng{the chief}, \eng{the elders}, \eng{the mediator}), d'un lieu (\eng{the village council}, \eng{the court}) ou d'une étape du règlement (\eng{to negotiate}, \eng{to reconcile}, \eng{to settle}) ?
\item \textbf{Apprends le mot avec sa collocation} : on dit \eng{to settle a dispute} (régler un différend), \eng{to reach an agreement} (parvenir à un accord), \eng{to file a complaint} (déposer une plainte), jamais le mot isolé.
\item \textbf{Vérifie la nature grammaticale} : \eng{a dispute} (nom) / \eng{to dispute} (verbe) ; \eng{a settlement} (nom) / \eng{to settle} (verbe).
\item \textbf{Choisis le registre adapté} : \eng{to sort out a problem} est familier ; \eng{to resolve a conflict} est soutenu et convient mieux à une rédaction d'examen.
\item \textbf{Teste le mot dans une phrase complète} avant de l'utiliser : \eng{The elders settled the land dispute peacefully.} « Les anciens ont réglé le conflit foncier pacifiquement. »
\end{enumerate}
\end{methode}

\begin{exoresolu}[Compléter un texte à trous]
Énoncé — Complete the following passage with the correct words from the box: \emph{dispute, elders, mediator, agreement, complaint, peacefully}.

\eng{Last month, a serious land (1) \ldots\ broke out between two families in our village. The (2) \ldots\ of the community met both parties at the chief's palace. A neutral (3) \ldots\ listened to each side carefully. After long discussions, the two families finally reached an (4) \ldots\ and shook hands. The conflict was settled (5) \ldots\ , without going to court.}
\tcblower
Solution détaillée :
\begin{itemize}
\item (1) \textbf{dispute} : la collocation attendue est \eng{a land dispute} « un conflit foncier » ; le verbe \eng{broke out} « éclater » confirme qu'il s'agit d'un conflit.
\item (2) \textbf{elders} : \eng{the elders of the community} « les anciens de la communauté » ; ce sont eux qui traditionnellement réunissent les parties au palais du chef.
\item (3) \textbf{mediator} : l'indice est l'adjectif \eng{neutral} ; le médiateur est la personne neutre qui écoute les deux camps (\eng{each side}).
\item (4) \textbf{agreement} : collocation figée \eng{to reach an agreement} « parvenir à un accord » ; attention à l'article \eng{an} devant voyelle.
\item (5) \textbf{peacefully} : il faut un adverbe pour modifier le verbe \eng{was settled} ; \eng{peacefully} « pacifiquement » s'oppose à \eng{going to court}.
\end{itemize}
Texte complété : \eng{Last month, a serious land \textbf{dispute} broke out between two families in our village. The \textbf{elders} of the community met both parties at the chief's palace. A neutral \textbf{mediator} listened to each side carefully. After long discussions, the two families finally reached an \textbf{agreement} and shook hands. The conflict was settled \textbf{peacefully}, without going to court.}
Conclusion : toujours s'appuyer sur les indices grammaticaux (article, adverbe, collocation) présents dans la phrase.
\end{exoresolu}

\begin{methode}[Réemployer le lexique à l'oral et à l'écrit]
Pour réutiliser ce vocabulaire dans tes propres productions :
\begin{enumerate}
\item \textbf{Prépare une « banque de phrases »} par fonction : exposer le conflit (\eng{The dispute concerns\ldots}), proposer une solution (\eng{I suggest that both parties should\ldots}), conclure (\eng{Finally, they reconciled and\ldots}).
\item \textbf{Utilise les connecteurs logiques} : \eng{first}, \eng{then}, \eng{after that}, \eng{finally} pour raconter les étapes du règlement ; \eng{however}, \eng{therefore} pour argumenter.
\item \textbf{À l'oral} (jeu de rôle, débat) : emploie des formules de médiation : \eng{Let's listen to both sides.} « Écoutons les deux parties. » ; \eng{Could you explain your point of view?} « Pourriez-vous expliquer votre point de vue ? »
\item \textbf{À l'écrit} : varie les verbes (\eng{to settle}, \eng{to resolve}, \eng{to mediate}, \eng{to reconcile}) pour éviter les répétitions, signe d'un bon niveau à l'examen.
\item \textbf{Relis-toi} en vérifiant les prépositions : \eng{to listen \textbf{to}}, \eng{to agree \textbf{with} somebody}, \eng{to apologize \textbf{for}}.
\end{enumerate}
\end{methode}

\begin{exoresolu}[Dialogue de médiation à compléter et à jouer]
Énoncé — Complete the dialogue between a mediator and two neighbours, Mr Biya and Mr Etame, who disagree about a fence. Use: \emph{point of view, solution, apologize, both sides, compromise}.

\textbf{Mediator:} Good morning. I am here to listen to (1) \ldots\ . Mr Etame, could you explain your (2) \ldots\ ?

\textbf{Mr Etame:} Yes. Mr Biya built his fence on my land!

\textbf{Mediator:} I see. Mr Biya, do you have anything to say?

\textbf{Mr Biya:} I made a mistake. I (3) \ldots\ for that.

\textbf{Mediator:} Good. Now, we need a fair (4) \ldots\ . I propose a (5) \ldots\ : the fence will be moved one metre back, and Mr Etame will pay half of the cost.

\textbf{Both:} We agree!
\tcblower
Solution détaillée :
\begin{itemize}
\item (1) \textbf{both sides} : \eng{to listen to both sides} « écouter les deux parties » ; le médiateur doit entendre les deux voisins.
\item (2) \textbf{point of view} : \eng{to explain your point of view} « expliquer votre point de vue » ; le possessif \eng{your} impose un nom.
\item (3) \textbf{apologize} : collocation \eng{to apologize for something} « s'excuser de quelque chose » ; l'indice est \eng{I made a mistake} « j'ai fait une erreur ». (Présent simple performatif : \eng{I apologize}).
\item (4) \textbf{solution} : \eng{a fair solution} « une solution équitable » ; l'adjectif \eng{fair} précède le nom recherché.
\item (5) \textbf{compromise} : \eng{to propose a compromise} « proposer un compromis » ; chacun fait un geste (déplacer la clôture / payer la moitié).
\end{itemize}
Conclusion : dans un dialogue de médiation, repère qui parle et quelle fonction il remplit (écouter, s'excuser, proposer) pour choisir le mot juste.
\end{exoresolu}

\begin{methode}[Éviter les faux amis et les calques du français]
\begin{enumerate}
\item \textbf{Repère les faux amis du thème} : \eng{a dispute} = un différend (pas « une dispute » au sens de querelle d'enfants, qui se dit \eng{a quarrel}) ; \eng{to demand} = exiger (et non « demander », qui se dit \eng{to ask}) ; \eng{a process} = une procédure (et non « un procès », qui se dit \eng{a trial}).
\item \textbf{Méfie-toi des calques de structure} : « écouter quelqu'un » = \eng{to listen \textbf{to} somebody} (préposition obligatoire) ; « être d'accord avec » = \eng{to agree \textbf{with}} ; « discuter de quelque chose » = \eng{to discuss something} (SANS \eng{about} !).
\item \textbf{Vérifie le verbe dans un dictionnaire ou ta leçon} avant de l'écrire : demande-toi « quelle préposition suit ce verbe en anglais ? ».
\item \textbf{Traduis mentalement ta phrase en français littéral} : si le résultat sonne bizarre, c'est probablement un calque.
\item \textbf{Mémorise des paires contrastées} : \eng{to ask a question} / \eng{to demand justice} ; \eng{a quarrel} / \eng{a dispute}.
\end{enumerate}
\end{methode}

\begin{exoresolu}[Corriger les calques du français]
Énoncé — Each sentence contains a mistake caused by French interference. Correct it.

1. The chief listened the two parties carefully.

2. They discussed about the problem for two hours.

3. The mediator demanded a question to Mr Etame.

4. I am agree with the court's decision.
\tcblower
Solution détaillée :
\begin{enumerate}
\item \eng{The chief listened \textbf{to} the two parties carefully.} — Le verbe \eng{to listen} exige la préposition \eng{to} ; le français « écouter quelqu'un » n'en a pas, d'où le calque.
\item \eng{They discussed the problem for two hours.} — \eng{To discuss} est transitif direct : JAMAIS \eng{about} après. « Discuter de quelque chose » = \eng{to discuss something}.
\item \eng{The mediator \textbf{asked} Mr Etame a question.} — Faux ami : \eng{to demand} signifie « exiger » ; « poser une question » se dit \eng{to ask a question}.
\item \eng{I agree with the court's decision.} — \eng{To agree} est un verbe, pas un adjectif : on ne met pas \eng{am} devant. « Être d'accord » = \eng{to agree} (sans auxiliaire \eng{be}).
\end{enumerate}
Conclusion : trois pièges classiques des francophones — la préposition après \eng{listen}, le \eng{about} interdit après \eng{discuss}, et le verbe \eng{agree} employé sans \eng{be}.
\end{exoresolu}

\begin{methode}[Maîtriser les mots de la même famille, l'orthographe et la prononciation]
\begin{enumerate}
\item \textbf{Construis les familles de mots} autour d'une racine : \eng{to settle} (verbe) → \eng{settlement} (nom) → \eng{settled} (adjectif) ; \eng{to mediate} → \eng{mediation} → \eng{mediator} (la personne).
\item \textbf{Retiens les suffixes de personnes} : \eng{-or} / \eng{-er} désignent celui qui agit : \eng{mediator} (médiateur), \eng{arbitrator} (arbitre), \eng{negotiator} (négociateur).
\item \textbf{Attention à l'orthographe} : \eng{peaceful} (un seul \emph{l} à l'intérieur du mot, mais \eng{peacefully} avec \emph{ll}) ; \eng{argument} (sans \emph{e}, contrairement à \eng{argue}) ; \eng{agreement} (avec \emph{-ment}).
\item \textbf{Note la prononciation en lettres françaises} : \eng{dispute} « dis-pioute », \eng{mediator} « mi-di-é-teur », \eng{reconcile} « ré-kon-saïl », \eng{compromise} « kom-pro-maïz ».
\item \textbf{Vérifie l'accent tonique} : il tombe souvent sur la première syllabe des noms (\eng{DI-spute} comme nom) — cela change le sens et la nature du mot.
\end{enumerate}
\end{methode}

\begin{exoresolu}[Familles de mots : trouver la bonne forme]
Énoncé — Complete each sentence with the correct form of the word in brackets.

1. The \ldots\ of the conflict took several weeks. (settle)

2. A \ldots\ was appointed to help the two villages. (mediate)

3. Both parties signed the peace \ldots\ yesterday. (agree)

4. Their \ldots\ was finally resolved by the village council. (argue)
\tcblower
Solution détaillée :
\begin{enumerate}
\item \textbf{settlement} : après l'article \eng{the}, il faut un nom ; \eng{to settle} → \eng{settlement} « règlement ». Phrase : \eng{The settlement of the conflict took several weeks.} « Le règlement du conflit a pris plusieurs semaines. »
\item \textbf{mediator} : l'article \eng{a} et le sujet de \eng{was appointed} désignent une personne ; suffixe \eng{-or} → \eng{mediator} « médiateur ».
\item \textbf{agreement} : après le nom \eng{peace} utilisé comme classificateur (nom adjectif), il faut un nom ; \eng{to agree} → \eng{agreement} « accord ». \eng{A peace agreement} « un accord de paix ».
\item \textbf{argument} : après le possessif \eng{their}, il faut un nom ; attention à l'orthographe : \eng{argue} perd son \emph{e} final → \eng{argument} « dispute, querelle ».
\end{enumerate}
Conclusion : repère l'indice grammatical (article, possessif, nom classificateur) qui précède le trou pour déterminer la nature du mot attendu, puis applique la règle de formation et d'orthographe.
\end{exoresolu}

\begin{methode}[Rédiger un court texte sur un conflit communautaire]
Pour une composition ou un paragraphe guidé à l'examen :
\begin{enumerate}
\item \textbf{Introduction} : présente le conflit (qui ? où ? quel problème ?) : \eng{Last year, a land dispute opposed two families in my village.}
\item \textbf{Développement} : raconte les étapes du règlement dans l'ordre chronologique avec \eng{first, then, after that, finally} : plainte → réunion des anciens → médiation → accord.
\item \textbf{Conclusion} : donne le résultat et une leçon : \eng{Thanks to the elders' wisdom, peace was restored in the community.}
\item \textbf{Emploie au moins 6 mots du champ lexical} : \eng{dispute, complaint, elders, mediator, negotiate, agreement, reconcile, peacefully}.
\item \textbf{Relis} : accords, temps (prétérit pour le récit), prépositions, orthographe (\eng{peacefully}, \eng{agreement}, \eng{argument}).
\end{enumerate}
\end{methode}

\begin{exoresolu}[Paragraphe guidé]
Énoncé — Write a short paragraph (about 80 words) narrating how a conflict between two neighbours was settled in your community. Use at least five words from the lesson.
\tcblower
Solution détaillée — démarche : on suit le plan en trois parties (conflit → étapes → résultat), au prétérit, avec les connecteurs chronologiques et le lexique de la leçon.

Réponse rédigée :

\eng{Last year, a serious \textbf{dispute} broke out between two neighbours in my village over a piece of land. First, one of them filed a \textbf{complaint} to the chief. Then, the \textbf{elders} invited both parties to the palace and listened to each side carefully. A wise \textbf{mediator} helped them to \textbf{negotiate} for several hours. After that, they finally reached an \textbf{agreement}: the land was divided equally. Finally, the two neighbours \textbf{reconciled} and shook hands. Peace was restored \textbf{peacefully} in our community.}

« L'année dernière, un grave conflit a éclaté entre deux voisins de mon village à propos d'un terrain. D'abord, l'un d'eux a déposé une plainte auprès du chef. Ensuite, les anciens ont invité les deux parties au palais et ont écouté chaque camp attentivement. Un sage médiateur les a aidés à négocier pendant plusieurs heures. Après cela, ils sont finalement parvenus à un accord : le terrain a été partagé en parts égales. Enfin, les deux voisins se sont réconciliés et se sont serré la main. La paix a été rétablie pacifiquement dans notre communauté. »

Conclusion : le paragraphe contient 7 mots du champ lexical, des connecteurs chronologiques, le prétérit et aucune faute de préposition — les critères d'une excellente copie.
\end{exoresolu}

\begin{attention}[Les 5 erreurs qui coûtent des points au Bac]
\begin{enumerate}
\item \textbf{Oublier la préposition après \eng{listen}} : on écrit \eng{The elders listened \textbf{to} both parties.} et jamais \eng{listened both parties}. Calque direct du français « écouter quelqu'un ».
\item \textbf{Mettre \eng{about} après \eng{discuss}} : \eng{to discuss} est transitif direct. Écris \eng{They discussed the problem.} « Ils ont discuté du problème. » — jamais \eng{discussed about}.
\item \textbf{Confondre les faux amis} : \eng{to demand} = exiger (≠ « demander » = \eng{to ask}) ; \eng{a process} = une procédure (≠ « un procès » = \eng{a trial}) ; \eng{a dispute} = un différend sérieux (≠ une querelle d'enfants = \eng{a quarrel}).
\item \textbf{Écrire \eng{I am agree}} : \eng{agree} est un verbe ! On dit \eng{I agree with you.} « Je suis d'accord avec toi. » — sans auxiliaire \eng{be}.
\item \textbf{Fautes d'orthographe dans les familles de mots} : \eng{argument} (sans \emph{e}, bien que le verbe soit \eng{argue}) ; \eng{peacefully} (avec \emph{ll}) ; \eng{agreement} et \eng{settlement} (avec \emph{-ment}). Une faute sur ces mots fréquents du thème est toujours pénalisée.
\end{enumerate}
\end{attention}

\newpage

\coursec{Exercices}

\courssub{Je vérifie mes connaissances}

\begin{exercice}[QCM : la bonne collocation]
Pour chaque phrase, choisis la bonne réponse (A, B ou C).

\begin{enumerate}
\item The chief helped the two families to \ldots{} their dispute peacefully.\\
A. do \hspace{1cm} B. settle \hspace{1cm} C. make
\item The mediator asked both parties to \ldots{} a compromise.\\
A. reach \hspace{1cm} B. arrive \hspace{1cm} C. take
\item The elders \ldots{} a meeting to discuss the land conflict.\\
A. did \hspace{1cm} B. made \hspace{1cm} C. held
\item After the hearing, the traditional council \ldots{} its verdict.\\
A. said \hspace{1cm} B. gave \hspace{1cm} C. told
\item The two neighbours finally decided to \ldots{} the matter to court.\\
A. bring \hspace{1cm} B. carry \hspace{1cm} C. take
\end{enumerate}
\end{exercice}

\begin{textebox}[A traditional mediation in Bamenda]
Last month, a serious quarrel broke out between two farmers in a village near Bamenda. Both men claimed the same piece of farmland. Instead of going to court, they asked the village council of elders to settle the dispute. The elders listened carefully to both sides, visited the land and questioned the neighbours. After two days of discussion, they proposed a compromise: each farmer would keep the part his father had cultivated. The two men shook hands, shared a calabash of palm wine, and promised to live in peace. The whole village celebrated the reconciliation.
\end{textebox}

\begin{exercice}[Vrai ou faux ? Justifie ta réponse]
Lis le texte \emph{A traditional mediation in Bamenda} ci-dessus, puis dis si les affirmations sont vraies (TRUE) ou fausses (FALSE). Justifie chaque réponse par une courte citation du texte.

\begin{enumerate}
\item The two farmers went to court immediately.
\item The elders heard only one side of the story.
\item The settlement was based on what the fathers had cultivated.
\item The reconciliation ended with a celebration.
\end{enumerate}
\end{exercice}

\begin{exercice}[Vocabulaire : trouve le mot]
Complète chaque définition avec un mot du thème commençant par la lettre indiquée.

\begin{enumerate}
\item A person who helps two sides reach an agreement without judging them : m\ldots
\item An official decision given by a judge or a jury : v\ldots
\item An agreement in which each side accepts less than it wanted : c\ldots
\item A serious disagreement between people or groups : d\ldots
\item The restoration of friendly relations after a conflict : r\ldots
\item A person who brings a case against someone in court : p\ldots
\end{enumerate}
\end{exercice}

\begin{exercice}[Word formation : la bonne forme du mot]
Complète chaque phrase avec la forme correcte du mot entre parenthèses.

\begin{enumerate}
\item The two parties finally reached an \ldots{} after long talks. (agree)
\item The \ldots{} of the case will be announced tomorrow morning. (judge)
\item The \ldots{} between the company and the workers lasted three weeks. (negotiate)
\item Their \ldots{} surprised everyone because they had been enemies for years. (reconcile)
\item The chief is a \ldots{} man who always looks for dialogue. (peace)
\item The lawyer gave a clear \ldots{} of the facts to the court. (explain)
\end{enumerate}
\end{exercice}

\courssub{Je m'entraîne}

\begin{exercice}[Traduction : du français vers l'anglais]
Traduis les phrases suivantes en anglais correct, en utilisant le vocabulaire de la leçon.

\begin{enumerate}
\item Le chef du village a réglé le conflit entre les deux familles.
\item Les deux parties ont accepté de négocier au lieu d'aller au tribunal.
\item Le médiateur a proposé un compromis que tout le monde a accepté.
\item Après la réconciliation, les deux voisins ont enterré la hache de guerre.
\item Les anciens ont écouté les plaignants avant de rendre leur verdict.
\end{enumerate}
\end{exercice}

\begin{exercice}[Texte à trous : cloze test]
Complète le texte suivant avec les mots de la banque ci-dessous. Chaque mot ne peut être utilisé qu'une seule fois.

\begin{center}
\textbf{dispute — mediator — compromise — verdict — reconciliation}
\end{center}

In the village of Mankon, a (1) \ldots{} broke out between two quarters over the use of a new water well. The quarter heads could not agree, so the fon appointed a (2) \ldots{} to help them talk. After several meetings, both sides accepted a (3) \ldots{} : each quarter would use the well on alternate days. The fon then announced his final (4) \ldots{}, which everyone respected. A ceremony of (5) \ldots{} was organised, and the two quarters shared a meal together.
\end{exercice}

\begin{exercice}[Appariement : idiomes et équivalents français]
Relie chaque expression idiomatique anglaise (1--6) à son équivalent français (a--f).

\begin{center}
\begin{tabularx}{\linewidth}{|X|X|}
\hline
\rowcolor{popblueL} \textbf{English idiom} & \textbf{Équivalent français} \\
\hline
1. to bury the hatchet & a. jeter de l'huile sur le feu \\
\hline
2. to add fuel to the fire & b. enterrer la hache de guerre \\
\hline
3. to meet someone halfway & c. laver son linge sale en public \\
\hline
4. to wash one's dirty linen in public & d. couper la poire en deux \\
\hline
5. to pour oil on troubled waters & e. apaiser les tensions \\
\hline
6. to clear the air & f. dissiper le malentendu \\
\hline
\end{tabularx}
\end{center}
\end{exercice}

\begin{exercice}[Attention aux faux amis !]
Les phrases suivantes, écrites par des élèves francophones, contiennent chacune un faux ami ou un calque du français. Souligne l'erreur et réécris chaque phrase correctement.

\begin{enumerate}
\item I demand a question to the mediator.
\item The chief assisted the meeting yesterday.
\item The two brothers are in conflict over the succession of their father.
\item The judge condemned him to pay a compensation of fifty thousand francs.
\item We must find a solution to this actual problem.
\item The elders gave him a good advice about the land dispute.
\end{enumerate}
\end{exercice}

\courssub{Je me prépare au Bac}

\begin{textebox}[Settling disputes the traditional way]
In many Cameroonian communities, going to court is seen as a last resort. When a conflict arises — over land, inheritance, marriage or debts — people first turn to the family head or the village elders. This traditional system of dispute settlement has several advantages. It is free, it is fast, and it takes place in the language everyone understands.

The process usually begins with a complaint. The aggrieved person informs the family head, who summons both parties to a hearing. During the hearing, each side presents its case, and witnesses may be called. The elders listen patiently, because their role is not only to judge but also to reconcile. Their goal is to restore harmony in the community, not to punish one side.

Once the facts are clear, the elders propose a settlement. It may involve an apology, the payment of compensation, or the sharing of the disputed property. The losing party is often asked to provide a symbolic gift — a goat, a fowl or some palm wine — which both families then share. Eating together seals the reconciliation and shows that the conflict is truly over.

Today, some people fear that this system is disappearing as more cases go to modern courts. Yet many lawyers admit that traditional mediation remains essential, especially in rural areas where courts are far away and legal fees are high.
\end{textebox}

\begin{exercice}[Compréhension écrite : reading comprehension]
Lis le texte \emph{Settling disputes the traditional way} ci-dessus, puis réponds aux questions en anglais, par des phrases complètes.

\textbf{Questions}
\begin{enumerate}
\item According to the text, why do Cameroonians prefer traditional dispute settlement? Give two reasons. (2 marks)
\item What is the first step of the traditional settlement process? (2 marks)
\item What is the double role of the elders during a hearing? (2 marks)
\item What does the symbolic gift represent? (2 marks)
\item Find in the text a word or phrase that means : (a) ``a formal accusation'' (paragraph 2) ; (b) ``money paid to repair a wrong'' (paragraph 3). (2 marks)
\end{enumerate}
\end{exercice}

\begin{exercice}[Traduction : version et thème]
\textbf{A. Version.} Traduis le passage suivant en français.

\eng{The elders listen patiently, because their role is not only to judge but also to reconcile. Their goal is to restore harmony in the community, not to punish one side.}

\textbf{B. Thème.} Traduis les phrases suivantes en anglais.

\begin{enumerate}
\item Quand un conflit éclate, les villageois s'adressent d'abord au chef de famille.
\item Les anciens ont proposé un compromis que les deux parties ont accepté.
\item Ce système traditionnel est gratuit, rapide et compris par tout le monde.
\end{enumerate}
\end{exercice}

\begin{exercice}[Expression écrite : guided writing]
Your two best friends, Ali and Brenda, have quarrelled over a borrowed bicycle and no longer speak to each other. As their common friend, you decided to help them reconcile.

Write a paragraph of \textbf{80 to 100 words} in which you :
\begin{itemize}
\item explain the cause of the dispute ;
\item describe the steps you took to mediate between them ;
\item say how the reconciliation finally took place.
\end{itemize}

Use at least \textbf{five words or expressions} from the lesson (dispute, mediator, to settle, compromise, to reconcile, to bury the hatchet, etc.) and underline them in your paragraph.
\end{exercice}

\begin{methode}[Réussir ce paragraphe guidé]
\begin{enumerate}
\item \textbf{Planifie} : note trois idées (cause du conflit, étapes de la médiation, réconciliation) avant de rédiger.
\item \textbf{Choisis tes temps} : le récit se fait au prétérit simple (\eng{quarrelled, decided, listened, agreed}) ; le plus-que-parfait situe une action antérieure (\eng{Brenda had lent Ali her bicycle}).
\item \textbf{Place le lexique} : sélectionne d'abord cinq mots de la leçon, puis construis une phrase naturelle autour de chacun — jamais l'inverse.
\item \textbf{Relie les idées} avec des connecteurs : \eng{first, then, after that, finally}.
\item \textbf{Relis-toi} : vérifie les accords sujet--verbe, les articles et le nombre de mots (80 à 100).
\end{enumerate}
\end{methode}

\begin{exemplebox}[Modèle de production (95 mots)]
\eng{Last month, a \underline{dispute} broke out between Ali and Brenda: she had lent him her bicycle, and he returned it damaged. They stopped talking to each other. As their common friend, I acted as a \underline{mediator}. First, I listened to each of them separately; then I brought them together and helped them \underline{settle} the problem. They finally reached a \underline{compromise}: Ali would pay for the repairs, and Brenda would forgive him. They shook hands and \underline{buried the hatchet}. I was proud to have helped them \underline{reconcile}.}
\end{exemplebox}

\newpage

\coursec{Corrigés des exercices}

\begin{corrige}[Exercice 1 — QCM : la bonne collocation]
\begin{enumerate}
\item \textbf{B. settle} — collocation : \eng{to settle a dispute} (régler un différend). \eng{To do a dispute} et \eng{to make a dispute} n'existent pas en anglais.
\item \textbf{A. reach} — collocation : \eng{to reach a compromise} (parvenir à un compromis). \eng{To arrive} exige la préposition \eng{at} : \eng{to arrive at a compromise}.
\item \textbf{C. held} — \eng{to hold a meeting} (tenir une réunion) ; verbe irrégulier : hold -- held -- held.
\item \textbf{B. gave} — \eng{to give a verdict} (rendre un verdict). \eng{To say} et \eng{to tell} ne se construisent pas avec \eng{verdict}.
\item \textbf{C. take} — \eng{to take a matter to court} (porter une affaire devant les tribunaux).
\end{enumerate}
\end{corrige}

\begin{corrige}[Exercice 2 — Vrai ou faux ?]
\begin{enumerate}
\item \textbf{FALSE} — “Instead of going to court, they asked the village council of elders to settle the dispute.” (Ils n'ont pas saisi le tribunal, ils ont fait appel aux anciens.)
\item \textbf{FALSE} — “The elders listened carefully to both sides.” (Les anciens ont écouté les deux parties, pas une seule.)
\item \textbf{TRUE} — “each farmer would keep the part his father had cultivated.” (Le partage repose sur ce que les pères avaient cultivé.)
\item \textbf{TRUE} — “The whole village celebrated the reconciliation.” (Le village entier a fêté la réconciliation.)
\end{enumerate}
\end{corrige}

\begin{corrige}[Exercice 3 — Vocabulaire : trouve le mot]
\begin{enumerate}
\item \textbf{mediator} (un médiateur) — nom de personne formé sur le verbe \eng{to mediate}.
\item \textbf{verdict} (un verdict) — décision officielle d'un juge ou d'un jury.
\item \textbf{compromise} (un compromis) — chaque partie accepte moins que ce qu'elle voulait.
\item \textbf{dispute} (un différend, un litige) — désaccord sérieux.
\item \textbf{reconciliation} (une réconciliation) — retour à des relations amicales.
\item \textbf{plaintiff} (un plaignant) — attention à l'orthographe : \emph{plaintiff}, avec \emph{-ff} final ; à ne pas confondre avec \eng{defendant} (le prévenu, la partie défenderesse).
\end{enumerate}
\end{corrige}

\begin{corrige}[Exercice 4 — Word formation]
\begin{enumerate}
\item \textbf{agreement} — nom formé avec le suffixe \emph{-ment} : to agree $\rightarrow$ an agreement. \eng{The two parties finally reached an agreement after long talks.}
\item \textbf{judgement} — nom : to judge $\rightarrow$ judgement (orthographe britannique ; l'américain \emph{judgment} est aussi admis). \eng{The judgement of the case will be announced tomorrow morning.}
\item \textbf{negotiation} — nom : to negotiate $\rightarrow$ negotiation. \eng{The negotiation between the company and the workers lasted three weeks.}
\item \textbf{reconciliation} — nom : to reconcile $\rightarrow$ reconciliation. \eng{Their reconciliation surprised everyone because they had been enemies for years.}
\item \textbf{peaceful} — adjectif : peace $\rightarrow$ peaceful (pacifique). \eng{The chief is a peaceful man who always looks for dialogue.}
\item \textbf{explanation} — nom : to explain $\rightarrow$ explanation ; attention, le \emph{i} de \emph{explain} disparaît : expla\textbf{n}ation. \eng{The lawyer gave a clear explanation of the facts to the court.}
\end{enumerate}
\end{corrige}

\begin{corrige}[Exercice 5 — Traduction : du français vers l'anglais]
\begin{enumerate}
\item \eng{The village chief settled the conflict between the two families.}
\item \eng{Both parties agreed to negotiate instead of going to court.}
\item \eng{The mediator proposed a compromise which everybody accepted.}
\item \eng{After the reconciliation, the two neighbours buried the hatchet.}
\item \eng{The elders listened to the plaintiffs before giving their verdict.}
\end{enumerate}
Points de vigilance : « régler » un conflit se dit \eng{to settle} (et non \emph{to regulate}) ; \eng{instead of} est suivi d'un nom ou d'un verbe en \emph{-ing} ; \eng{to bury the hatchet} est une expression figée, au singulier ; « rendre un verdict » se dit \eng{to give a verdict} (\eng{to return a verdict} est correct pour un jury).
\end{corrige}

\begin{corrige}[Exercice 6 — Texte à trous : cloze test]
\begin{enumerate}
\item \textbf{dispute} — collocation : \eng{a dispute broke out} (un différend a éclaté).
\item \textbf{mediator} — nom de personne précédé de l'article \eng{a} : le fon nomme un médiateur.
\item \textbf{compromise} — collocation : \eng{to accept a compromise} (accepter un compromis).
\item \textbf{verdict} — le fon annonce sa décision finale : \eng{his final verdict}.
\item \textbf{reconciliation} — \eng{a ceremony of reconciliation} (une cérémonie de réconciliation).
\end{enumerate}
Traduction du texte complété : « Dans le village de Mankon, un différend a éclaté entre deux quartiers à propos de l'usage d'un nouveau puits. Les chefs de quartier n'arrivaient pas à se mettre d'accord, alors le fon a nommé un médiateur pour les aider à dialoguer. Après plusieurs réunions, les deux parties ont accepté un compromis : chaque quartier utiliserait le puits un jour sur deux. Le fon a ensuite annoncé son verdict final, que tout le monde a respecté. Une cérémonie de réconciliation a été organisée, et les deux quartiers ont partagé un repas ensemble. »
\end{corrige}

\begin{corrige}[Exercice 7 — Appariement : idiomes]
1 -- b ; 2 -- a ; 3 -- d ; 4 -- c ; 5 -- e ; 6 -- f.

Exemples d'emploi :
\begin{itemize}
\item \eng{After years of quarrelling, the two families finally buried the hatchet.} « Après des années de querelles, les deux familles ont finalement enterré la hache de guerre. »
\item \eng{His rude reply only added fuel to the fire.} « Sa réponse grossière n'a fait que jeter de l'huile sur le feu. »
\item \eng{The mediator persuaded both sides to meet each other halfway.} « Le médiateur a persuadé les deux parties de couper la poire en deux. »
\end{itemize}
\end{corrige}

\begin{corrige}[Exercice 8 — Attention aux faux amis !]
\begin{enumerate}
\item Erreur : \emph{demand a question}. Correction : \eng{I asked the mediator a question.} — \eng{to demand} signifie « exiger » ; « demander » se dit \eng{to ask}, et « poser une question » \eng{to ask a question}.
\item Erreur : \emph{assisted the meeting}. Correction : \eng{The chief attended the meeting yesterday.} — \eng{to assist} signifie « aider » ; « assister à » se dit \eng{to attend}.
\item Erreur : \emph{succession}. Correction : \eng{The two brothers are in conflict over their father's inheritance.} — \eng{succession} est un terme juridique rare ; pour « héritage », on dit \eng{inheritance}.
\item Erreur : \emph{condemned}. Correction : \eng{The judge ordered him to pay compensation of fifty thousand francs.} — \eng{to condemn} signifie « condamner (à une peine) » ; ici le juge « ordonne » un paiement : \eng{to order someone to pay}. De plus, \eng{compensation} est indénombrable : pas d'article \emph{a}.
\item Erreur : \emph{actual}. Correction : \eng{We must find a solution to this current problem.} — \eng{actual} signifie « réel, véritable » ; « actuel » se dit \eng{current} ou \eng{present}.
\item Erreur : \emph{a good advice}. Correction : \eng{The elders gave him some good advice about the land dispute.} — \eng{advice} est indénombrable : jamais de \emph{a} ni de pluriel \emph{advices}.
\end{enumerate}
\end{corrige}

\begin{corrige}[Exercice 9 — Reading comprehension]
Barème indicatif : 2 points par question (1 point pour l'information exacte, 1 point pour la correction grammaticale de la phrase complète), soit 10 points.

\begin{enumerate}
\item \eng{Cameroonians prefer traditional dispute settlement because it is free and fast, and because it takes place in a language everyone understands.} (Deux raisons suffisent parmi : gratuit, rapide, langue comprise par tous.)
\item \eng{The process begins with a complaint: the aggrieved person informs the family head, who summons both parties to a hearing.}
\item \eng{The elders' role is not only to judge but also to reconcile the two parties.}
\item \eng{The symbolic gift seals the reconciliation: it is shared by both families and shows that the conflict is truly over.}
\item (a) \eng{a complaint} ; (b) \eng{compensation}.
\end{enumerate}
Points de vigilance : répondre par des phrases complètes avec sujet et verbe ; ne pas recopier de longs passages sans les adapter à la question ; respecter les temps (le texte est au présent de vérité générale).
\end{corrige}

\begin{corrige}[Exercice 10 — Traduction : version et thème]
\textbf{A. Version.} « Les anciens écoutent patiemment, car leur rôle n'est pas seulement de juger, mais aussi de réconcilier. Leur objectif est de rétablir l'harmonie dans la communauté, et non de punir l'une des parties. »

Points de vigilance : \eng{not only \ldots{} but also} se rend par « non seulement \ldots{} mais aussi » ; \eng{to restore harmony} = « rétablir l'harmonie » ; \eng{one side} = « l'une des parties ».

\textbf{B. Thème.}
\begin{enumerate}
\item \eng{When a conflict breaks out, the villagers first turn to the family head.}
\item \eng{The elders proposed a compromise which both parties accepted.}
\item \eng{This traditional system is free, fast and understood by everybody.}
\end{enumerate}
Points de vigilance : « éclater » (pour un conflit) se dit \eng{to break out} (broke out, broken out) ; « s'adresser à » se dit \eng{to turn to} ou \eng{to go to}, jamais \emph{to address to} ; « gratuit » se dit \eng{free} (et non \emph{gratuit}) ; « compris par » = participe passé \eng{understood} (understand -- understood -- understood).
\end{corrige}

\begin{corrige}[Exercice 11 — Expression écrite : guided writing]
Barème indicatif (sur 10 points) : respect du nombre de mots et des trois consignes (3 points) ; emploi d'au moins cinq mots ou expressions de la leçon, soulignés (2 points) ; correction grammaticale et temps verbaux (3 points) ; cohérence et connecteurs (2 points).

Proposition de rédaction modèle (95 mots) :

\eng{Last month, a \underline{dispute} broke out between Ali and Brenda: she had lent him her bicycle, and he returned it damaged. They stopped talking to each other. As their common friend, I acted as a \underline{mediator}. First, I listened to each of them separately; then I brought them together and helped them \underline{settle} the problem. They finally reached a \underline{compromise}: Ali would pay for the repairs, and Brenda would forgive him. They shook hands and \underline{buried the hatchet}. I was proud to have helped them \underline{reconcile}.}

Traduction du modèle : « Le mois dernier, un différend a éclaté entre Ali et Brenda : elle lui avait prêté son vélo, et il le lui a rendu abîmé. Ils ont cessé de se parler. En tant qu'ami commun, j'ai joué le rôle de médiateur. D'abord, j'ai écouté chacun d'eux séparément ; ensuite, je les ai réunis et je les ai aidés à régler le problème. Ils sont finalement parvenus à un compromis : Ali paierait les réparations, et Brenda lui pardonnerait. Ils se sont serré la main et ont enterré la hache de guerre. J'étais fier de les avoir aidés à se réconcilier. »

Points de vigilance : le récit se fait au prétérit simple (\eng{broke out, returned, listened, reached}) ; le plus-que-parfait situe l'action antérieure (\eng{she had lent him her bicycle}) ; \eng{to act as a mediator} (et non \emph{to play a mediator}) ; \eng{to stop} + \emph{-ing} : \eng{they stopped talking} ; vérifier que les cinq mots du lexique sont bien soulignés.
\end{corrige}

\newpage

\coursec{Fiche bilan}

\courssub{Fiche Bilan — Lesson 9 : Settling Community Disputes}

\begin{popfigure}
\begin{tikzpicture}[
  mindmap,
  concept color=popblue,
  text=white,
  font=\small\bfseries,
  level 1/.style={level distance=4.5cm, sibling angle=360/7},
  level 2/.style={level distance=3cm, sibling angle=360/4},
  every node/.style={concept, execute at begin node=\hskip0pt},
  root concept/.append style={concept color=popdark, minimum size=2.8cm},
  child concept/.append style={minimum size=2.2cm},
]
\node[root concept, align=center]{Settling\\ Community\\ Disputes}
  child[concept color=poporange] { node[child concept, align=center]{Types of\\ Conflicts}
    child { node[concept, concept color=poporange!70, minimum size=1.6cm] {Land / Boundary} }
    child { node[concept, concept color=poporange!70, minimum size=1.6cm] {Inheritance / Succession} }
    child { node[concept, concept color=poporange!70, minimum size=1.6cm] {Chieftaincy / Leadership} }
    child { node[concept, concept color=poporange!70, minimum size=1.6cm] {Noise / Nuisance} }
  }
  child[concept color=popgreen] { node[child concept] {Actors}
    child { node[concept, concept color=popgreen!70, minimum size=1.6cm] {Chief / Fon / Lamido} }
    child { node[concept, concept color=popgreen!70, minimum size=1.6cm] {Notables / Elders} }
    child { node[concept, concept color=popgreen!70, minimum size=1.6cm] {Mediator / Arbitrator} }
    child { node[concept, concept color=popgreen!70, minimum size=1.6cm] {Quarter Head} }
  }
  child[concept color=poppurple] { node[child concept] {Places}
    child { node[concept, concept color=poppurple!70, minimum size=1.6cm] {Palace / Traditional Court} }
    child { node[concept, concept color=poppurple!70, minimum size=1.6cm] {Quarter Office} }
    child { node[concept, concept color=poppurple!70, minimum size=1.6cm] {Council Hall} }
    child { node[concept, concept color=poppurple!70, minimum size=1.6cm] {Customary Court} }
  }
  child[concept color=poppink] { node[child concept] {Process Verbs}
    child { node[concept, concept color=poppink!70, minimum size=1.6cm] {Mediate / Arbitrate} }
    child { node[concept, concept color=poppink!70, minimum size=1.6cm] {Reconcile / Settle} }
    child { node[concept, concept color=poppink!70, minimum size=1.6cm] {Rule / Adjudicate} }
    child { node[concept, concept color=poppink!70, minimum size=1.6cm] {Testify / Compromise} }
  }
  child[concept color=popgold] { node[child concept] {Collocations}
    child { node[concept, concept color=popgold!70, minimum size=1.6cm] {Reach a verdict} }
    child { node[concept, concept color=popgold!70, minimum size=1.6cm] {Bury the hatchet} }
    child { node[concept, concept color=popgold!70, minimum size=1.6cm] {Out-of-court settlement} }
    child { node[concept, concept color=popgold!70, minimum size=1.6cm] {Binding decision} }
  }
  child[concept color=popblue] { node[child concept] {Word Families}
    child { node[concept, concept color=popblue!70, minimum size=1.6cm] {Mediate $\to$ Mediation} }
    child { node[concept, concept color=popblue!70, minimum size=1.6cm] {Arbitrate $\to$ Arbitration} }
    child { node[concept, concept color=popblue!70, minimum size=1.6cm] {Adjudicate $\to$ Adjudicator} }
    child { node[concept, concept color=popblue!70, minimum size=1.6cm] {Dispute $\to$ Disputant} }
  }
  child[concept color=popmuted] { node[child concept] {False Friends}
    child { node[concept, concept color=popmuted!70, minimum size=1.6cm] {Prejudice $\neq$ Préjugé} }
    child { node[concept, concept color=popmuted!70, minimum size=1.6cm] {Sentence $\neq$ Phrase} }
    child { node[concept, concept color=popmuted!70, minimum size=1.6cm] {Hearing $\neq$ Audition} }
    child { node[concept, concept color=popmuted!70, minimum size=1.6cm] {Injunction $\neq$ Injection} }
  };
\end{tikzpicture}
\legende{Carte mentale récapitulative du lexique « Settling Community Disputes »}
\end{popfigure}

\begin{bilan}

\textbf{1. Lexique clé — Anglais / Français}

\begin{center}
\begin{tabularx}{\linewidth}{|X|X|X|}
\hline
\rowcolor{popblueL}
\textbf{English} & \textbf{Français} & \textbf{Collocation / Contexte} \\
\hline
\cle{Dispute} / \cle{Conflict} & Différend / Conflit & A \eng{boundary dispute}, a \eng{land conflict} \\
\cle{Grudge} / \cle{Grudge-holder} & Rancune / Celui qui garde rancune & To \eng{bear a grudge against sb} \\
\cle{Altercation} & Altercation, dispute violente & A \eng{heated altercation} \\
\cle{Mediation} / \cle{Mediator} & Médiation / Médiateur & \eng{Peer mediation}, a \eng{neutral mediator} \\
\cle{Arbitration} / \cle{Arbitrator} & Arbitrage / Arbitre & \eng{Binding arbitration}, an \eng{arbitrator's award} \\
\cle{Reconciliation} & Réconciliation & A \eng{reconciliation ceremony}, to \eng{seek reconciliation} \\
\cle{Verdict} / \cle{Ruling} / \cle{Award} & Verdict / Jugement / Sentence arbitrale & To \eng{deliver a verdict}, a \eng{final ruling} \\
\cle{Compromise} / \cle{Settlement} & Compromis / Règlement & To \eng{reach a compromise}, an \eng{out-of-court settlement} \\
\cle{Testimony} / \cle{Witness} & Témoignage / Témoin & To \eng{give testimony}, an \eng{eyewitness} \\
\cle{Oath} / \cle{Swear} & Serment / Jurer & To \eng{take an oath}, to \eng{swear on the Bible/Quran} \\
\cle{Customary law} / \cle{Traditional ruler} & Droit coutumier / Chef traditionnel & \eng{Customary court}, the \eng{Fon / Lamido / Chief} \\
\cle{Quarter head} / \cle{Notable} & Chef de quartier / Notable & The \eng{quarter head} convened the meeting \\
\cle{Boundary} / \cle{Beacon} / \cle{Landmark} & Frontière / Borne / Repère & To \eng{plant beacons}, a \eng{landmark tree} \\
\hline
\end{tabularx}
\end{center}

\textbf{2. Règles de grammaire et formation de mots (à connaître par cœur)}

\begin{center}
\begin{tabularx}{\linewidth}{|X|X|X|}
\hline
\rowcolor{popblueL}
\textbf{Règle / Suffixe} & \textbf{Exemples (Verbe $\to$ Nom / Agent)} & \textbf{Remarque} \\
\hline
\cle{-tion / -sion} (Action/Process) & Mediate $\to$ \eng{Mediation} \newline Arbitrate $\to$ \eng{Arbitration} \newline Reconcile $\to$ \eng{Reconciliation} \newline Adjudicate $\to$ \eng{Adjudication} & Prononciation : « \emph{mé-di-é-ïe-chone} », « \emph{ar-bi-tra-ïe-chone} ». Attention à l'orthographe : \textbf{mediation} (1 t), \textbf{arbitration}. \\
\cle{-or / -er} (Agent) & Mediate $\to$ \eng{Mediator} \newline Arbitrate $\to$ \eng{Arbitrator} \newline Adjudicate $\to$ \eng{Adjudicator} \newline Settle $\to$ \eng{Settler} & \textbf{-or} fréquent pour verbes latins en -ate. \textbf{-er} pour verbes germaniques (settle). \\
\cle{-ee} (Patient/Bénéficiaire) & \eng{Appointee}, \eng{Awardee} & Celui qui reçoit l'action. Rare ici mais possible : \eng{the awardee} (bénéficiaire de la sentence). \\
\cle{Verbe + Particule (Phrasal)} & \eng{Settle down} (s'installer / calmer) \newline \eng{Call off} (annuler) \newline \eng{Back down} (céder, renoncer) \newline \eng{Iron out} (aplanir, régler des détails) & \eng{They ironed out their differences.} (Ils ont aplani leurs différends.) \\
\hline
\end{tabularx}
\end{center}

\textbf{3. Expressions idiomatiques et Faux amis — Pièges à éviter}

\begin{center}
\begin{tabularx}{\linewidth}{|X|X|X|}
\hline
\rowcolor{popblueL}
\textbf{Expression / Mot} & \textbf{Sens correct} & \textbf{Piège / Faux ami} \\
\hline
\eng{To bury the hatchet} & Enterrer la hache de guerre (faire la paix) & Ne pas dire « \eng{bury the axe} ». \\
\eng{To meet halfway} & Faire un compromis, se rencontrer à mi-chemin & Ne pas confondre avec « \eng{meet in the middle} » (littéral). \\
\eng{A bone of contention} & Une pomme de discorde & \eng{Contention} = dispute, pas « contention » (tension musculaire). \\
\eng{To smoke the peace pipe} & Fumer le calumet de la paix (image amérindienne utilisée métonymiquement) & Expression figée. \\
\eng{Prejudice} & Préjudice (juridique) / Parti pris (biais) & \textbf{Faux ami} : \eng{Prejudice} $\neq$ « préjugé » seul (souvent \eng{bias}). En droit : \eng{without prejudice} (sans préjudice). \\
\eng{Sentence} & Peine (juridique) / Phrase (grammaire) & \textbf{Polysémie} : \eng{The judge passed a heavy sentence.} (peine). \\
\eng{Hearing} & Audience (procès) / Audition & \textbf{Faux ami} : \eng{Hearing} $\neq$ « ouïe » (\eng{hearing} = sens, mais \eng{hearing} jurid. = audience). \\
\eng{Injunction} & Injonction (ordre de justice) & \textbf{Faux ami} : $\neq$ « injection » (\eng{injection}). \\
\eng{To rule} & Juger, statuer, rendre ordonnance & $\neq$ « régler » (\eng{regulate}) ni « gouverner » seul. \eng{The court ruled in his favor.} \\
\hline
\end{tabularx}
\end{center}

\textbf{4. Méthodes express}

\end{bilan}

\begin{methode}[Rédiger un compte-rendu de médiation (Minutes of Mediation)]
\begin{enumerate}
\item \textbf{Heading:} \eng{Minutes of the Mediation Session held at [Place] on [Date]}.
\item \textbf{Parties present:} \eng{Complainant: [Name] / Respondent: [Name] / Mediator: [Name] / Witnesses: [Names]}.
\item \textbf{Background:} \eng{Brief statement of the dispute (land boundary, noise, etc.)}. Utiliser le \textbf{Past Simple} pour les faits, \textbf{Present Perfect} pour la situation actuelle (\eng{The dispute has arisen...}).
\item \textbf{Proceedings:} \eng{The mediator opened the session... Each party stated their case...}.
\item \textbf{Agreement / Resolution:} \eng{The parties agreed to... / The mediator ruled that...}. Utiliser \eng{Shall} pour les engagements (\eng{Both parties shall respect...}).
\item \textbf{Signatures:} \eng{Signed: \_\_\_\_\_\_\_\_\_\_ (Complainant) \_\_\_\_\_\_\_\_\_\_ (Respondent) \_\_\_\_\_\_\_\_\_\_ (Mediator)}.
\end{enumerate}
\end{methode}

\begin{methode}[Jouer une médiation (Role-play Speaking)]
\begin{enumerate}
\item \textbf{Préparation (2 min) :} Lire la fiche de rôle (Chef, Partie A, Partie B, Médiateur). Noter 3 arguments clés + 1 concession possible.
\item \textbf{Ouverture :} Le médiateur pose le cadre : \eng{“We are here to find a peaceful settlement. Let’s listen to each other without interrupting.”}
\item \textbf{Échange :} Utiliser \eng{I feel... when... because...} (Communication Non Violente). Le médiateur reformule : \eng{“So you are saying that...”}.
\item \textbf{Négociation :} Proposer : \eng{“What if we...?” / “Would you accept...?” / “Let’s meet halfway.”}
\item \textbf{Clôture :} Lecture de l'accord oral ou écrit. \eng{“Do we have a deal?”} Serrement de main / \eng{“Let’s bury the hatchet.”}
\end{enumerate}
\end{methode}



\begin{aretenir}[Checklist avant le Bac]
$\square$ Je sais définir \cle{dispute}, \cle{mediation}, \cle{arbitration}, \cle{reconciliation} et donner un exemple de collocation pour chacun.\\
$\square$ Je distingue \cle{mediator} (facilite) d'\cle{arbitrator} (décide) et d'\cle{adjudicator} (juge formel).\\
$\square$ Je forme correctement les noms en \textbf{-tion} (\eng{mediation}, \eng{arbitration}, \eng{reconciliation}) et les agents en \textbf{-or} (\eng{mediator}, \cle{arbitrator}).\\
$\square$ J'utilise les phrasal verbs clés : \eng{iron out differences}, \eng{back down}, \eng{call off a strike/protest}.\\
$\square$ J'évite les faux amis : \eng{prejudice} (préjudice/biais), \eng{sentence} (peine), \eng{hearing} (audience), \eng{injunction} (injonction).\\
$\square$ Je connais 3 expressions idiomatiques : \eng{bury the hatchet}, \eng{meet halfway}, \eng{a bone of contention}.\\
$\square$ Je sais identifier les acteurs locaux : \eng{Quarter Head}, \eng{Notables}, \eng{Traditional Ruler (Fon/Lamido/Chief)}, \eng{Customary Court}.\\
$\square$ Je sais rédiger les grandes lignes de \eng{Minutes of Mediation} (en-tête, parties, faits, accord, signatures).\\
$\square$ Je prononce correctement : \eng{mediation} [mi-di-éi-chn], \eng{arbitration} [ar-bi-trei-chn], \eng{reconciliation} [re-kon-si-li-éi-chn].\\
$\square$ Je peux jouer un dialogue de médiation en utilisant \eng{I feel...}, \eng{Would you accept...?}, \eng{Let's compromise.}\\
$\square$ Je sais traduire : « Règlement à l'amiable » $\to$ \eng{Out-of-court settlement / Amicable settlement}.\\
$\square$ Je sais traduire : « Porter plainte » $\to$ \eng{Lodge a complaint / File a suit} (pas \eng{make a plainte}).
\end{aretenir}

\findecours

\end{document}
